% body.tex -- shared content for both build targets.
% The wrapper sets \newif\ifanon and \anontrue / \anonfalse before % body.tex -- shared content for both build targets.
% The wrapper sets \newif\ifanon and \anontrue / \anonfalse before % body.tex -- shared content for both build targets.
% The wrapper sets \newif\ifanon and \anontrue / \anonfalse before % body.tex -- shared content for both build targets.
% The wrapper sets \newif\ifanon and \anontrue / \anonfalse before \input{body}.

% ---- identity macros (switch on \ifanon) ----
\ifanon
  \newcommand{\artref}{the artifact}
  \newcommand{\artobj}{the artifact}
\else
  \newcommand{\artref}{logos v0.1.0}
  \newcommand{\artobj}{the logos artifact (v0.1.0)}
\fi

% ---- PDF metadata: named for arXiv, stripped for the anonymous TACL PDF ----
\ifanon
  \hypersetup{pdfauthor={},pdftitle={},pdfsubject={},pdfkeywords={},pdfcreator={},pdfproducer={}}
\else
  \hypersetup{pdfauthor={Kyriakos Papadopoulos},
    pdftitle={Falsification-First Decipherment: A Decontaminated Inference
Framework for Testing Undeciphered-Script Claims, with Linear A as a
Worked Null},
    pdfcreator={},pdfproducer={}}
\fi

\title{Falsification-First Decipherment: A Decontaminated Inference
Framework for Testing Undeciphered-Script Claims, with Linear A as a
Worked Null}

\ifanon
  \author{}
\else
  \author{
    Kyriakos Papadopoulos \\
    Independent Researcher \\
    Netherlands \\
    ORCID 0009-0003-0995-2518
  }
\fi

\date{}

\begin{document}
\maketitle
\ifanon\else
\begin{center}\small\emph{Preprint --- under review; not peer
reviewed.}\end{center}
\fi

\begin{abstract}
Decipherment research has a false-positive problem and no shared
instrument: a small corpus, an unknown language, and many candidate
readings let an analyst ``translate'' almost anything. The flagship
survey lists fifteen challenges of script and data; none addresses the
orthogonal inference-control layer --- overfitting, multiple testing,
contamination --- a candidate sixteenth challenge. We build a
falsification-first framework for it: a fabricated-language
negative-control family, multiple-comparison deflation, a literature
index separating published from not-indexed signs, and pre-registration
with a mechanically-graded gate. This is a methods paper; we assert no
decipherment. The machinery calibrates: on a tested null the gate's
false-graduation rate is 0.6\% (3/500); the identical morphology test
finding no power on Linear A recovers Mycenaean (Linear B) inflection
(9/16 affixes above the bigram floor), supporting the interpretation
that Linear A's short forms contribute materially to its lack of power.
Across probes it returns differentiated, calibrated outcomes --- null,
no-power, supported-structural, circular, and contamination-sensitive.
Results are consistent with identifiability, rather than corpus size,
being the more fundamental present constraint. The June-2026 *301
``cracked'' claim is a qualitative illustration of the failure modes the
gate is designed to detect.
\end{abstract}

\section{Introduction}

Decipherment has a false-positive problem, and the field has no agreed
instrument for measuring it. The history of undeciphered-script research
is, in large part, a history of uncontrolled model selection: a tiny
corpus, an unknown language, and a large space of candidate readings let
a determined analyst ``translate'' almost anything, with internal
consistency on the cherry-picked set standing in for confirmation. The
mechanism is not any single bad match but uncorrected multiplicity: with
enough candidate roots and a small enough corpus, a pile of
individually-plausible matches is the \emph{expected} outcome under pure
chance, so each individual comparison looks suggestive while the
conclusion is nonsense. A standard demonstration of spurious
comparison: one can ``prove'' \emph{English is a Semitic language} by
matching a handful of English words to Proto-Semitic roots while quietly
ignoring the far larger number that do not match. The same move, applied
to a real undeciphered script, has produced a long line of decipherment
claims that no one can either confirm or refute.

That the discipline lacks a shared standard here is not our editorializing
but the field's own assessment. The flagship computational-decipherment
survey --- \citeauthor{braovic2024}'s \citeyearpar{braovic2024}
systematic review of Bronze Age Aegean and Cypriot scripts ---
enumerates fifteen ``major challenges'' of script, language, and data
(\S2.1) \citep{braovic2024}. Those fifteen are
primarily properties of the \emph{object and its data} --- the script,
the language, the surviving corpus, and the resource/comparison-data
conditions. What
governs whether a \emph{claimed reading} is believable is a different
kind of thing: a property of the inference methodology and its
epistemics --- overfitting, multiple testing, false-positive control,
and contamination. We therefore frame this axis, once and modestly, as a
candidate \textbf{sixteenth challenge} --- a deliberate rhetorical foil
to the survey's fifteen, not a claim to a new taxonomic slot or to sole
discovery. It is \emph{under-instrumented} rather than unrecognised
(\S2.4): our contribution is to
\emph{consolidate and operationalize} that axis for decipherment, not to
invent it. It is the methodological axis the object-level taxonomy was
never meant to cover, and it is distinct from the survey's own
acknowledged evaluation gap (\S2.1) --- an \emph{evaluation-target}
problem, separable from the false-positive /
multiplicity problem; it is the latter this paper instruments.

\textbf{Contribution.} We build the machinery for that axis: a
falsification-first, decontaminated inference framework for testing
undeciphered-script decipherment claims \citep{braovic2024}. The harness
applies fail-closed validation to every positive and enforces four
concrete defences: a fabricated-language canary that supplies a
\textbf{generator-conditional negative-control distribution} (a
false-positive floor conditional on the fabricated-language generator ---
here a first-order sign-bigram Markov model; a broader control ensemble
spanning unigram, root-template, and real-language generators is
disclosed as future work, not built here); multiple-comparison deflation
for the number of signs, roots, and families a search was free to try; a
literature index that partitions published (\texttt{L\_indexed}) from
literature-unseen (\texttt{L\_not\_indexed}) signs so that reproducing a
known value cannot masquerade as discovery; and pre-registration with a
held-out gate whose verdict is computed mechanically, never by the model
that proposed the reading. The evaluation gap gets a concrete, if
power-limited, answer: the recurring libation \emph{formula} --- roughly
50 inscribed objects spread across peak sanctuaries and caves at more
than five findspots --- supports a natural leave-one-site-out
cross-validation graded from held-out data --- explicitly a
low-statistical-power test, reported as such rather than as decisive
(\S2.1; Appendix~A.3) \citep{braovic2024}.

\textbf{This is a methods paper, and Linear A is its worked null.} We do
not claim a decipherment. The framework is designed to withhold
validation when a proposed reading does not clear pre-registered negative
controls, multiplicity adjustment, contamination checks, and held-out
evaluation; when nothing survives, the calibrated null --- what the
corpus can and cannot support --- is itself the deliverable. The working
conclusion on current data is that the results are \textbf{consistent
with identifiability} being the more fundamental present constraint ---
not a demonstrated corpus-size threshold: there is no known related
language to map onto, the multiple-testing burden of candidate readings
goes uncorrected, and the sign inventory is mixed and partly lossy. We do
not claim additional data would be irrelevant --- a null can arise from
identifiability, low power, or corpus size alike. A lossy, non-injective
sign-to-sound channel need not admit a decipherment at any corpus size,
so we make no claim that ``there is enough text to pin a reading''; the
rigorous demonstration of \emph{that} constraint is the result the field
can use. The June-2026 AI-assisted ``Linear A cracked'' claim --- reading
the *301-anchored libation verb as an extinct Semitic form --- serves
throughout as a qualitative illustration of the failure modes the gate is
designed to detect, analysed from its public, archived presentation as a
claim, not a person, and neither reproducible nor held-out validated.
The artifact is open and archived (\S9.7).

\textbf{Roadmap.} Section~2 situates logos against the Braovi\'{c} et al.
taxonomy and the computational-decipherment lineage, including the
cognate-ID baseline. Section~3 specifies the harness --- canary,
deflation, literature index, and the pre-registration/held-out gate.
Section~4 develops the information-floor diagnostic and states its
identifiability limits. Section~5 works the June-2026 *301 claim as a
case study of the failure modes. Sections~6--8 report the pre-registered
nulls: morphology (\S6); metrology, phonology, and cross-script signals
(\S7); and contamination and sufficiency probes (\S8). Section~9 draws
the discussion, limitations, and conclusion.

\section{Related work}

\subsection{Position relative to the survey taxonomy}

The reference frame for computational work on the Bronze Age Aegean
scripts is \citet[\emph{Computational Linguistics}
50(2):725--779]{braovic2024}, the field's first survey focused
specifically on Aegean and Cypriot scripts; its only comparable
predecessor is \citet{ferrara2022}, and the broader
machine-learning-for-ancient-languages landscape is surveyed by
\citet[\emph{Computational Linguistics} 49(3)]{sommerschield2023}.
Braovi\'{c} et al.'s Section~4 enumerates fifteen ``major challenges'' ---
from unknown writing system, reading direction and language, through
small dataset, allographs and incomplete vocabulary, to unavailable
parallel data and hardware --- \textbf{none} of which concerns
overfitting, multiple testing, false-positive control, contamination, or
distinguishing a \emph{fitted} decipherment from a \emph{real} one. The
same gap runs through Ferrara and Tamburini's
\citeyearpar[\emph{Lingue e Linguaggio} XXI.2:239--259]{ferrara2022}
\S3.5, which reviews the cognate-matching lineage logos builds on ---
\citet{snyder2010} on Ugaritic$\leftrightarrow$Hebrew (the 29/30-sign
benchmark our canary calibrates against), \citet{bergkirkpatrick2011},
and \citet{luo2019, luo2021} --- and frames Linear A, via Gelb and
Whiting's typology \citep{gelb1975}, as a Type III script (unknown script
\emph{and} language) whose ``necessary validation'' is ``altogether
impossible,'' yet nowhere treats how a match is to be \emph{deflated} for
the multiplicity it was selected from. Two surveys, the same unfilled
axis. We do not score this as an oversight: their challenges are
properties of the object and its data, whereas the axis logos foregrounds
is a property of inference \emph{methodology} --- named-but-unfilled
rather than absent. logos supplies the machinery to police it (\S3).

Braovi\'{c} et al. do gesture at a related but separable difficulty:
evaluation. Because computational analyses of an undeciphered script have
``nothing to compare them against,'' they conclude (Sec.~6.3) that
``their evaluation remains an open problem.'' That is a distinct problem
from overfitting and multiplicity control, and logos addresses it
separately. Its held-out design is a concrete response --- the libation
formula recurs across a small set of peak-sanctuary and cave findspots,
permitting a leave-one-site-out cross-validation adjudicated mechanically
by \texttt{verdict.py} --- with the low-power caveat flagged at the
outset (\S1). We locate our components in the survey's own
method taxonomy (Sec.~7.1) and, where prior art exists, defer to it
rather than claim novelty.

\subsection{The current philological synthesis}

\citet{salgarella2025}, \emph{Writing in Bronze Age Crete: ``Minoan''
Linear A} (Cambridge Elements), by the author of the SigLA database, is
the most current authoritative one-author overview. It concludes that
Minoan is a language \textbf{isolate} --- it ``does not belong in any
known language family'' (the families used for comparison being
Indo-European, Semitic and Afro-Asiatic) --- typologically agglutinative
with a provisional Verb--Subject--Object order, and it records that ``no
bilingual text has so far been found'' in a corpus of roughly 2,500
inscriptions. Crucially, Salgarella ranks the etymological/comparative
method below combinatorial-contextual analysis and names
\textbf{digital/statistical analysis as ``the next desideratum in the
field.''} logos positions itself as exactly that desideratum --- but
under discipline, treating the isolate verdict as a central
\emph{identifiability} constraint rather than a target to be argued away.

\subsection{Prior computational attempts and the search trap}

The seminal automatic-decipherment result is \citet{luo2019}, which
recovers a character mapping from a lost script to a \emph{known related}
language via neural minimum-cost flow. It auto-deciphered
Ugaritic$\rightarrow$Hebrew ($+5.5\%$ over prior state of the art) and
translated 67.3\% of Linear B$\rightarrow$Greek cognates, but did not
address Linear A and would fail there by construction: with no known
cognate plaintext, the objective has nothing to optimise toward. This
cognate-language precondition is the empirical form of our
identifiability argument. The cognate-ID baseline logos actually runs is
\emph{not} a reimplementation of Luo's neural flow: it is a published
matcher, \citeauthor{tamburini2025}'s \citeyearpar{tamburini2025}
CSA\_OptMatcher (coupled simulated annealing), which reports the same
\citet{luo2019} accuracy metric. We adopt it as the non-neural baseline
against which any Linear A cognate claim must clear.

Against that disciplined baseline stand the in-domain instances of the
search trap: the dictionary-fishing programmes of \citet{eu2019} and
\citet{loh2020}, the latter explicitly a ``brute force attack'' on
dictionaries. These are the concrete Linear-A cases the multiplicity
correction is built to deflate. On the visual side, cross-script
comparison is \emph{not} our novelty: \citet{daggumati2019, daggumati2023}
already applied CNN+SVM cross-script similarity, and \citet{corazza2022}
(``Sign2Vec'') clustered signs unsupervised; our defensible contribution
is only the controlled image-versus-sequence \emph{asymmetry} under one
harness.

\subsection{The discipline lineage}

Two older strands supply the honesty machinery. \citet{packard1974},
\emph{Minoan Linear A}, constructed nine deliberately fictitious
decipherments as a null control --- and the genuine Ventris-value transfer
beat them only about 2:1 --- anticipating our permutation-null canary. The
statistical lineage is selective inference and multiple-comparison
correction \citep{bailey2014}, the correction for the number of
hypotheses $\times$ signs $\times$ roots $\times$ families tried that
turns a raw match-rate into a defensible one; its transfer from the
finance/backtest setting to script decipherment is specified in
Section~3. logos's empty cell is the intersection none of these occupy: a
falsification-first, decontaminated inference framework that couples a
published cognate-ID matcher \citep[reporting the Luo 2019
metric]{tamburini2025} with preregistration, multiplicity deflation and
mechanical held-out verdicts, released openly (\S9.7), with Linear A
worked as the honest null rather than a claimed crack.

\section{Methods: the discipline harness}

The harness is a comparison-layer pipeline that a reader can reimplement.
A candidate reading enters as a \emph{proposal} and passes through four
ordered stages --- pre-registration, deflation, decontamination, held-out
graduation --- under a cross-cutting persistence invariant: every decision
is computed from stored artifacts, never from a model's narration. Model-derived signals --- LLM theses, learned embeddings,
cognate-matcher assignments --- enter capped at confidence $\le 0.75$,
structurally unable to \emph{decide} a verdict; only
mechanically-verified held-out reads rank higher (\S1--2).

\subsection{Pre-registration}

Before any model is built or run, the predictions are frozen. The
morphology pre-registration fixes, for each hypothesis, the claimed affix,
its source, the falsifiable test, and the acceptance threshold, frozen
in-repository before each run so the model cannot be tuned to them after
the fact (the anchoring mechanics and the per-affix acceptance criteria
are given in Appendix~A.1). Outcomes are enumerated in advance --- CONFIRM, EXTEND,
REFUTE/NULL, or NO POWER --- so that a negative is a publishable result,
not a failed experiment. Any change is a new dated pre-registration,
never an edit.

\subsection{Deflation}

Internal fit on the derivation set is treated as no evidence. Every match
rate is deflated for everything tried: \texttt{n\_hypotheses} $\times$
\texttt{n\_signs} $\times$ \texttt{n\_roots}. The nulls are pipeline-level
--- a frequency-banded permutation null (after \citealp{packard1974}) and
the order-statistic bar of \S3.4 (implementation detail in Appendix~D).
Multiplicity is instrumented, not hand-passed, by a deduplicating
\texttt{SearchLog} (invariant 12; Appendix~D). The strong structural statistic,
\texttt{S\_morph}, earns credit only through a productivity gate --- an
affix must attach to $\ge$2 distinct residual stems, so a single formula
word copied across inscriptions scores zero (calibration in
Appendix~A.2, Table~\ref{tab:calib}).

\subsection{Decontamination}

A large open-weight model may have encountered longstanding published
readings by Gordon and Best; reproduction of indexed historical proposals
is therefore treated as potential shared-source contamination rather than
independent evidence. The June-2026 Di Mino claim is analysed separately
as a contemporary public case and is not assumed to be present in the
model's pretraining data --- it is indexed so that future claims and
models cannot reproduce it as discovery. Three mechanisms address this.

\textbf{The literature index} records published Linear A sign-value and
correspondence claims; any proposal matching the index is tagged
\texttt{literature\_match} and quarantined --- it may be tested but can
never count as discovery. The index partitions the sign inventory into
\texttt{L\_indexed} (referenced by $\ge$1 published proposal) and
\texttt{L\_not\_indexed} (untouched). Discovery claims may rest only on
\texttt{L\_not\_indexed} held-out success, since regurgitation can only
return what is in the literature. (The conservative seed, the adversarial
population, and the census are given in Appendix~C.5,
Table~\ref{tab:canary}.)

\textbf{The \texttt{L\_fake} canary} is a phonotactically-plausible
invented lexicon, calibrated to a real candidate language
(Appendix~C.5), generated after the model's release and never published
by us; exact-string collisions with real lexica are measured
(Appendix~C.5). Running the full pipeline
against it yields a \textbf{generator-conditional} false-positive floor:
any apparent lexical support it produces is spurious with respect to the
target historical-language hypothesis, and memorisation cannot
be invoked because there is nothing to memorise. A real candidate's match
must outperform this \emph{one calibrated class of negative controls} by
the corrected margin --- support means the candidate beat this
generator-conditional control family (here a first-order sign-bigram /
root-template generator), not that it defeated every possible false-language
process.

\textbf{LLM-ablation} compares the pipeline with and without the model's
proposals to isolate contamination (procedure in Appendix~C.1).

\subsection{The held-out graduation gate}

A reading graduates only mechanically, from held-out data, as a
conjunction of pre-registered clauses computed by \texttt{scripts/verdict.py}
--- never by the model that proposed the reading. The operative deflation
clause is an \textbf{order-statistic bar}: the observed held-out statistic
must exceed the \emph{expected maximum} of the multiplicity-null taken over
the number of trials the search was free to run (the E[max] over
\texttt{n\_trials} deflation), a \textbf{normal-order-statistic
approximation} in which \textbf{$n =$ \texttt{n\_trials} is the \emph{count of
candidates the search tried}, not an estimated number of independent
hypotheses}. The explicit formula and its Gaussian/independence assumptions
are given in Appendix~D; the error-control claim is carried not by the formula
but by the direct empirical calibration below. This is a direct
correction for selection across the hypothesis
space, and it fails closed on a degenerate null (Appendix~D). The framework yields
four primary evidential outcomes: \textbf{GRADUATE} (every clause holds and the
search multiplicity was instrumented), \textbf{REJECT} (a clause failed, or
the statistic sits below the corrected bar), \textbf{NULL\_PUBLISHED} (the
statistic is indistinguishable from the \texttt{L\_fake} null --- a
calibrated non-result, publishable as such), and \textbf{NO POWER} ---
the last assigned \emph{upstream}, by the pre-registered module outcome
when the experiment cannot adjudicate (\S3.1), before the gate sequence
runs; the gate itself additionally emits a fail-closed \textbf{INCOMPLETE}
status when search multiplicity is un-instrumented. These outcomes concern \emph{evidential
support, not truth}: because Linear A has no ground truth, REJECT / NULL /
NO POWER mean a reading is unsupported or insufficiently validated on
held-out data, \textbf{never that the underlying linguistic proposition is
false}. The verdict follows a fixed decision sequence
(\texttt{verdict.py}): all pre-registered clauses hold on an above-null
match $\rightarrow$ \textbf{GRADUATE}; else a match clearing the
substantive clauses but on an un-instrumented search $\rightarrow$
\textbf{INCOMPLETE}; else an observed statistic not exceeding
the null mean $\rightarrow$ \textbf{NULL\_PUBLISHED}; else $\rightarrow$
\textbf{REJECT}.
We calibrate the refusal rule directly: under the
tested best-of-100 random-map null --- an analyst trying 100 random
candidate maps and submitting the best --- 3 of 500 replicates graduated,
an observed false-graduation rate of 0.6\% (one-sided exact
Clopper--Pearson 95\% upper bound $\approx 1.54\%$), with 487/500 rejected
outright and 10 published as calibrated nulls (Appendix~D); this is an
empirical
error-control result conditional on that null-generating process, not a
universal guarantee. A second
pre-registered clause, \textbf{\texttt{generalizes\_to\_not\_indexed}},
requires the reading to recover literature-unseen (\texttt{L\_not\_indexed})
signs --- so reproducing already-published values cannot by itself
graduate a claim; the clause fails closed (its pre-committed thresholds
are in Appendix~D). Graduation is defined as validated \emph{novel}
discovery; independent held-out support for an already-published reading
is recorded as a distinct outcome and does not graduate. The remaining
pre-registered clauses (\texttt{L\_fake} separation, corrected margin,
order-statistic bar) are conjoined as specified in Appendix~D. These are
\emph{gate parameters} --- pre-registered
choices for this corpus, not universal constants. Validation is
leave-one-site-out, not ordinary \emph{k}-fold (Appendix~A.1). The gate is
AND-monotone --- clauses can only make it stricter --- so hardening never
lets a previously-failing reading graduate.

\subsection{Grade from persisted artifacts}

The invariant across every gate: verdicts are computed from persisted
rows, never from a model's account of its own outcome. The verdict writer
is grep-clean of any model call and is the sole writer of the verdicts
table.

\begin{table*}[t]
\centering
\begin{tabular}{@{}p{2.6cm}p{3.0cm}p{2.4cm}p{3.0cm}p{3.3cm}@{}}
\toprule
\textbf{Probe} & \textbf{Control / comparator} & \textbf{Real outcome} &
\textbf{Null / control floor} & \textbf{Verdict} \\
\midrule
Graduation gate (\S3.4) & --- & --- & best-of-100 random-map null &
3/500 false graduations (0.6\%; calibrated under the tested null) \\
Morphology (\S6) & Linear B 9/16 & Linear A 4/16 & \texttt{L\_fake}
bigram 6/16 & NO POWER \\
Segmentation (\S6) & --- & micro-F1 0.436 & random 0.389 & supported (gap
CI excludes 0) \\
Metrology (\S7.1) & planted 0.9 & 0.000 & permutation 0.113 & null
(agrees with the field) \\
Phonology (\S7.2) & planted (fires $\ge$ 2) & 0.080 / 0.140 & 0.065 /
0.201 & null (data-limited) \\
Image alignment (\S7.3) & font-swap survives & 0.410 & image chance
0.0135 & circular demonstration (capped) \\
Contamination (\S8.1) & model-free 0/40 & 14/40 (JA-NE) & 0/40 &
exploratory (public-exposure gradient) \\
\bottomrule
\end{tabular}
\caption{One harness, calibrated epistemic outcomes across probes.}
\label{tab:harness}
\end{table*}

The same falsification-first harness returns \emph{different} calibrated
verdicts across probes (Table~\ref{tab:harness}) --- the paper's
contribution in one view. Each number is detailed in the cited section.

\section{The information floor}

Before any decipherment method is graded, the corpus must be sized against
what it can identify. Two quantities matter: how much text exists, and how
many free parameters a reading spends against it. Confusing the first for
the second is the origin of the ``too little text'' folk claim --- and, we
argue, of its equally unearned inverse, the optimism that a big-enough
corpus makes a reading identifiable.

\textbf{The corpus.} The current authoritative synthesis
\citep{salgarella2025} gives an envelope of roughly 2,500 total Linear A
finds, of which about 1,400 are readable inscriptions, overwhelmingly
concentrated in the single LM IB destruction horizon (findplace and
dating detail in Appendix~C.4). That near-single horizon is why a
site-held-out split, not a chronological one, is the only realistic
partition. No bilingual text has ever been found
\citep[\S8]{salgarella2025}.

\textbf{Corpus sizing and the unicity computation (numeric detail in
Appendix~C.4).} The recurring sign-count confusion resolves into several
distinct sign-count denominators and two occurrence channels that must be
kept distinct (Table~\ref{tab:denoms}; Appendix~C.4). On the parsed
syllable-sign budget, Shannon's unicity distance sits everywhere far
below the corpus size (Table~\ref{tab:denoms};
\citealp[\S1--3]{salgarella2025}); we retain
this only as a toy-model \emph{precondition} diagnostic and explicitly
withdraw it as any sufficiency or refutation claim. Qualitatively, the
figure cannot settle identifiability either way (Appendix~C.4): unicity
distance presupposes a known plaintext-language oracle we do not have.
Corpus size, in short, cannot decide the question.

\textbf{What actually binds.} The results are consistent with
\emph{identifiability} being the more fundamental present constraint rather
than corpus size \citep[\S8]{salgarella2025}; we do not establish that more
data would be irrelevant. It has three parts: (i) the \textbf{absence of a
known cognate language} --- the map has a value-space but no anchored
target (Luo 2019's limit); (ii) the \textbf{multiple-testing / search
trap} --- that a unique consistent map may exist in principle neither
guarantees an algorithm finds it nor stops a cherry-picked false positive
across the parameters tried (invariant 8); and (iii) a \textbf{mixed
inventory} of syllabograms, logograms, and numerals that a real reading
must first separate. These, not the count of surviving signs, are what a
decipherment must overcome. Accordingly, when a candidate lexicon grows
without bound on a fixed occurrence budget --- free parameters exceeding
the floor --- the reading is underdetermined, and the framework must say
so.

\section{A qualitative audit of the *301 claim}

The framework of Section~3 is abstract until it meets a concrete claim. In
June 2026 a widely circulated ``Linear A cracked'' announcement
\citep{dimino2026} argued that Linear A records an extinct Semitic
language, reviving Gordon's hypothesis \citep{gordon1957, gordon1966}. Its
load-bearing anchor is a single reading: the recurring peak-sanctuary
libation formula word A-TA-I-*301-WA-JA, in which the Linear-A-only sign
*301 is \emph{itself} assigned a /na/-initial value, so the verb root
reads as West-Semitic N-W-Y ``to dwell'' (\emph{nawaya}) and is thereafter
matched to later Hebrew liturgy. We analyse this \textbf{as a claim, on its
public evidence} --- not the person --- from an immutable archived copy
carried in \artobj. It is a clean, self-contained instance of the failure
modes the gate is built to catch, and the three-part refutation below rests
wherever possible on the claim's \textbf{own published table} and on
primary sources \citep{davis2013, duhoux1992, thomas2020, salgarella2025,
steele2017} --- none advanced in logos's own voice, and the paper
introduces no new sign reading. We flag the meta-point once: this is a
literature-and-primary-source argument, not itself a logos held-out result
--- the graduation gate and \texttt{L\_fake} canary are validated as
instruments (unit tests; the known-answer Ugaritic$\leftrightarrow$Hebrew
surrogate) but have \textbf{not been run on this claim}. Scoring Gordon's
equations and the Di Mino lexicon against the \texttt{L\_fake} floor and
the order-statistic bar is directly runnable with the components of \S3 and
is the priority of the planned revision; that the refutation has not been
independently vetted by Aegean epigraphers is a disclosed limitation
(\S9.6), not a standard the work claims to meet.

\textbf{(a) The segmentation is not novel.} Root i-*301 plus affixation is
the reading of \citet{davis2013}; an independent segmentation by
\citet{thomas2020} reaches the same core and is tested head-to-head against
it (the current synthesis, \citealp{salgarella2025}, does not itself adopt
Thomas). On the Davis and Thomas segmentations, the affix-stacking is
evidenced by multi-attestation forms --- ta-na-i-*301-u-ti-nu (IO Za 6) and
a-ta-i-*301-de-ka (ZA Zb 3) --- that vary the material around a stable
i-*301 core. (These sigla are confirmed against Davis 2013: \texttt{TL Za
1}, \texttt{IO Za 6}, and \texttt{ZA Zb 3}.) The claim re-segments a
structure already established in the literature. Under the decontamination
layer this matters procedurally, not just historically: a correspondence
that reproduces a published segmentation is treated as potential
shared-source contamination,
tagged \texttt{literature\_match} and barred from counting as discovery
(\artref, \S C.1).

\textbf{(b) The gloss is contested and weaker.} The mainstream value of the
verb is ``give/dedicate'' --- Duhoux's \citeyearpar{duhoux1992} gloss,
adopted by \citet{davis2013} as a \emph{structural} reading that assigns
*301 \textbf{no phonetic value at all}, and endorsed by the current
synthesis \citep[\S7.2]{salgarella2025}. ``Dwell'' is a semantically
distant, unmotivated substitution. The point is not that ``give'' is
certainly right, but that a \emph{structural} reading needing no
sound-value for *301 already accounts for the data the phonetic claim is
invoked to explain.

\textbf{(c) The primary descriptions of this form's morphology do not
support the Semitic label.} The claim's own table segments the verb as
prefix-stacking and agglutinative --- [prefix A ``I''] [stem-morpheme TA]
[stem-vowel I] [root *301]. We deliberately do \textbf{not} rest on an
\emph{a-priori} typological contradiction here, and we retract the earlier
draft's claim that a prefix-stacked verb is ``internally incoherent'' with
a Semitic reading: Semitic verbal morphology includes productive prefix
conjugations, so a prefix on the verb is not by itself un-Semitic. The
point is empirical, not typological, and it comes from the primary
literature. \citet[p.~38 n.~13]{davis2013} characterises i-*301's
morphology as ``neither IE nor Afroasiatic'' --- and Semitic \emph{is}
Afroasiatic, so a decade earlier the primary description already found
\textbf{no Afroasiatic morphological signature} in this very form. (We rest
on what the footnote states --- the absence of an Afroasiatic signature ---
not on a stronger ``ruled out'' than a hedged footnote can bear.)
Independently, \citet[\S8]{salgarella2025} concludes Minoan is
agglutinative (after Duhoux 1978) and a language isolate. The claim's own
segmentation is thus consistent with the isolate typology, while the
specific Semitic assignment finds no morphological support in the primary
descriptions of the form.

\textbf{The free-parameter receipt.} The decisive point for identifiability
is that *301 has no cross-script anchor: in the claim's own table its
AB-number column is ``---'' and its Linear B correspondence ``(none)''.
Even the standard practice of reading Linear A homomorphs with their Linear
B values is a backward projection to be handled with care
\citep[already cautioned by Duhoux 1989, apud Davis 2013, p.~38
n.~11]{steele2017} --- and an A-only sign such as *301 has no homomorph to
project from at all, so nothing cross-script constrains its value. The
single sign on which the entire decipherment pivots is thus a \textbf{free
parameter} --- its value was assigned to fit the desired root, not
constrained by any independent evidence. This is precisely the case the
minimum-description-length / identifiability \emph{diagnostic} flags: a
value chosen to make one word read out carries no description-length saving
that survives the multiplicity of alternatives it was selected from (cf.
Section~3; the diagnostic is reported next to every claim, and the
operative graduation clause --- the order-statistic bar over exactly that
multiplicity --- is one a hand-fitted value cannot clear). For this reading
the obstacle is identifiability rather than corpus size --- no known cognate
language to constrain the assignment, uncorrected multiplicity, and a sign
whose value is chosen rather than anchored (a free parameter is unanchored
at any corpus size).

Read against Section~3, the claim instantiates two distinct failure modes
at once. It is \textbf{conclusion-driven model specification} --- the
table's \emph{structure} says agglutinative and prefixing, but the
\emph{label} says Semitic, and the label wins because the desired
conclusion (a Semitic prayer) is a fixed prior rather than an emergent
reading. And it is a \textbf{free-parameter / identifiability} failure ---
an unanchored sign-value tuned to the target, on a corpus with no known
cognate language to constrain it. logos indexes the reading (*301 = /na/)
precisely so a model cannot regurgitate it as discovery, and quarantines
rather than dismisses it. The lesson generalises: a match rate on a
hand-picked word, absent a committed prediction and multiplicity
accounting, is not evidence. The *301 claim is a qualitative illustration
of the failure modes the gate is designed to detect.

\section{Results I: morphology induction}

We test whether Linear A's pre-registered affixes and word divisions
reflect genuine morphology. The apparent affixation does not clear the
operative \texttt{L\_fake} bigram floor, so the verdict is NO POWER and
\texttt{has\_validated\_morphology} = False. The identical test
\textbf{validates on Mycenaean Greek (Linear B)}, whose longer words carry
uncontested inflection: it confirms affixation at rate 0.562 --- above the
same bigram floor ($\sim$0.30) --- with \texttt{has\_morphology\_power} =
True, so the control rules out a universally insensitive detector and
supports the reading that Linear A's shorter forms contribute materially to
its lack of power (Table~\ref{tab:morphrates}; Appendix~A.1). One held-out positive stands and survives
its own uncertainty: a DP-unigram segmenter recovers the scribe's own word
divisions at micro-F1 0.436 versus a random-boundary baseline of 0.389, a
gap whose site-clustered 95\% bootstrap CI [0.021, 0.099] excludes zero
(P(gap $>$ 0) = 0.998; Appendix~A.1). This
separates real recoverable structure from fitted morphology --- the paper's
core discipline.

\section{Results II: metrology, phonology, and the cross-script asymmetry}

Three further pre-registered probes exercise the harness against distinct
field claims. Each returns a calibrated null or a truth-capped positive;
together they map where the Linear A corpus does and does not carry
recoverable signal, and they are consistent with identifiability rather
than effort being the more fundamental present constraint.

\subsection{Metrology: a parse-coverage null that agrees with the field}

Direction D poses the accounting tablets as a constraint-optimisation
over the fraction-sign values (Appendix~B.1; Table~\ref{tab:metro}).
Held-out fraction balance
does not beat the shuffle null
($p = 1.0$), so the verdict is NULL; the only value the method determines,
the fraction sign J = 1/2, is the long-standing field consensus
\citep{bennett1950}, not a logos discovery. This is a null the field itself
reached \citep{corazza2021} --- the harness crediting a surviving value to
its originator rather than re-branding it.

\subsection{Distributional phonology: a data-limited null bounded by a
power curve}

Does a sign's distributional context (which signs sit beside it) carry
information about its phonetic value? Tested on the known AB signs, neither
the consonant class nor the vowel class clears chance (permutation
$p = 0.40$ / $0.86$). The verdict is a data-limited NULL: the sound path is
not recoverable from Linear A alone, and no sound value is imputed for any
sign (Table~\ref{tab:phon}; Appendix~B.2).

\subsection{The cross-script asymmetry: image alignment recovers shared
anchors, but the recovery is circular by construction}

Do image embeddings recover the cross-script A$\leftrightarrow$B
correspondence that sequence and cognate co-occurrence cannot? Classical
shape descriptors reach direct-NN \textbf{0.410} on held-out
A$\leftrightarrow$B anchors while sequence and cognate co-occurrence stay
at chance. But the recovery is circular by construction: the
A$\leftrightarrow$B anchor values were themselves assigned by
twentieth-century epigraphers on sign-shape similarity, so this is a
shape-genealogy engineering demonstration, truth-layer-capped $\le$
\textbf{0.75}, not a positive decipherment claim (\S9.6;
Table~\ref{tab:xscript}).

\section{Results III: contamination probes}

One failure mode is invisible to the false-positive / multiplicity axis of
\S7: a decipherment that \emph{looks} novel but silently reproduces
published readings a language model has memorised. This section reports the
two decontamination probes that target it --- both run as method
demonstrations, returning completed (if deliberately bounded) findings.
Throughout, where a clean number was not obtainable we report the
partial-null and the confound rather than a headline.

\subsection{LLM-ablation: does a large open-weight model regurgitate the
literature?}

In the LLM-ablation, \texttt{qwen2.5:72b} reproduces famous
published readings --- Gordon's JA-NE = ``wine'' in 14/40 seeds --- that
the model-free edit-distance baseline never reaches (0/40), yet never
reproduces obscure vessel equations (0/40; Table~\ref{tab:repro}).
Verdict: a reproduction
gradient consistent with differential public exposure --- though the
experiment does not isolate memorisation from correlated item-level
differences --- the memorisation confound the discipline layer must gate
out.

\subsection{\texttt{L\_not\_indexed}: indexed versus not-indexed
abstention}

Does the model's behaviour on \emph{published} signs generalise to signs
unseen in the indexed literature? \texttt{qwen2.5:72b} valued the
memorisable known signs in $\sim$20/40 seeds but abstained $\sim$97\% of
the time on the literature-unseen *-series signs ($\sim$1/40;
Table~\ref{tab:abstain}) --- an
abstention asymmetry consistent with differential public exposure (the
\S8.1 caveat applies). With no ground truth on the not-indexed signs,
\textbf{consistency is not correctness}: it bounds how the model behaves,
never whether it is right.

\section*{Registered extension: a known-answer sufficiency curve}
\addcontentsline{toc}{section}{Registered extension: a known-answer
sufficiency curve}

Is a Linear-A-scale corpus large enough to recover a decipherment if a
known cognate existed? Because every benchmark is handed a real cognate
signal Linear A lacks, any recovery is an \textbf{optimistic benchmark,
conditional on a known cognate relationship}: at-floor at Linear-A scale
$\Rightarrow$ information-insufficient even given a cognate; above-chance
$\Rightarrow$ identifiability (cognate-absence, uncorrected multiplicity)
rather than corpus size would be implicated; neither branch is a
decipherability claim. The completed 2{,}000-step curve (near-converged:
Luvian$\to$Hittite 0.462 vs.\ published 0.475, Phoenician$\to$Ugaritic
0.760 vs.\ 0.805; every recovery a \emph{lower bound}) splits the
branches at the sweep's Linear-A-scale locator (600--700 distinct
word-forms, syllabic normalization): Linear~B$\to$Greek, the primary
analog, stays at its chance floor ($0.044\pm0.023$, $n{=}276$; floor
$\approx0.004$), while Cypriot$\to$Greek sits marginally above
($0.066$, $n{=}346$). Under-convergence keeps the at-floor branch
non-definitive; full per-size and per-cell artifacts are deposited
(Appendix~C.3).

\section{Discussion, limitations, and conclusion}

\subsection{The discipline is the contribution}

The deliverable of this work is not a reading of Linear A. It is a
\emph{harness} --- a literature index with an indexed versus
literature-unseen partition, a fabricated-language control corpus,
model-ablation contamination checks, multiple-comparison correction, and
pre-registration, all graded mechanically from persisted artifacts ---
together with a body of pre-registered, calibrated nulls, each with its
borderline case named. The stratified morphology re-run is the template: at
the pre-registered $\alpha = 0.01$, no pre-registered affix is cross-stratum
stable above the bigram floor, robust to the genre-bucketing choice
(Tables~\ref{tab:strata}--\ref{tab:sens}; Appendix~A.1). A null that
names its threshold, its sensitivity, and its
near-miss is a stronger referee object than a flat negative.

The axis this harness targets is the \textbf{false-positive / multiplicity
axis} --- the mechanism by which a determined analyst ``translates'' almost
anything on a small corpus --- a property of inference methodology, not of
the object or its data (\S1--2).

\subsection{Two named failure modes for the methods literature}

\textbf{(i) Narrative capture.} Structure says X; the label says the
desired Y; the label wins. In the June-2026 *301 claim the published table
segments an agglutinative, prefix-stacking verb while the assigned family
label is Semitic, and the desired output operates as a fixed prior, so
the conclusion conforms to the prior rather than emerging from
constrained readings (\S5). This
is distinct from statistical overfitting and must be named separately.
Notably, the mismatch is empirical --- between the label and the primary
descriptions of the form (\S5) --- not an a-priori typological
impossibility.

\textbf{(ii) Free parameters exceeding the information floor.} A lexicon
reported as ``408 terms and counting'' on a fixed corpus has more degrees
of freedom than the data can identify; a decipherment whose lexicon grows
weekly is underdetermined by construction. The minimum-description-length /
unicity \emph{diagnostic} (\texttt{k $\le$ U\_floor}) is built to flag
exactly this (it is reported next to every claim rather than enforced as a
graduation clause; \S4): the single sign the whole
reading pivots on is the textbook free-parameter case (\S5), its value
assigned to fit the desired root; the refutation rests on the claim's own
published table, archived in \artobj.

\subsection{The post-hoc evaluation sequence}

These failure modes trace to sequencing. The *301 claim publicly commits to
a post-hoc evaluation sequence: declare ``officially deciphered,'' then vet
later. logos's documented inverse is \textbf{pre-register $\rightarrow$
held-out $\rightarrow$ locked artifact} (\S3.1, \S3.4; the held-out
libation-formula corpus and its low-power caveat are in Appendix~A.3).

\subsection{Identifiability appears the more fundamental present constraint}

We take care not to write a defeatist framing, and equally not to
over-claim from corpus size: the toy-model unicity distance
($U \approx 204$--$415$ signs; Table~\ref{tab:denoms}, \S4) is a
\emph{precondition} diagnostic
only, explicitly withdrawn as any sufficiency or refutation claim ---
``$U \ll N$'' does \textbf{not} license the statement that there is enough
text to pin a decipherment (\S4; Appendix~C.4). The
results are consistent with \emph{identifiability} being the more
fundamental present constraint --- the absence of a known cognate language
to map to, an uncorrected multiple-testing search burden, and a mixed sign
inventory --- rather than a demonstrated corpus-size threshold; a null can
arise from identifiability, low power, or size alike, and we do not
establish that more data would be irrelevant.
\citet[\S8]{salgarella2025}
reaches the field-level version of this by a different route (\S2.2).

\subsection{Study-wide error budget}

Multiplicity discipline operates at two levels, and we are explicit about
both. Within a single probe, matches are deflated for everything
searched (\S3.2--3.4). Across the study, the unit of multiplicity is the
\textbf{per-probe pre-registration}: the \textbf{five executed probes}
(morphology, metrology, phonology, cross-script image alignment, and
contamination) each carry their own pre-committed $\alpha$, and every probe
is reported regardless of outcome. \textbf{Error control is defined within
each pre-registered probe; we claim no global family-wise guarantee across
probes.} Nothing in the battery is read as a decipherment.

\subsection{Limitations}

The nulls are honestly bounded, not decisive. A dependence-adjusted
trial-count correction of the morphology z-test is deferred; the deferral
would only strengthen the null (Appendix~A.1). No held-out
\emph{decipherment} has
been achieved, and the study reports \textbf{no positive decipherment
claim} --- the deliverable is calibrated absence, not a reading. Nothing
here claims a rejected reading is \emph{false} (\S3.4): every REJECT /
NULL in this study bounds \emph{evidential support}
on held-out data, not truth. The one probe that recovers an above-chance
number, cross-script sign-image alignment, is a font-robust but circular
\emph{engineering demonstration} (a font-swap control refutes the
typeface-artifact worry; \S7.3, Appendix~B.3), truth-capped at
$\le 0.75$, never a reading. No
phonetic claim is made. The philological points in \S5 rest on the
published literature cited there rather than on independent
Aegean-epigraphic
review, which we disclose as a limitation; because those points are cited
rather than original, the argument does not depend on any new sign reading
of our own. Chronological cross-validation is largely foreclosed by the
single-horizon corpus (\S4; Appendix~A.1). And the corpus itself
--- small, partly lossy, with no bilingual yet found
\citep[\S8]{salgarella2025} --- sets the ceiling.

\subsection{Conclusion and reproducibility}

The framework withholds validation when a proposed reading does not clear
the pre-registered gates (\S1, \S3); the contribution is the
framework and the calibrated null itself --- the instrument and the honest
limits of what the corpus can support --- not a reading, and the calibrated
null is a safeguard against unsupported claims.
\ifanon
The code, pre-registrations, and all findings are open under the MIT
licence and available as an anonymized artifact (code and pre-registrations
released on acceptance).
\else
The code, pre-registrations, and all findings are open under the MIT
licence and archived on Zenodo (concept DOI 10.5281/zenodo.21087572;
repository \url{https://github.com/papadopouloskyriakos/logos}).
\fi
The licensed epigraphic corpora are not redistributed; each source's
licence and retrieval path is documented (data-provenance; Software \&
data note), so the harness is reproducible against a
locally obtained corpus. This work is independent and not affiliated with,
or endorsed by, any of the cited projects or their authors.

\bibliographystyle{acl_natbib}
\nocite{*}
\bibliography{refs}

\appendix

\section{Morphology (replication detail)}

\paragraph{A.1 Morphology: rates, floors, word-length, stratification,
sensitivity.} The probe is pre-registered, unsupervised morphology and
segmentation induction over the full Linear A corpus
(Table~\ref{tab:denoms}). Target claims are the recurring affixes and
positional grammar of the Davis/Thomas/Duhoux--Val\'{e}rio tradition:
prefixes and suffixes (\texttt{i-}, \texttt{a-}, \texttt{-te/-ti}, and
Davis's \texttt{i-*301} verb stem). Pre-specification anchoring (\S3.1):
pre-specifications were frozen in-repository before each run (commit
history); the external immutable deposit
\ifanon
([artifact DOI withheld for review])
\else
(Zenodo concept DOI 10.5281/zenodo.21087572)
\fi
certifies the current state of the registrations, not temporal priority.
Each pre-registered affix
(drawn from \citet{thomas2020} and Duhoux/Val\'{e}rio) must recur on
$\ge$2 distinct stems across independent inscriptions, above a
within-form permutation null, and survive multiplicity correction at a
target of \emph{p} $< 0.01$ corrected. Every candidate was graded against two
null floors: a within-word sign-order permutation (shuffle) floor, and an
\texttt{L\_fake} floor built from a first-order sign-bigram (Markov) corpus
with zero lexical overlap with the real inscriptions.

\emph{Pooled run.} Of the 16 pre-registered affixes, 4 beat the within-form
permutation null and the deflated multiplicity bar (\texttt{i-*301},
\texttt{a-}, \texttt{i-}, \texttt{-te/-ti}), clearing the shuffle floor
decisively (rates and Wilson intervals in Table~\ref{tab:morphrates}).
The apparent standout \texttt{i-*301} sits far above the permutation null
--- but that figure is computed against the weaker within-form permutation
null only, and it does not survive the operative \texttt{L\_fake} bigram
floor, under which the whole set fails: the fabricated Markov corpus
confirms the same affixes at a \emph{higher} rate than the real corpus
(Table~\ref{tab:morphrates}). The outcome is classified NO POWER, not
evidence against morphology; \texttt{has\_morphology\_power} is False and
the module reports the null (the Linear-B positive control that validates
the detector is reported below).

\emph{Word-length distribution.} Linear A words are short
(Table~\ref{tab:morphrates}), and on such words ``morphology'' and
``bigram order'' are statistically indistinguishable. This reconciles
with Fuls's 3.3-sign figure once the denominator is matched:
\citet[57, \S6.5]{fuls2015} computes the mean over a word \emph{list} of
``554 complete sign sequences, where duplicates are reduced to one''
(GORILA; \citealp{godart1985}) --- distinct word \emph{types}, not
tokens; restricting our corpus to distinct multi-sign types likewise
recovers 3.07. Fuls reads the 3.3-sign length as evidence that Minoan is
\emph{polysynthetic}; our direct test finds the apparent affixation
bigram-reproducible with no power, and we report the null. The direct
positive control --- the identical bigram-floor test on Mycenaean Greek
(Linear B), whose longer words carry uncontested inflectional morphology
--- was RUN on the D\=AMOS corpus (parsed by the same
\texttt{\_damos\_wordforms} extractor as \S7.3) against a pre-registered
16-affix Mycenaean panel \citep{ventris1953}. It fires
(Table~\ref{tab:morphrates}), carried by the securely-grammatical
suffixes -jo, -o-jo (gen sg -oio), -de (allative), -qe (enclitic
``and''), -o-i (dat pl), so \texttt{has\_morphology\_power = True}. The
identical test that returns NO POWER on Linear A's short words recovers
Mycenaean morphology on longer ones (a sensitivity check: an ad-hoc
parser yielding mean 2.19 signs/word did \emph{not} fire --- power tracks
word length exactly as the short-word argument predicts), which rules out
a universally insensitive detector and supports the interpretation that
Linear A's shorter forms contribute materially to its lack of power
(\texttt{scripts/\allowbreak comparison/\allowbreak linb\_morphology\_control.py}).

\emph{Segmentation positive.} Under leave-one-site-out cross-validation (52
sites; not k-fold, given formulaic dependence), a DP-unigram
(Goldwater-style) segmenter recovers the scribe's own word divisions above
a random-boundary baseline; a site-clustered bootstrap (resample the 52
sites with replacement, 4,000 draws) puts the real-minus-random gap above
zero, and the paired gap, not the wide marginal per-segmenter CIs, is the
operative test (Table~\ref{tab:morphrates}). Morfessor over-segments the
short syllabic corpus and is reported as not usable rather than hidden
(Table~\ref{tab:morphrates}).

\emph{Stratified re-run.} The strata reconcile exactly to the
1,341-inscription corpus total (Table~\ref{tab:strata}). The
99-inscription ``libation'' \emph{genre} stratum is not the same object
as the libation-\emph{formula} corpus (\S A.3). Following
\citet{salgarellajudson2024}, chronological cross-validation was explicitly
declined (988/1,341 $\approx$ 74\% single-horizon LM IB), and site --- not
scribal hand, which is not individuated for Linear A
\citep{salgarella2019} --- is the independence unit, gated by presence at
$\ge$2 distinct sites. The result is \texttt{has\_validated\_morphology} =
False (\S9.1). Stratification also showed
that even \texttt{i-*301} is bigram-reproducible in the repetitive libation
register, where the Markov corpus manufactures the apparent verb
morphology; Davis's slot grammar is not refuted, merely not separable from
sign bigrams here. Where formula-bearing sites drive a signal, power is
intrinsically low: a leave-one-site-out CV over so few formula findspots
(an effective $\sim$5-fold site partition) cannot resolve a modest effect.

\emph{Sensitivity, disclosed.} The honest 2$\times$2 over genre bucketing
$\times$ $\alpha$ threshold is Table~\ref{tab:sens}; the doubly-relaxed
corner's \texttt{i-} (locative) and \texttt{-te/-ti} (ablative) are
precisely the Duhoux/Val\'{e}rio H3 nominal affixes that
\citet{salgarella2025} \S8 independently endorses --- the named borderline
candidates, fragile to \emph{both} the $\alpha$ threshold and a defensible
genre choice (the $\alpha = 0.01$ null is itself robust to the bucketing).
The present analytical approximation uses attempted-candidate
count rather than a dependence-adjusted effective trial count --- a
recorded, deferred limitation whose correction would only strengthen the
null.

\paragraph{A.2 \texttt{S\_morph} detector calibration (\S3.2).} After the productivity-gate fix (\S3.2 --- an affix must attach to
$\ge$2 distinct residual stems), \texttt{S\_morph}'s measured false-positive
rate on within-form-shuffled corpora is consistent with the nominal
one-sided $z \ge 2$ rate within sampling error, and its planted-affix
power is a saturated point, not a full sensitivity curve
(Table~\ref{tab:calib}; contrast the phonology power curve, \S7). This is
an \textbf{indicative calibration check, not the operative gate}: 6/200
is too coarse to certify an $\alpha = 0.01$ threshold on its own, and the
gate applies the order-statistic bar over the far larger permutation and
candidate counts the deflation machinery tracks (\S3.4).

\paragraph{A.3 Held-out design: the libation-formula corpus (\S9.3).} A
leave-one-site-out split over so few formula sites (\S1) has
low statistical power --- and site-level CV is in any case necessary, not
sufficient, since within-site scribal-hand correlation must also be controlled
\citep{salgarellajudson2024}. This formula corpus should not be conflated with
the 99 stone-vessel inscriptions of the ``libation'' genre stratum (\S6): they
are different objects.

\section{Metrology, phonology, and cross-script (methods and controls)}

\paragraph{B.1 Metrology.} Direction D parses HT (LM I) documents into
balance units \texttt{[recipient][commodity logogram][integer][fraction
signs]} and solves the fraction-sign values by RANSAC over exact-\texttt{Fraction}
Gaussian elimination (robust to the unbalanced-tablet outliers that defeat
least-squares), requiring line items to balance the preserved KU-RO totals;
validation is grouped k-fold by document (no tablet straddles train and
test) against a permutation null. Results are in Table~\ref{tab:metro}: of
$\sim$7 fraction-bearing constraints only one is satisfied by the solved
system, the sole determinable value is J = 1/2, held-out fraction balance
is not separated from the permutation null, and the synthetic positive
control (planted values, all recovered) separates. The null traces to sparse
automated parse coverage (multi-commodity blocks, illegible number glyphs,
KI-RO deficit sub-lists, editorial restorations); \citet[p.~2]{corazza2021}
abandoned the balance-to-totals method for the same reason. The surviving J
= 1/2 (\S7.1) gains Schrijver's
\citeyearpar{schrijver2014} three further tablets (HT Zd 156, PE 1, HT
123a): at least four corroborating documents.

\paragraph{B.2 Distributional phonology.} Track 2 asks whether a sign's
distributional context --- which signs sit beside it --- carries
information about its phonetic value, measured on the known \texttt{AB}
signs whose Linear B sound values can be borrowed. Features are
PPMI-weighted, L2-normalised left/right co-occurrence counts; labels are
the Linear-B-transferred (C, V) parsed from the GORILA name, so features
never see the labels. The test is leave-one-out 1-NN against a
label-permutation null, \v{S}id\'{a}k-deflated over the two tests; 50 of
$\sim$55 AB signs met the $\ge$3-occurrence threshold.

Per-test results are in Table~\ref{tab:phon}. The honest verdict is
\emph{data-limited},
not \emph{no bridge exists}, because the positive control disciplines the
reading: the positive-control power curve over planted strengths shows
the test firing at strength $\ge$ 2 but not at 1. The test
works but is not infinitely sensitive at n = 50 / 13 classes, so we cannot
distinguish ``no distribution$\rightarrow$phonetics bridge'' from ``too few
signs to detect a weak one'' (\S7.2, \S B.3).

\paragraph{B.3 Cross-script image alignment.} The palaeographic probe
renders 88 Linear-A and 74 Linear-B glyphs and tests held-out recovery of
shared A$\leftrightarrow$B anchors. Classical descriptors (HOG + Hu moments
+ shape-context) sit far above image chance, and a from-scratch I-JEPA is
excluded from the claims as over-parameterized for $\sim$90 glyphs
(results and controls in Table~\ref{tab:xscript}). Two controls decide how
to read the image-vs-sequence asymmetry (\S7.3), and together they
\emph{downgrade} it from a positive to a demonstration.

\emph{First, a font-swap control.} Rendering both scripts in one typeface
(Aegean, George Douros) risks making the alignment a single-designer
house-style artifact rather than a fact about the signs. We re-render Linear
A in Noto Sans Linear A and Linear B in Noto Sans Linear B --- two
independent type designers, each font covering only its own Unicode block
(Noto-A carries zero Linear B glyphs and vice-versa), so no single hand
draws both scripts. The recovery survives and is if anything stronger
(Table~\ref{tab:xscript}; both legs $\gg$ chance): the font-artifact
hypothesis is \emph{refuted}, and ``distances reflect shape, not font'' is
now demonstrated rather than asserted
(\texttt{scripts/\allowbreak palaeo/\allowbreak font\_control.py}).

\emph{Second, and decisively, the circularity the font control cannot
remove.} Linear B was historically adapted from Linear A, and the
A$\leftrightarrow$B transcription \emph{values} that define the anchor set
were assigned by twentieth-century epigraphers largely on the basis of
sign-shape similarity, so the recovery re-derives a shape-defined
correspondence (\S7.3): it confirms a known palaeographic fact (borrowed
signs look alike) and validates a font-robust descriptor pipeline, but it
reads no Minoan. It remains truth-layer-capped in any case --- the
\citet{ferrara2021} ``Archanes Formula'' hard negative (\texttt{CH 095}
resembles \texttt{AB 60}/\emph{ra} but is not its phonetic counterpart)
shows a plausible shape match can be phonetically wrong, so a shape match is
a $\le 0.75$ signal (invariant 5), never a reading. The honest asymmetry is
thus: the image path re-derives a shape-defined mapping (expected, circular,
capped), while the sequence and cognate paths \textbf{remain at chance even
after ingesting the full D\=AMOS Mycenaean corpus} --- 13,562 wordforms,
$\approx$15$\times$ the earlier cognate file, 56 anchors
(Table~\ref{tab:xscript}) --- while the \emph{positive control
recovers}, so the harness works and the null is a real null: the
A$\leftrightarrow$B phonetic correspondence is simply not carried by
cross-script sign co-occurrence, even at full Linear-B scale
(\S B.2). And the A-only signs on which
positive-decipherment claims would pivot have no cross-script
anchor at all (\S5), so images cannot impute them even in principle.

\section{Contamination and sufficiency (detail)}

\paragraph{C.1 LLM-ablation reproduction gradient.} The ablation
(\S3.3) is the remaining gate on the contamination number, matching a
model's cognate glosses against the indexed lexical claims. We ran it
against \texttt{qwen2.5:72b} --- a large open-weight model, a far
stronger memoriser than a local 12B --- on a rented H100, with an
NW-Semitic candidate and a Hebrew skeleton lexicon for the model-free arm
(scale and quantities in Table~\ref{tab:canary}) (\artref). The
first-pass contamination rate is an \textbf{artifact, reported as such}:
the mechanical arm never shares a literature-matching correspondence with
the LLM (their value vocabularies barely overlap), so the rate is forced
to 1.0 by construction and is not reported as a contamination result
(Table~\ref{tab:canary}; \artref).

\emph{Generation settings (replication).} Checkpoint \texttt{qwen2.5:72b}
as served by Ollama (Ollama's default \texttt{q4\_K\_M} quantization;
Ollama version not recorded), on an ephemeral vast.ai H100 SXM 80\,GB
instance (torn down after the run). One \texttt{/api/generate} call per
seed, stream off, timeout 180\,s; temperature 0.8; sampling seed = the
replicate seed; top-p, top-k, maximum output length, and stop conditions
left at server defaults (not recorded). Each
seed's 40 held-out forms are drawn deterministically
(\texttt{numpy.default\_rng(seed)}) and rendered
as canonical sign names in a fixed JSON-request prompt
(\texttt{\_PROMPT\_TEMPLATE} in
\texttt{scripts/\allowbreak comparison/\allowbreak ablation.py}). Parsing:
the JSON \texttt{partial\_map} is extracted; keys not naming a sign shown
to the model are dropped (the proposer cannot invent signs); a dead or
garbled call contributes an empty proposal without crashing the batch
(fail-closed). The gate-2 rerun
prepends the 7 indexed probe forms to every seed. Per-call logs:
\texttt{runtime/ablation/llm\_calls.jsonl}; harness commits
\texttt{0582966} (ablation) and \texttt{9cd6d04} (gate-2).

The gate-2 rerun, reconciled to a common consonant-skeleton vocabulary,
yields the honest deliverable: the per-item reproduction gradient over 40
seeds (Table~\ref{tab:repro}; \artref). An
external red-team correction split \texttt{KU-RO} between the
Semitic-root route (\emph{kull}) and the general-knowledge meaning
``total'' (confirmed by tablet arithmetic), so \texttt{KU-RO} = ``total''
is \textbf{excluded} as a witness and \texttt{KU-RO} = \emph{kull} is
downweighted (Table~\ref{tab:repro}; \artref). The surviving gradient ---
famous readings reproduced, obscure ones not --- is what the ablation
detects (\S8.1). The gate-2 rate itself remains confounded by the baseline's
representational ceiling and is not reported as a contamination measure
(Table~\ref{tab:canary}).

\paragraph{C.2 \texttt{L\_not\_indexed} abstention asymmetry.} The
\texttt{L\_not\_indexed} probe shows a sign in real context (\S8.2); the
not-indexed signs are absent from the indexed
sources, hence not memorisable from them (never ``impossible to have seen''
in any absolute sense) (\artref). Across 40 seeds, \texttt{qwen2.5:72b}
abstains $\sim$97\% of
the time on the untouched signs while confidently valuing the memorisable
ones (Table~\ref{tab:abstain}). The apparent known$-$not-indexed
consistency gap is coverage-confounded and was flagged, not asserted: the
matched-\emph{n} robustness check returned \texttt{no\_power}, so the
rigorous comparison could not run (Table~\ref{tab:abstain}).
The interpretable signal is the abstention asymmetry itself, read under
the standing rule of \S8.2 (\artref).

\paragraph{C.3 Information-sufficiency CSA sweep (design and deposit).}
The sweep downsamples known-answer benchmarks --- Linear B /
Mycenaean Greek as the primary syllabic analog, sourced through D\=AMOS,
and further syllabic corpora --- and measures held-out cognate-ID
recovery against per-size permutation floors, using the published
cognate-ID matcher of \S2.3 \citep{tamburini2025}, validated in the
artifact against its published Table-3 values (\artref). The reading
rules and the reported curve are in the body (Registered extension); the
full per-size report and per-cell artifacts are deposited in the
repository (\texttt{results/csa/}).
The phonological / sound path is no longer data-deferred: even the full
ingested D\=AMOS corpus leaves the cross-script distributional alignment
at chance (\S7.3, \S B.3).

\paragraph{C.4 Information floor: corpus denominators and the unicity
computation (\S4).} The corpus envelope, the four
sign-count denominators, the two occurrence channels, and the measured
unicity range are tabulated in Table~\ref{tab:denoms}
\citep[Schoep 2002;][\S1--3, \S6]{salgarella2025}; the envelope spans ca.
1800--1450 BCE (Middle Minoan II--Late Minoan IB; \S4). Each denominator
must be labelled where it is used (Table~\ref{tab:denoms}). The channel through
which structure is observed is a \emph{sign-occurrence} count, not a
document count or a sign inventory (invariant 7): the broader $\sim$7,400
occurrences count every sign type, while the narrower
$N \approx 5{,}792$ parsed syllable-signs --- after logograms and
numerals are stripped --- is what the unicity computation consumes. Two
further properties tighten the budget: Linear A words are short and
formulaic (the six-position libation-formula template;
\citealp[\S7]{salgarella2025}), and a single language/horizon dominates.

On the logos-normalized corpus, Shannon's unicity-distance criterion
yields $U \approx 204$--$415$ signs across the whole plausible $(V, P)$
parameter space, everywhere far below $N$ (Table~\ref{tab:denoms}) ---
retained only as a toy-model precondition diagnostic (\S4). Unicity
distance presupposes a \emph{known plaintext-language oracle}; because we
estimate the redundancy $D$ from the ciphertext itself, that oracle
collapses into Linear A up to relabelling, so the computation
demonstrates only measurable recurrent structure, not that any reading is
identifiable; and a lossy, non-injective sign$\rightarrow$sound channel
may possess no unicity distance at all (\S9.4).

\paragraph{C.5 Literature index and \texttt{L\_fake} canary (\S3.3).} \emph{Index seed and adversarial population.}
The seed is deliberately conservative --- a sign wrongly omitted is wrongly
granted \texttt{L\_not\_indexed} status, the dangerous direction --- and
every verifiably-published disputed entry is indexed to \emph{catch}
regurgitation, never as an
accepted reading. Population is adversarial: of the seeded West-Semitic
lexical proposals, five are Gordon's equations and a-sa-sa-ra-me = ``oh
Asherah!'' is \citet{best1981}, not Gordon; provenance-unverifiable
candidates --- \texttt{qa-pa} = kappu and a Semitic \texttt{ki-ro}, whose
attribution to Gordon is not supported by the primary account
\citep{rendsburg1996} --- were excluded rather than indexed. Exposure
indexing is independent of evidential validity --- a false public claim
can still contaminate. Generated from the
seed (invariant 12), the current census (Table~\ref{tab:canary}) is a
deliberately conservative, non-exhaustive seed, not yet the full \S C.1
literature census.
\emph{\texttt{L\_fake} calibration.} The canary is calibrated to a real
candidate language's phoneme frequencies, root-template structure, and
lexicon size (\S3.3). Its own divergence from the calibration
targets is reported, never minimised; with a Semitic root-template sampler
and an external Hebrew (ETCBC/bhsa) reject set, the root-template
total-variation distance fell markedly (Table~\ref{tab:canary}), and any
residual real-lexicon collision is measured and published rather than
asserted to be zero.

\section{The order-statistic graduation bar: formula, assumptions,
calibration}

The operative deflation clause (\S3.4) is a normal-order-statistic
\textbf{approximation} to the expected maximum of the multiplicity-null over
the attempted candidates. For a null of mean $\mu_0$ and standard deviation
$\sigma_0$ and $n$ attempted candidates,
\[
  \resizebox{0.95\columnwidth}{!}{$\displaystyle
  E[\max] \approx \mu_0 + \sigma_0\bigl[(1-\gamma)\,\Phi^{-1}(1-\tfrac{1}{n})
  + \gamma\,\Phi^{-1}(1-\tfrac{1}{ne})\bigr]$}
\]
with $\gamma = 0.5772\ldots$ the Euler--Mascheroni constant --- the
blended expected-maximum expression used by \citet{bailey2014}; the
general order-statistics treatment is \citet[\emph{Order Statistics}, 3rd
ed.]{david2003}. The
pipeline exposes the bar as
\texttt{expected\_max\_order\_stat}($\mu_0, \sigma_0, N$); the
deduplicating \texttt{SearchLog}
(\S3.2, invariant 12) reports the exact distinct-candidate count
\texttt{n\_trials} a scanner produces. The bar fails closed on a degenerate null
($\sigma_0 \le 0 \Rightarrow E[\max] = \mu_0 \Rightarrow$ no graduation).
The \texttt{generalizes\_to\_not\_indexed} clause (\S3.4) requires recovery
above a pre-committed threshold and on a minimum number of distinct
not-indexed signs, and fails closed when no threshold or no supported
not-indexed sign is present. The full clause conjunction: the candidate
must separate from the \texttt{L\_fake} negative-control distribution,
its corrected margin must be cleared, and the observed statistic must
beat the order-statistic bar over the Packard/random-lexeme multiplicity
null.

\textbf{Assumptions and the direction of their error.} The approximation is
exact only under (i) an approximately Gaussian null and (ii) independence
among the $n$ attempted candidates. Neither holds exactly. (i) The
permutation / \texttt{L\_fake} nulls are not perfectly Gaussian; the blended
Euler-$\gamma$ form is a first-order correction, not a guarantee. (ii) The
candidates a real search tries are typically \textbf{positively correlated}
(overlapping sign-value assignments). Under a common-variance Gaussian
comparison model, positive dependence lowers the expected maximum relative
to independence (a Slepian-type comparison; \citealp{slepian1962}), so
using the
raw attempted-candidate count $n$ in place of a smaller effective
independent-trial count places the bar \emph{above} the correlated truth;
this does \textbf{not} establish conservatism for arbitrary heterogeneous
or non-Gaussian dependence --- the empirical best-of-100 calibration
therefore carries the operational error-control claim. We use the
raw deduplicated candidate count; a dependence-adjusted effective count is a
recorded, deferred refinement (\S9.6) whose correction would only strengthen
the nulls.

\textbf{Calibration and its interval.} Because the analytic bar rests on
assumptions that do not hold exactly, the error-control claim is carried not
by the formula but by the direct empirical calibration (\S3.4): a
Monte-Carlo of the full gate on a best-of-100 random-map null-generating
process (\texttt{scripts/\allowbreak gate\_null\_calibration.py}). The
interval on that calibration is the \textbf{one-sided exact
Clopper--Pearson} bound (the Beta-quantile inversion of the binomial), which
needs no normal approximation to the calibration count --- for 3 graduations
in 500 replicates, a 95\% upper bound of $\approx 1.54\%$; the verdict
partition of the 500 replicates is in Table~\ref{tab:calib} (no
INCOMPLETE, the search being instrumented).
\citet{bailey2014} supply the selective-inference lineage (\S2.4) and the
blended expression above; the operational error-control claim nonetheless
rests on the empirical calibration, not on the formula.

\textbf{Cognate-ID baseline (\S3.4).} The match-rate the deflation is
applied to is supplied by the published cognate-ID matcher identified in
\S2.3 (reproduced in the artifact), so the gate does not depend on any
single citation.

\begin{center}\rule{0.5\columnwidth}{0.4pt}\end{center}

\noindent\emph{Software \& data.}
\ifanon
The artifact. Code, pre-registrations, and findings are released as an
anonymized artifact on acceptance. Licensed corpora (GORILA-, SigLA-,
lineara.xyz-, D\=AMOS-derived) are not redistributed; a data-provenance note
documents each source's licence and retrieval path.
\else
logos (v0.1.0). Zenodo. Concept DOI 10.5281/zenodo.21087572; repository
\url{https://github.com/papadopouloskyriakos/logos}. Licensed corpora
(GORILA-, SigLA-, lineara.xyz-, D\=AMOS-derived) are not redistributed; see
\texttt{docs/data-provenance.md}.
\fi

\section{Complementary tables}

\begin{table}[h]
\centering\small
\resizebox{\columnwidth}{!}{%
\begin{tabular}{@{}lrrrl@{}}
\toprule
\textbf{Stratum} & \textbf{Inscr.} & \textbf{Words} & \textbf{Sites} &
\textbf{Role} \\
\midrule
admin & 290 & 1,486 & 14 & induction \\
libation (stone vessels) & 99 & 370 & 15 & induction \\
other & 129 & 187 & 40 & induction \\
seal & 740 & --- & --- & excluded \\
emptied by abbrev.-strip & 83 & --- & --- & excluded \\
\midrule
total & 1,341 & & & \\
\bottomrule
\end{tabular}}
\caption{Morphology corpus strata (Appendix~A.1). Seals are structurally
excluded as an abbreviation channel; the strata reconcile exactly to the
1,341-inscription corpus total.}
\label{tab:strata}
\end{table}

\begin{table}[h]
\centering\small
\resizebox{\columnwidth}{!}{%
\begin{tabular}{@{}lll@{}}
\toprule
\textbf{Bucketing} & $\alpha=0.01$ & $\alpha=0.05$ \\
\midrule
strict (stone vessels only) & NULL & NULL \\
loose (stone objects + inked) & NULL & \texttt{i-}, \texttt{-te/-ti}
validate \\
\bottomrule
\end{tabular}}
\caption{Morphology sensitivity $2\times2$ (Appendix~A.1): genre bucketing
$\times$ significance threshold. Only the doubly-relaxed corner validates
the H3 nominal affixes \texttt{i-} (locative) and \texttt{-te/-ti}
(ablative).}
\label{tab:sens}
\end{table}

\begin{table}[h]
\centering\small
\resizebox{\columnwidth}{!}{%
\begin{tabular}{@{}ll@{}}
\toprule
\textbf{Quantity} & \textbf{Value} \\
\midrule
Linear A confirm rate (16 affixes) & 0.250 (4/16; [0.10, 0.50]) \\
\texttt{L\_fake} bigram-floor rate & 0.375 (6/16; [0.18, 0.61]) \\
\texttt{i-*301} vs.\ within-form null & $z \approx 8.0$ (6 environments) \\
\midrule
Linear B (D\=AMOS) confirm rate & 0.562 (9/16) \\
Linear B within-word-shuffle floor & $\sim$0.03 \\
Linear B \texttt{L\_fake} bigram floor & $\sim$0.30 (5 draws) \\
\midrule
word-token mean length & 1.84 signs (median/mode 1) \\
word tokens $\le$2 signs / 1 sign & 76.1\% / 56.5\% \\
multi-sign \emph{type} mean length & 3.07 (Fuls 3.3, 554-type list) \\
Linear B mean length & 3.23 (13,562 wordforms) \\
\midrule
DP-unigram segmenter micro-F1 & 0.436 \\
random-boundary baseline & 0.389 (base rate 0.406) \\
real$-$random gap, 95\% CI & [0.021, 0.099]; P(gap${>}$0) = 0.998 \\
Morfessor & 0.156 (over-segments; not usable) \\
\bottomrule
\end{tabular}}
\caption{Morphology-induction replication quantities (Appendix~A.1):
confirm rates with Wilson 95\% intervals, word-length denominators, and
the segmentation positive with its site-clustered bootstrap gap (4,000
draws over 52 sites).}
\label{tab:morphrates}
\end{table}

\begin{table}[h]
\centering\small
\resizebox{\columnwidth}{!}{%
\begin{tabular}{@{}ll@{}}
\toprule
\textbf{Quantity} & \textbf{Value} \\
\midrule
documents / balance units & 32 / 35 \\
integer-level balance & 7/28 (25\%) \\
determinable fraction value & J = 1/2 (HT104: $2 \cdot J = 1$) \\
held-out fraction balance (real) & 0.0 \\
permutation null & mean 0.113 (max 0.143); $p = 1.0$ \\
planted positive control & 0.9 vs.\ null 0.22; $p = 0.003$ \\
\bottomrule
\end{tabular}}
\caption{Metrology (Direction D) replication quantities (Appendix~B.1).
The real corpus is not separated from the permutation null; the planted
control separates, so the null is not a machinery bug.}
\label{tab:metro}
\end{table}

\begin{table}[h]
\centering\small
\resizebox{\columnwidth}{!}{%
\begin{tabular}{@{}lrrrr@{}}
\toprule
\textbf{Test} & \textbf{Classes} & \textbf{Acc.} & \textbf{Null} & $p$ \\
\midrule
consonant class & 13 & 0.080 & 0.065 & 0.40 \\
vowel class & 5 & 0.140 & 0.201 & 0.86 \\
\bottomrule
\end{tabular}}
\caption{Distributional-phonology per-test results (Appendix~B.2):
leave-one-out 1-NN vs.\ a label-permutation null, \v{S}id\'{a}k-deflated
over the two tests. The positive-control power curve fires ($p<.05$) at
planted strength $\ge 2$ but not at 1.}
\label{tab:phon}
\end{table}

\begin{table}[h]
\centering\small
\resizebox{\columnwidth}{!}{%
\begin{tabular}{@{}lll@{}}
\toprule
\textbf{Leg} & \textbf{Accuracy} & \textbf{Floor} \\
\midrule
classical, direct-NN & 0.410 [0.18, 0.64] & 0.0135 \\
classical, Procrustes & 0.185 [0.00, 0.36] & 0.0135 \\
I-JEPA (3 seeds; excluded) & $0.324 \pm 0.025$ & 0.0135 \\
font-swap, direct-NN & $0.367 \rightarrow 0.416$ & $\gg$ chance \\
font-swap, Procrustes & $0.156 \rightarrow 0.195$ & $\gg$ chance \\
sequence/cognate (D\=AMOS) & 0.025 & 0.011 \\
\quad positive control & 0.947 & --- \\
\bottomrule
\end{tabular}}
\caption{Cross-script sign-image alignment and controls (Appendix~B.3).
Font-swap rows give same-font $\rightarrow$ independent-fonts accuracy on
the control's matched anchor set; the sequence/cognate row is the best of
five alignment methods over 56 anchors
(\texttt{clearly\_above\_chance} = False).}
\label{tab:xscript}
\end{table}

\begin{table}[h]
\centering\small
\resizebox{\columnwidth}{!}{%
\begin{tabular}{@{}llrr@{}}
\toprule
\textbf{Item} & \textbf{Source} & \textbf{LLM} & \textbf{Baseline} \\
\midrule
\texttt{JA-NE} = \emph{yn} ``wine'' & Gordon & 14/40 & 0/40 \\
\texttt{A-SA-SA-RA-ME} & \citet{best1981} & 2/40 & 0/40 \\
\texttt{KA-RO-PA} & Gordon & 0/40 & 0/40 \\
\texttt{SU-PA-RA} & Gordon & 0/40 & 0/40 \\
\texttt{KU-RO} = \emph{kull} (downwt.) & Gordon & 8/40 & --- \\
\texttt{KU-RO} = ``total'' (excl.) & gen.\ knowledge & 7/40 & --- \\
\bottomrule
\end{tabular}}
\caption{LLM-ablation per-item reproduction over 40 seeds,
\texttt{qwen2.5:72b} vs.\ the model-free edit-distance baseline
(Appendix~C.1). The two \texttt{KU-RO} routes overlap on one seed; three
single-word items are among Gordon's five West-Semitic equations, while
\texttt{A-SA-SA-RA-ME} is Best's.}
\label{tab:repro}
\end{table}

\begin{table}[h]
\centering\small
\resizebox{\columnwidth}{!}{%
\begin{tabular}{@{}lrr@{}}
\toprule
\textbf{Sign set} & \textbf{Signs} & \textbf{Valued (of 40)} \\
\midrule
known AB (indexed) & 20 & $\sim$20 \\
not-indexed *-series & 11 & $\sim$1 ($\sim$97\% abstention) \\
\bottomrule
\end{tabular}}
\caption{\texttt{L\_not\_indexed} abstention asymmetry (Appendix~C.2).
The apparent consistency gap 0.601 (permutation $p = 0.0005$) is
coverage-confounded; the matched-$n$ check returned \texttt{no\_power}
(0 of 11 not-indexed signs reached $n \ge 10$).}
\label{tab:abstain}
\end{table}

\begin{table}[h]
\centering\small
\resizebox{\columnwidth}{!}{%
\begin{tabular}{@{}ll@{}}
\toprule
\textbf{Denominator / channel} & \textbf{Value} \\
\midrule
token-frequent syllabograms & 97 \\
full syllabary estimate & $\sim$150 \\
working repertoire ($V$; conflates syll./subscr./logograms) & 259 \\
incl.\ ideograms (GORILA: 390) & $\sim$350 \\
\midrule
sign occurrences, all types & $\sim$7,400 \\
parsed syllable-signs (unicity input) & $N \approx 5{,}792$ \\
\midrule
corpus envelope & $\sim$2,500 finds; $\sim$1,400 readable \\
findplaces (logos-normalized) & 56 ($\rightarrow$ 52 sites) \\
readable inscriptions & 1,341 \\
distinct word-forms & 539 \\
\midrule
unicity distance over $(V,P)$ space & $U \approx 204$--$415$ signs \\
\bottomrule
\end{tabular}}
\caption{Corpus denominators, occurrence channels, and the unicity range
(Appendix~C.4). The $\sim$7,400 all-type occurrence channel never enters
the unicity arithmetic, which consumes only the parsed syllable-signs.}
\label{tab:denoms}
\end{table}

\begin{table}[h]
\centering\small
\resizebox{\columnwidth}{!}{%
\begin{tabular}{@{}ll@{}}
\toprule
\textbf{Calibration quantity} & \textbf{Value} \\
\midrule
\texttt{S\_morph} FPR (within-form shuffle) & 3.0\% (6/200; [1.4, 6.4]\%) \\
nominal one-sided $z \ge 2$ rate & $\approx$2.3\% \\
\texttt{S\_morph} power (planted affix) & 100\% (50/50; saturated) \\
\midrule
gate false-graduation (best-of-100 null) & 0.6\% (3/500; CP $\le$ 1.54\%) \\
gate verdict partition (500 replicates) & 3 GRAD / 10 NULL / 487 REJ \\
\bottomrule
\end{tabular}}
\caption{Detector and gate calibration quantities (Appendices~A.2
and~D): Wilson interval for the \texttt{S\_morph} false-positive rate;
one-sided exact Clopper--Pearson 95\% upper bound for the gate's
false-graduation rate under the tested best-of-100 random-map null.}
\label{tab:calib}
\end{table}

\begin{table}[h]
\centering\small
\resizebox{\columnwidth}{!}{%
\begin{tabular}{@{}ll@{}}
\toprule
\textbf{Quantity} & \textbf{Value} \\
\midrule
literature-index census & 63 claims / 62 entries (55 single) \\
root-template TV distance & $\sim$0.84 $\rightarrow$ $\sim$0.33
  ($\sim$0.23 standalone) \\
\midrule
ablation scale & 40 seeds $\times$ 40 held-out forms \\
model-free lexicon & 7,682 Hebrew skeletons \\
LLM errors & 0 (40/40 seeds) \\
value-vocabulary overlap & 10 of 134 \\
first-pass contamination rate & 1.000 (33/33 defined; artifact) \\
gate-2 rate & mean 0.646 (confounded; not reported) \\
\bottomrule
\end{tabular}}
\caption{Literature-index, \texttt{L\_fake}-canary, and LLM-ablation
quantities (Appendices~C.1 and~C.5). The TV distance is measured
end-to-end on the Ugaritic$\leftrightarrow$Hebrew gold; the first-pass
contamination rate is forced to 1.0 by construction and reported as an
artifact.}
\label{tab:canary}
\end{table}

\end{document}
.

% ---- identity macros (switch on \ifanon) ----
\ifanon
  \newcommand{\artref}{the artifact}
  \newcommand{\artobj}{the artifact}
\else
  \newcommand{\artref}{logos v0.1.0}
  \newcommand{\artobj}{the logos artifact (v0.1.0)}
\fi

% ---- PDF metadata: named for arXiv, stripped for the anonymous TACL PDF ----
\ifanon
  \hypersetup{pdfauthor={},pdftitle={},pdfsubject={},pdfkeywords={},pdfcreator={},pdfproducer={}}
\else
  \hypersetup{pdfauthor={Kyriakos Papadopoulos},
    pdftitle={Falsification-First Decipherment: A Decontaminated Inference
Framework for Testing Undeciphered-Script Claims, with Linear A as a
Worked Null},
    pdfcreator={},pdfproducer={}}
\fi

\title{Falsification-First Decipherment: A Decontaminated Inference
Framework for Testing Undeciphered-Script Claims, with Linear A as a
Worked Null}

\ifanon
  \author{}
\else
  \author{
    Kyriakos Papadopoulos \\
    Independent Researcher \\
    Netherlands \\
    ORCID 0009-0003-0995-2518
  }
\fi

\date{}

\begin{document}
\maketitle
\ifanon\else
\begin{center}\small\emph{Preprint --- under review; not peer
reviewed.}\end{center}
\fi

\begin{abstract}
Decipherment research has a false-positive problem and no shared
instrument: a small corpus, an unknown language, and many candidate
readings let an analyst ``translate'' almost anything. The flagship
survey lists fifteen challenges of script and data; none addresses the
orthogonal inference-control layer --- overfitting, multiple testing,
contamination --- a candidate sixteenth challenge. We build a
falsification-first framework for it: a fabricated-language
negative-control family, multiple-comparison deflation, a literature
index separating published from not-indexed signs, and pre-registration
with a mechanically-graded gate. This is a methods paper; we assert no
decipherment. The machinery calibrates: on a tested null the gate's
false-graduation rate is 0.6\% (3/500); the identical morphology test
finding no power on Linear A recovers Mycenaean (Linear B) inflection
(9/16 affixes above the bigram floor), supporting the interpretation
that Linear A's short forms contribute materially to its lack of power.
Across probes it returns differentiated, calibrated outcomes --- null,
no-power, supported-structural, circular, and contamination-sensitive.
Results are consistent with identifiability, rather than corpus size,
being the more fundamental present constraint. The June-2026 *301
``cracked'' claim is a qualitative illustration of the failure modes the
gate is designed to detect.
\end{abstract}

\section{Introduction}

Decipherment has a false-positive problem, and the field has no agreed
instrument for measuring it. The history of undeciphered-script research
is, in large part, a history of uncontrolled model selection: a tiny
corpus, an unknown language, and a large space of candidate readings let
a determined analyst ``translate'' almost anything, with internal
consistency on the cherry-picked set standing in for confirmation. The
mechanism is not any single bad match but uncorrected multiplicity: with
enough candidate roots and a small enough corpus, a pile of
individually-plausible matches is the \emph{expected} outcome under pure
chance, so each individual comparison looks suggestive while the
conclusion is nonsense. A standard demonstration of spurious
comparison: one can ``prove'' \emph{English is a Semitic language} by
matching a handful of English words to Proto-Semitic roots while quietly
ignoring the far larger number that do not match. The same move, applied
to a real undeciphered script, has produced a long line of decipherment
claims that no one can either confirm or refute.

That the discipline lacks a shared standard here is not our editorializing
but the field's own assessment. The flagship computational-decipherment
survey --- \citeauthor{braovic2024}'s \citeyearpar{braovic2024}
systematic review of Bronze Age Aegean and Cypriot scripts ---
enumerates fifteen ``major challenges'' of script, language, and data
(\S2.1) \citep{braovic2024}. Those fifteen are
primarily properties of the \emph{object and its data} --- the script,
the language, the surviving corpus, and the resource/comparison-data
conditions. What
governs whether a \emph{claimed reading} is believable is a different
kind of thing: a property of the inference methodology and its
epistemics --- overfitting, multiple testing, false-positive control,
and contamination. We therefore frame this axis, once and modestly, as a
candidate \textbf{sixteenth challenge} --- a deliberate rhetorical foil
to the survey's fifteen, not a claim to a new taxonomic slot or to sole
discovery. It is \emph{under-instrumented} rather than unrecognised
(\S2.4): our contribution is to
\emph{consolidate and operationalize} that axis for decipherment, not to
invent it. It is the methodological axis the object-level taxonomy was
never meant to cover, and it is distinct from the survey's own
acknowledged evaluation gap (\S2.1) --- an \emph{evaluation-target}
problem, separable from the false-positive /
multiplicity problem; it is the latter this paper instruments.

\textbf{Contribution.} We build the machinery for that axis: a
falsification-first, decontaminated inference framework for testing
undeciphered-script decipherment claims \citep{braovic2024}. The harness
applies fail-closed validation to every positive and enforces four
concrete defences: a fabricated-language canary that supplies a
\textbf{generator-conditional negative-control distribution} (a
false-positive floor conditional on the fabricated-language generator ---
here a first-order sign-bigram Markov model; a broader control ensemble
spanning unigram, root-template, and real-language generators is
disclosed as future work, not built here); multiple-comparison deflation
for the number of signs, roots, and families a search was free to try; a
literature index that partitions published (\texttt{L\_indexed}) from
literature-unseen (\texttt{L\_not\_indexed}) signs so that reproducing a
known value cannot masquerade as discovery; and pre-registration with a
held-out gate whose verdict is computed mechanically, never by the model
that proposed the reading. The evaluation gap gets a concrete, if
power-limited, answer: the recurring libation \emph{formula} --- roughly
50 inscribed objects spread across peak sanctuaries and caves at more
than five findspots --- supports a natural leave-one-site-out
cross-validation graded from held-out data --- explicitly a
low-statistical-power test, reported as such rather than as decisive
(\S2.1; Appendix~A.3) \citep{braovic2024}.

\textbf{This is a methods paper, and Linear A is its worked null.} We do
not claim a decipherment. The framework is designed to withhold
validation when a proposed reading does not clear pre-registered negative
controls, multiplicity adjustment, contamination checks, and held-out
evaluation; when nothing survives, the calibrated null --- what the
corpus can and cannot support --- is itself the deliverable. The working
conclusion on current data is that the results are \textbf{consistent
with identifiability} being the more fundamental present constraint ---
not a demonstrated corpus-size threshold: there is no known related
language to map onto, the multiple-testing burden of candidate readings
goes uncorrected, and the sign inventory is mixed and partly lossy. We do
not claim additional data would be irrelevant --- a null can arise from
identifiability, low power, or corpus size alike. A lossy, non-injective
sign-to-sound channel need not admit a decipherment at any corpus size,
so we make no claim that ``there is enough text to pin a reading''; the
rigorous demonstration of \emph{that} constraint is the result the field
can use. The June-2026 AI-assisted ``Linear A cracked'' claim --- reading
the *301-anchored libation verb as an extinct Semitic form --- serves
throughout as a qualitative illustration of the failure modes the gate is
designed to detect, analysed from its public, archived presentation as a
claim, not a person, and neither reproducible nor held-out validated.
The artifact is open and archived (\S9.7).

\textbf{Roadmap.} Section~2 situates logos against the Braovi\'{c} et al.
taxonomy and the computational-decipherment lineage, including the
cognate-ID baseline. Section~3 specifies the harness --- canary,
deflation, literature index, and the pre-registration/held-out gate.
Section~4 develops the information-floor diagnostic and states its
identifiability limits. Section~5 works the June-2026 *301 claim as a
case study of the failure modes. Sections~6--8 report the pre-registered
nulls: morphology (\S6); metrology, phonology, and cross-script signals
(\S7); and contamination and sufficiency probes (\S8). Section~9 draws
the discussion, limitations, and conclusion.

\section{Related work}

\subsection{Position relative to the survey taxonomy}

The reference frame for computational work on the Bronze Age Aegean
scripts is \citet[\emph{Computational Linguistics}
50(2):725--779]{braovic2024}, the field's first survey focused
specifically on Aegean and Cypriot scripts; its only comparable
predecessor is \citet{ferrara2022}, and the broader
machine-learning-for-ancient-languages landscape is surveyed by
\citet[\emph{Computational Linguistics} 49(3)]{sommerschield2023}.
Braovi\'{c} et al.'s Section~4 enumerates fifteen ``major challenges'' ---
from unknown writing system, reading direction and language, through
small dataset, allographs and incomplete vocabulary, to unavailable
parallel data and hardware --- \textbf{none} of which concerns
overfitting, multiple testing, false-positive control, contamination, or
distinguishing a \emph{fitted} decipherment from a \emph{real} one. The
same gap runs through Ferrara and Tamburini's
\citeyearpar[\emph{Lingue e Linguaggio} XXI.2:239--259]{ferrara2022}
\S3.5, which reviews the cognate-matching lineage logos builds on ---
\citet{snyder2010} on Ugaritic$\leftrightarrow$Hebrew (the 29/30-sign
benchmark our canary calibrates against), \citet{bergkirkpatrick2011},
and \citet{luo2019, luo2021} --- and frames Linear A, via Gelb and
Whiting's typology \citep{gelb1975}, as a Type III script (unknown script
\emph{and} language) whose ``necessary validation'' is ``altogether
impossible,'' yet nowhere treats how a match is to be \emph{deflated} for
the multiplicity it was selected from. Two surveys, the same unfilled
axis. We do not score this as an oversight: their challenges are
properties of the object and its data, whereas the axis logos foregrounds
is a property of inference \emph{methodology} --- named-but-unfilled
rather than absent. logos supplies the machinery to police it (\S3).

Braovi\'{c} et al. do gesture at a related but separable difficulty:
evaluation. Because computational analyses of an undeciphered script have
``nothing to compare them against,'' they conclude (Sec.~6.3) that
``their evaluation remains an open problem.'' That is a distinct problem
from overfitting and multiplicity control, and logos addresses it
separately. Its held-out design is a concrete response --- the libation
formula recurs across a small set of peak-sanctuary and cave findspots,
permitting a leave-one-site-out cross-validation adjudicated mechanically
by \texttt{verdict.py} --- with the low-power caveat flagged at the
outset (\S1). We locate our components in the survey's own
method taxonomy (Sec.~7.1) and, where prior art exists, defer to it
rather than claim novelty.

\subsection{The current philological synthesis}

\citet{salgarella2025}, \emph{Writing in Bronze Age Crete: ``Minoan''
Linear A} (Cambridge Elements), by the author of the SigLA database, is
the most current authoritative one-author overview. It concludes that
Minoan is a language \textbf{isolate} --- it ``does not belong in any
known language family'' (the families used for comparison being
Indo-European, Semitic and Afro-Asiatic) --- typologically agglutinative
with a provisional Verb--Subject--Object order, and it records that ``no
bilingual text has so far been found'' in a corpus of roughly 2,500
inscriptions. Crucially, Salgarella ranks the etymological/comparative
method below combinatorial-contextual analysis and names
\textbf{digital/statistical analysis as ``the next desideratum in the
field.''} logos positions itself as exactly that desideratum --- but
under discipline, treating the isolate verdict as a central
\emph{identifiability} constraint rather than a target to be argued away.

\subsection{Prior computational attempts and the search trap}

The seminal automatic-decipherment result is \citet{luo2019}, which
recovers a character mapping from a lost script to a \emph{known related}
language via neural minimum-cost flow. It auto-deciphered
Ugaritic$\rightarrow$Hebrew ($+5.5\%$ over prior state of the art) and
translated 67.3\% of Linear B$\rightarrow$Greek cognates, but did not
address Linear A and would fail there by construction: with no known
cognate plaintext, the objective has nothing to optimise toward. This
cognate-language precondition is the empirical form of our
identifiability argument. The cognate-ID baseline logos actually runs is
\emph{not} a reimplementation of Luo's neural flow: it is a published
matcher, \citeauthor{tamburini2025}'s \citeyearpar{tamburini2025}
CSA\_OptMatcher (coupled simulated annealing), which reports the same
\citet{luo2019} accuracy metric. We adopt it as the non-neural baseline
against which any Linear A cognate claim must clear.

Against that disciplined baseline stand the in-domain instances of the
search trap: the dictionary-fishing programmes of \citet{eu2019} and
\citet{loh2020}, the latter explicitly a ``brute force attack'' on
dictionaries. These are the concrete Linear-A cases the multiplicity
correction is built to deflate. On the visual side, cross-script
comparison is \emph{not} our novelty: \citet{daggumati2019, daggumati2023}
already applied CNN+SVM cross-script similarity, and \citet{corazza2022}
(``Sign2Vec'') clustered signs unsupervised; our defensible contribution
is only the controlled image-versus-sequence \emph{asymmetry} under one
harness.

\subsection{The discipline lineage}

Two older strands supply the honesty machinery. \citet{packard1974},
\emph{Minoan Linear A}, constructed nine deliberately fictitious
decipherments as a null control --- and the genuine Ventris-value transfer
beat them only about 2:1 --- anticipating our permutation-null canary. The
statistical lineage is selective inference and multiple-comparison
correction \citep{bailey2014}, the correction for the number of
hypotheses $\times$ signs $\times$ roots $\times$ families tried that
turns a raw match-rate into a defensible one; its transfer from the
finance/backtest setting to script decipherment is specified in
Section~3. logos's empty cell is the intersection none of these occupy: a
falsification-first, decontaminated inference framework that couples a
published cognate-ID matcher \citep[reporting the Luo 2019
metric]{tamburini2025} with preregistration, multiplicity deflation and
mechanical held-out verdicts, released openly (\S9.7), with Linear A
worked as the honest null rather than a claimed crack.

\section{Methods: the discipline harness}

The harness is a comparison-layer pipeline that a reader can reimplement.
A candidate reading enters as a \emph{proposal} and passes through four
ordered stages --- pre-registration, deflation, decontamination, held-out
graduation --- under a cross-cutting persistence invariant: every decision
is computed from stored artifacts, never from a model's narration. Model-derived signals --- LLM theses, learned embeddings,
cognate-matcher assignments --- enter capped at confidence $\le 0.75$,
structurally unable to \emph{decide} a verdict; only
mechanically-verified held-out reads rank higher (\S1--2).

\subsection{Pre-registration}

Before any model is built or run, the predictions are frozen. The
morphology pre-registration fixes, for each hypothesis, the claimed affix,
its source, the falsifiable test, and the acceptance threshold, frozen
in-repository before each run so the model cannot be tuned to them after
the fact (the anchoring mechanics and the per-affix acceptance criteria
are given in Appendix~A.1). Outcomes are enumerated in advance --- CONFIRM, EXTEND,
REFUTE/NULL, or NO POWER --- so that a negative is a publishable result,
not a failed experiment. Any change is a new dated pre-registration,
never an edit.

\subsection{Deflation}

Internal fit on the derivation set is treated as no evidence. Every match
rate is deflated for everything tried: \texttt{n\_hypotheses} $\times$
\texttt{n\_signs} $\times$ \texttt{n\_roots}. The nulls are pipeline-level
--- a frequency-banded permutation null (after \citealp{packard1974}) and
the order-statistic bar of \S3.4 (implementation detail in Appendix~D).
Multiplicity is instrumented, not hand-passed, by a deduplicating
\texttt{SearchLog} (invariant 12; Appendix~D). The strong structural statistic,
\texttt{S\_morph}, earns credit only through a productivity gate --- an
affix must attach to $\ge$2 distinct residual stems, so a single formula
word copied across inscriptions scores zero (calibration in
Appendix~A.2, Table~\ref{tab:calib}).

\subsection{Decontamination}

A large open-weight model may have encountered longstanding published
readings by Gordon and Best; reproduction of indexed historical proposals
is therefore treated as potential shared-source contamination rather than
independent evidence. The June-2026 Di Mino claim is analysed separately
as a contemporary public case and is not assumed to be present in the
model's pretraining data --- it is indexed so that future claims and
models cannot reproduce it as discovery. Three mechanisms address this.

\textbf{The literature index} records published Linear A sign-value and
correspondence claims; any proposal matching the index is tagged
\texttt{literature\_match} and quarantined --- it may be tested but can
never count as discovery. The index partitions the sign inventory into
\texttt{L\_indexed} (referenced by $\ge$1 published proposal) and
\texttt{L\_not\_indexed} (untouched). Discovery claims may rest only on
\texttt{L\_not\_indexed} held-out success, since regurgitation can only
return what is in the literature. (The conservative seed, the adversarial
population, and the census are given in Appendix~C.5,
Table~\ref{tab:canary}.)

\textbf{The \texttt{L\_fake} canary} is a phonotactically-plausible
invented lexicon, calibrated to a real candidate language
(Appendix~C.5), generated after the model's release and never published
by us; exact-string collisions with real lexica are measured
(Appendix~C.5). Running the full pipeline
against it yields a \textbf{generator-conditional} false-positive floor:
any apparent lexical support it produces is spurious with respect to the
target historical-language hypothesis, and memorisation cannot
be invoked because there is nothing to memorise. A real candidate's match
must outperform this \emph{one calibrated class of negative controls} by
the corrected margin --- support means the candidate beat this
generator-conditional control family (here a first-order sign-bigram /
root-template generator), not that it defeated every possible false-language
process.

\textbf{LLM-ablation} compares the pipeline with and without the model's
proposals to isolate contamination (procedure in Appendix~C.1).

\subsection{The held-out graduation gate}

A reading graduates only mechanically, from held-out data, as a
conjunction of pre-registered clauses computed by \texttt{scripts/verdict.py}
--- never by the model that proposed the reading. The operative deflation
clause is an \textbf{order-statistic bar}: the observed held-out statistic
must exceed the \emph{expected maximum} of the multiplicity-null taken over
the number of trials the search was free to run (the E[max] over
\texttt{n\_trials} deflation), a \textbf{normal-order-statistic
approximation} in which \textbf{$n =$ \texttt{n\_trials} is the \emph{count of
candidates the search tried}, not an estimated number of independent
hypotheses}. The explicit formula and its Gaussian/independence assumptions
are given in Appendix~D; the error-control claim is carried not by the formula
but by the direct empirical calibration below. This is a direct
correction for selection across the hypothesis
space, and it fails closed on a degenerate null (Appendix~D). The framework yields
four primary evidential outcomes: \textbf{GRADUATE} (every clause holds and the
search multiplicity was instrumented), \textbf{REJECT} (a clause failed, or
the statistic sits below the corrected bar), \textbf{NULL\_PUBLISHED} (the
statistic is indistinguishable from the \texttt{L\_fake} null --- a
calibrated non-result, publishable as such), and \textbf{NO POWER} ---
the last assigned \emph{upstream}, by the pre-registered module outcome
when the experiment cannot adjudicate (\S3.1), before the gate sequence
runs; the gate itself additionally emits a fail-closed \textbf{INCOMPLETE}
status when search multiplicity is un-instrumented. These outcomes concern \emph{evidential
support, not truth}: because Linear A has no ground truth, REJECT / NULL /
NO POWER mean a reading is unsupported or insufficiently validated on
held-out data, \textbf{never that the underlying linguistic proposition is
false}. The verdict follows a fixed decision sequence
(\texttt{verdict.py}): all pre-registered clauses hold on an above-null
match $\rightarrow$ \textbf{GRADUATE}; else a match clearing the
substantive clauses but on an un-instrumented search $\rightarrow$
\textbf{INCOMPLETE}; else an observed statistic not exceeding
the null mean $\rightarrow$ \textbf{NULL\_PUBLISHED}; else $\rightarrow$
\textbf{REJECT}.
We calibrate the refusal rule directly: under the
tested best-of-100 random-map null --- an analyst trying 100 random
candidate maps and submitting the best --- 3 of 500 replicates graduated,
an observed false-graduation rate of 0.6\% (one-sided exact
Clopper--Pearson 95\% upper bound $\approx 1.54\%$), with 487/500 rejected
outright and 10 published as calibrated nulls (Appendix~D); this is an
empirical
error-control result conditional on that null-generating process, not a
universal guarantee. A second
pre-registered clause, \textbf{\texttt{generalizes\_to\_not\_indexed}},
requires the reading to recover literature-unseen (\texttt{L\_not\_indexed})
signs --- so reproducing already-published values cannot by itself
graduate a claim; the clause fails closed (its pre-committed thresholds
are in Appendix~D). Graduation is defined as validated \emph{novel}
discovery; independent held-out support for an already-published reading
is recorded as a distinct outcome and does not graduate. The remaining
pre-registered clauses (\texttt{L\_fake} separation, corrected margin,
order-statistic bar) are conjoined as specified in Appendix~D. These are
\emph{gate parameters} --- pre-registered
choices for this corpus, not universal constants. Validation is
leave-one-site-out, not ordinary \emph{k}-fold (Appendix~A.1). The gate is
AND-monotone --- clauses can only make it stricter --- so hardening never
lets a previously-failing reading graduate.

\subsection{Grade from persisted artifacts}

The invariant across every gate: verdicts are computed from persisted
rows, never from a model's account of its own outcome. The verdict writer
is grep-clean of any model call and is the sole writer of the verdicts
table.

\begin{table*}[t]
\centering
\begin{tabular}{@{}p{2.6cm}p{3.0cm}p{2.4cm}p{3.0cm}p{3.3cm}@{}}
\toprule
\textbf{Probe} & \textbf{Control / comparator} & \textbf{Real outcome} &
\textbf{Null / control floor} & \textbf{Verdict} \\
\midrule
Graduation gate (\S3.4) & --- & --- & best-of-100 random-map null &
3/500 false graduations (0.6\%; calibrated under the tested null) \\
Morphology (\S6) & Linear B 9/16 & Linear A 4/16 & \texttt{L\_fake}
bigram 6/16 & NO POWER \\
Segmentation (\S6) & --- & micro-F1 0.436 & random 0.389 & supported (gap
CI excludes 0) \\
Metrology (\S7.1) & planted 0.9 & 0.000 & permutation 0.113 & null
(agrees with the field) \\
Phonology (\S7.2) & planted (fires $\ge$ 2) & 0.080 / 0.140 & 0.065 /
0.201 & null (data-limited) \\
Image alignment (\S7.3) & font-swap survives & 0.410 & image chance
0.0135 & circular demonstration (capped) \\
Contamination (\S8.1) & model-free 0/40 & 14/40 (JA-NE) & 0/40 &
exploratory (public-exposure gradient) \\
\bottomrule
\end{tabular}
\caption{One harness, calibrated epistemic outcomes across probes.}
\label{tab:harness}
\end{table*}

The same falsification-first harness returns \emph{different} calibrated
verdicts across probes (Table~\ref{tab:harness}) --- the paper's
contribution in one view. Each number is detailed in the cited section.

\section{The information floor}

Before any decipherment method is graded, the corpus must be sized against
what it can identify. Two quantities matter: how much text exists, and how
many free parameters a reading spends against it. Confusing the first for
the second is the origin of the ``too little text'' folk claim --- and, we
argue, of its equally unearned inverse, the optimism that a big-enough
corpus makes a reading identifiable.

\textbf{The corpus.} The current authoritative synthesis
\citep{salgarella2025} gives an envelope of roughly 2,500 total Linear A
finds, of which about 1,400 are readable inscriptions, overwhelmingly
concentrated in the single LM IB destruction horizon (findplace and
dating detail in Appendix~C.4). That near-single horizon is why a
site-held-out split, not a chronological one, is the only realistic
partition. No bilingual text has ever been found
\citep[\S8]{salgarella2025}.

\textbf{Corpus sizing and the unicity computation (numeric detail in
Appendix~C.4).} The recurring sign-count confusion resolves into several
distinct sign-count denominators and two occurrence channels that must be
kept distinct (Table~\ref{tab:denoms}; Appendix~C.4). On the parsed
syllable-sign budget, Shannon's unicity distance sits everywhere far
below the corpus size (Table~\ref{tab:denoms};
\citealp[\S1--3]{salgarella2025}); we retain
this only as a toy-model \emph{precondition} diagnostic and explicitly
withdraw it as any sufficiency or refutation claim. Qualitatively, the
figure cannot settle identifiability either way (Appendix~C.4): unicity
distance presupposes a known plaintext-language oracle we do not have.
Corpus size, in short, cannot decide the question.

\textbf{What actually binds.} The results are consistent with
\emph{identifiability} being the more fundamental present constraint rather
than corpus size \citep[\S8]{salgarella2025}; we do not establish that more
data would be irrelevant. It has three parts: (i) the \textbf{absence of a
known cognate language} --- the map has a value-space but no anchored
target (Luo 2019's limit); (ii) the \textbf{multiple-testing / search
trap} --- that a unique consistent map may exist in principle neither
guarantees an algorithm finds it nor stops a cherry-picked false positive
across the parameters tried (invariant 8); and (iii) a \textbf{mixed
inventory} of syllabograms, logograms, and numerals that a real reading
must first separate. These, not the count of surviving signs, are what a
decipherment must overcome. Accordingly, when a candidate lexicon grows
without bound on a fixed occurrence budget --- free parameters exceeding
the floor --- the reading is underdetermined, and the framework must say
so.

\section{A qualitative audit of the *301 claim}

The framework of Section~3 is abstract until it meets a concrete claim. In
June 2026 a widely circulated ``Linear A cracked'' announcement
\citep{dimino2026} argued that Linear A records an extinct Semitic
language, reviving Gordon's hypothesis \citep{gordon1957, gordon1966}. Its
load-bearing anchor is a single reading: the recurring peak-sanctuary
libation formula word A-TA-I-*301-WA-JA, in which the Linear-A-only sign
*301 is \emph{itself} assigned a /na/-initial value, so the verb root
reads as West-Semitic N-W-Y ``to dwell'' (\emph{nawaya}) and is thereafter
matched to later Hebrew liturgy. We analyse this \textbf{as a claim, on its
public evidence} --- not the person --- from an immutable archived copy
carried in \artobj. It is a clean, self-contained instance of the failure
modes the gate is built to catch, and the three-part refutation below rests
wherever possible on the claim's \textbf{own published table} and on
primary sources \citep{davis2013, duhoux1992, thomas2020, salgarella2025,
steele2017} --- none advanced in logos's own voice, and the paper
introduces no new sign reading. We flag the meta-point once: this is a
literature-and-primary-source argument, not itself a logos held-out result
--- the graduation gate and \texttt{L\_fake} canary are validated as
instruments (unit tests; the known-answer Ugaritic$\leftrightarrow$Hebrew
surrogate) but have \textbf{not been run on this claim}. Scoring Gordon's
equations and the Di Mino lexicon against the \texttt{L\_fake} floor and
the order-statistic bar is directly runnable with the components of \S3 and
is the priority of the planned revision; that the refutation has not been
independently vetted by Aegean epigraphers is a disclosed limitation
(\S9.6), not a standard the work claims to meet.

\textbf{(a) The segmentation is not novel.} Root i-*301 plus affixation is
the reading of \citet{davis2013}; an independent segmentation by
\citet{thomas2020} reaches the same core and is tested head-to-head against
it (the current synthesis, \citealp{salgarella2025}, does not itself adopt
Thomas). On the Davis and Thomas segmentations, the affix-stacking is
evidenced by multi-attestation forms --- ta-na-i-*301-u-ti-nu (IO Za 6) and
a-ta-i-*301-de-ka (ZA Zb 3) --- that vary the material around a stable
i-*301 core. (These sigla are confirmed against Davis 2013: \texttt{TL Za
1}, \texttt{IO Za 6}, and \texttt{ZA Zb 3}.) The claim re-segments a
structure already established in the literature. Under the decontamination
layer this matters procedurally, not just historically: a correspondence
that reproduces a published segmentation is treated as potential
shared-source contamination,
tagged \texttt{literature\_match} and barred from counting as discovery
(\artref, \S C.1).

\textbf{(b) The gloss is contested and weaker.} The mainstream value of the
verb is ``give/dedicate'' --- Duhoux's \citeyearpar{duhoux1992} gloss,
adopted by \citet{davis2013} as a \emph{structural} reading that assigns
*301 \textbf{no phonetic value at all}, and endorsed by the current
synthesis \citep[\S7.2]{salgarella2025}. ``Dwell'' is a semantically
distant, unmotivated substitution. The point is not that ``give'' is
certainly right, but that a \emph{structural} reading needing no
sound-value for *301 already accounts for the data the phonetic claim is
invoked to explain.

\textbf{(c) The primary descriptions of this form's morphology do not
support the Semitic label.} The claim's own table segments the verb as
prefix-stacking and agglutinative --- [prefix A ``I''] [stem-morpheme TA]
[stem-vowel I] [root *301]. We deliberately do \textbf{not} rest on an
\emph{a-priori} typological contradiction here, and we retract the earlier
draft's claim that a prefix-stacked verb is ``internally incoherent'' with
a Semitic reading: Semitic verbal morphology includes productive prefix
conjugations, so a prefix on the verb is not by itself un-Semitic. The
point is empirical, not typological, and it comes from the primary
literature. \citet[p.~38 n.~13]{davis2013} characterises i-*301's
morphology as ``neither IE nor Afroasiatic'' --- and Semitic \emph{is}
Afroasiatic, so a decade earlier the primary description already found
\textbf{no Afroasiatic morphological signature} in this very form. (We rest
on what the footnote states --- the absence of an Afroasiatic signature ---
not on a stronger ``ruled out'' than a hedged footnote can bear.)
Independently, \citet[\S8]{salgarella2025} concludes Minoan is
agglutinative (after Duhoux 1978) and a language isolate. The claim's own
segmentation is thus consistent with the isolate typology, while the
specific Semitic assignment finds no morphological support in the primary
descriptions of the form.

\textbf{The free-parameter receipt.} The decisive point for identifiability
is that *301 has no cross-script anchor: in the claim's own table its
AB-number column is ``---'' and its Linear B correspondence ``(none)''.
Even the standard practice of reading Linear A homomorphs with their Linear
B values is a backward projection to be handled with care
\citep[already cautioned by Duhoux 1989, apud Davis 2013, p.~38
n.~11]{steele2017} --- and an A-only sign such as *301 has no homomorph to
project from at all, so nothing cross-script constrains its value. The
single sign on which the entire decipherment pivots is thus a \textbf{free
parameter} --- its value was assigned to fit the desired root, not
constrained by any independent evidence. This is precisely the case the
minimum-description-length / identifiability \emph{diagnostic} flags: a
value chosen to make one word read out carries no description-length saving
that survives the multiplicity of alternatives it was selected from (cf.
Section~3; the diagnostic is reported next to every claim, and the
operative graduation clause --- the order-statistic bar over exactly that
multiplicity --- is one a hand-fitted value cannot clear). For this reading
the obstacle is identifiability rather than corpus size --- no known cognate
language to constrain the assignment, uncorrected multiplicity, and a sign
whose value is chosen rather than anchored (a free parameter is unanchored
at any corpus size).

Read against Section~3, the claim instantiates two distinct failure modes
at once. It is \textbf{conclusion-driven model specification} --- the
table's \emph{structure} says agglutinative and prefixing, but the
\emph{label} says Semitic, and the label wins because the desired
conclusion (a Semitic prayer) is a fixed prior rather than an emergent
reading. And it is a \textbf{free-parameter / identifiability} failure ---
an unanchored sign-value tuned to the target, on a corpus with no known
cognate language to constrain it. logos indexes the reading (*301 = /na/)
precisely so a model cannot regurgitate it as discovery, and quarantines
rather than dismisses it. The lesson generalises: a match rate on a
hand-picked word, absent a committed prediction and multiplicity
accounting, is not evidence. The *301 claim is a qualitative illustration
of the failure modes the gate is designed to detect.

\section{Results I: morphology induction}

We test whether Linear A's pre-registered affixes and word divisions
reflect genuine morphology. The apparent affixation does not clear the
operative \texttt{L\_fake} bigram floor, so the verdict is NO POWER and
\texttt{has\_validated\_morphology} = False. The identical test
\textbf{validates on Mycenaean Greek (Linear B)}, whose longer words carry
uncontested inflection: it confirms affixation at rate 0.562 --- above the
same bigram floor ($\sim$0.30) --- with \texttt{has\_morphology\_power} =
True, so the control rules out a universally insensitive detector and
supports the reading that Linear A's shorter forms contribute materially to
its lack of power (Table~\ref{tab:morphrates}; Appendix~A.1). One held-out positive stands and survives
its own uncertainty: a DP-unigram segmenter recovers the scribe's own word
divisions at micro-F1 0.436 versus a random-boundary baseline of 0.389, a
gap whose site-clustered 95\% bootstrap CI [0.021, 0.099] excludes zero
(P(gap $>$ 0) = 0.998; Appendix~A.1). This
separates real recoverable structure from fitted morphology --- the paper's
core discipline.

\section{Results II: metrology, phonology, and the cross-script asymmetry}

Three further pre-registered probes exercise the harness against distinct
field claims. Each returns a calibrated null or a truth-capped positive;
together they map where the Linear A corpus does and does not carry
recoverable signal, and they are consistent with identifiability rather
than effort being the more fundamental present constraint.

\subsection{Metrology: a parse-coverage null that agrees with the field}

Direction D poses the accounting tablets as a constraint-optimisation
over the fraction-sign values (Appendix~B.1; Table~\ref{tab:metro}).
Held-out fraction balance
does not beat the shuffle null
($p = 1.0$), so the verdict is NULL; the only value the method determines,
the fraction sign J = 1/2, is the long-standing field consensus
\citep{bennett1950}, not a logos discovery. This is a null the field itself
reached \citep{corazza2021} --- the harness crediting a surviving value to
its originator rather than re-branding it.

\subsection{Distributional phonology: a data-limited null bounded by a
power curve}

Does a sign's distributional context (which signs sit beside it) carry
information about its phonetic value? Tested on the known AB signs, neither
the consonant class nor the vowel class clears chance (permutation
$p = 0.40$ / $0.86$). The verdict is a data-limited NULL: the sound path is
not recoverable from Linear A alone, and no sound value is imputed for any
sign (Table~\ref{tab:phon}; Appendix~B.2).

\subsection{The cross-script asymmetry: image alignment recovers shared
anchors, but the recovery is circular by construction}

Do image embeddings recover the cross-script A$\leftrightarrow$B
correspondence that sequence and cognate co-occurrence cannot? Classical
shape descriptors reach direct-NN \textbf{0.410} on held-out
A$\leftrightarrow$B anchors while sequence and cognate co-occurrence stay
at chance. But the recovery is circular by construction: the
A$\leftrightarrow$B anchor values were themselves assigned by
twentieth-century epigraphers on sign-shape similarity, so this is a
shape-genealogy engineering demonstration, truth-layer-capped $\le$
\textbf{0.75}, not a positive decipherment claim (\S9.6;
Table~\ref{tab:xscript}).

\section{Results III: contamination probes}

One failure mode is invisible to the false-positive / multiplicity axis of
\S7: a decipherment that \emph{looks} novel but silently reproduces
published readings a language model has memorised. This section reports the
two decontamination probes that target it --- both run as method
demonstrations, returning completed (if deliberately bounded) findings.
Throughout, where a clean number was not obtainable we report the
partial-null and the confound rather than a headline.

\subsection{LLM-ablation: does a large open-weight model regurgitate the
literature?}

In the LLM-ablation, \texttt{qwen2.5:72b} reproduces famous
published readings --- Gordon's JA-NE = ``wine'' in 14/40 seeds --- that
the model-free edit-distance baseline never reaches (0/40), yet never
reproduces obscure vessel equations (0/40; Table~\ref{tab:repro}).
Verdict: a reproduction
gradient consistent with differential public exposure --- though the
experiment does not isolate memorisation from correlated item-level
differences --- the memorisation confound the discipline layer must gate
out.

\subsection{\texttt{L\_not\_indexed}: indexed versus not-indexed
abstention}

Does the model's behaviour on \emph{published} signs generalise to signs
unseen in the indexed literature? \texttt{qwen2.5:72b} valued the
memorisable known signs in $\sim$20/40 seeds but abstained $\sim$97\% of
the time on the literature-unseen *-series signs ($\sim$1/40;
Table~\ref{tab:abstain}) --- an
abstention asymmetry consistent with differential public exposure (the
\S8.1 caveat applies). With no ground truth on the not-indexed signs,
\textbf{consistency is not correctness}: it bounds how the model behaves,
never whether it is right.

\section*{Registered extension: a known-answer sufficiency curve}
\addcontentsline{toc}{section}{Registered extension: a known-answer
sufficiency curve}

Is a Linear-A-scale corpus large enough to recover a decipherment if a
known cognate existed? Because every benchmark is handed a real cognate
signal Linear A lacks, any recovery is an \textbf{optimistic benchmark,
conditional on a known cognate relationship}: at-floor at Linear-A scale
$\Rightarrow$ information-insufficient even given a cognate; above-chance
$\Rightarrow$ identifiability (cognate-absence, uncorrected multiplicity)
rather than corpus size would be implicated; neither branch is a
decipherability claim. The completed 2{,}000-step curve (near-converged:
Luvian$\to$Hittite 0.462 vs.\ published 0.475, Phoenician$\to$Ugaritic
0.760 vs.\ 0.805; every recovery a \emph{lower bound}) splits the
branches at the sweep's Linear-A-scale locator (600--700 distinct
word-forms, syllabic normalization): Linear~B$\to$Greek, the primary
analog, stays at its chance floor ($0.044\pm0.023$, $n{=}276$; floor
$\approx0.004$), while Cypriot$\to$Greek sits marginally above
($0.066$, $n{=}346$). Under-convergence keeps the at-floor branch
non-definitive; full per-size and per-cell artifacts are deposited
(Appendix~C.3).

\section{Discussion, limitations, and conclusion}

\subsection{The discipline is the contribution}

The deliverable of this work is not a reading of Linear A. It is a
\emph{harness} --- a literature index with an indexed versus
literature-unseen partition, a fabricated-language control corpus,
model-ablation contamination checks, multiple-comparison correction, and
pre-registration, all graded mechanically from persisted artifacts ---
together with a body of pre-registered, calibrated nulls, each with its
borderline case named. The stratified morphology re-run is the template: at
the pre-registered $\alpha = 0.01$, no pre-registered affix is cross-stratum
stable above the bigram floor, robust to the genre-bucketing choice
(Tables~\ref{tab:strata}--\ref{tab:sens}; Appendix~A.1). A null that
names its threshold, its sensitivity, and its
near-miss is a stronger referee object than a flat negative.

The axis this harness targets is the \textbf{false-positive / multiplicity
axis} --- the mechanism by which a determined analyst ``translates'' almost
anything on a small corpus --- a property of inference methodology, not of
the object or its data (\S1--2).

\subsection{Two named failure modes for the methods literature}

\textbf{(i) Narrative capture.} Structure says X; the label says the
desired Y; the label wins. In the June-2026 *301 claim the published table
segments an agglutinative, prefix-stacking verb while the assigned family
label is Semitic, and the desired output operates as a fixed prior, so
the conclusion conforms to the prior rather than emerging from
constrained readings (\S5). This
is distinct from statistical overfitting and must be named separately.
Notably, the mismatch is empirical --- between the label and the primary
descriptions of the form (\S5) --- not an a-priori typological
impossibility.

\textbf{(ii) Free parameters exceeding the information floor.} A lexicon
reported as ``408 terms and counting'' on a fixed corpus has more degrees
of freedom than the data can identify; a decipherment whose lexicon grows
weekly is underdetermined by construction. The minimum-description-length /
unicity \emph{diagnostic} (\texttt{k $\le$ U\_floor}) is built to flag
exactly this (it is reported next to every claim rather than enforced as a
graduation clause; \S4): the single sign the whole
reading pivots on is the textbook free-parameter case (\S5), its value
assigned to fit the desired root; the refutation rests on the claim's own
published table, archived in \artobj.

\subsection{The post-hoc evaluation sequence}

These failure modes trace to sequencing. The *301 claim publicly commits to
a post-hoc evaluation sequence: declare ``officially deciphered,'' then vet
later. logos's documented inverse is \textbf{pre-register $\rightarrow$
held-out $\rightarrow$ locked artifact} (\S3.1, \S3.4; the held-out
libation-formula corpus and its low-power caveat are in Appendix~A.3).

\subsection{Identifiability appears the more fundamental present constraint}

We take care not to write a defeatist framing, and equally not to
over-claim from corpus size: the toy-model unicity distance
($U \approx 204$--$415$ signs; Table~\ref{tab:denoms}, \S4) is a
\emph{precondition} diagnostic
only, explicitly withdrawn as any sufficiency or refutation claim ---
``$U \ll N$'' does \textbf{not} license the statement that there is enough
text to pin a decipherment (\S4; Appendix~C.4). The
results are consistent with \emph{identifiability} being the more
fundamental present constraint --- the absence of a known cognate language
to map to, an uncorrected multiple-testing search burden, and a mixed sign
inventory --- rather than a demonstrated corpus-size threshold; a null can
arise from identifiability, low power, or size alike, and we do not
establish that more data would be irrelevant.
\citet[\S8]{salgarella2025}
reaches the field-level version of this by a different route (\S2.2).

\subsection{Study-wide error budget}

Multiplicity discipline operates at two levels, and we are explicit about
both. Within a single probe, matches are deflated for everything
searched (\S3.2--3.4). Across the study, the unit of multiplicity is the
\textbf{per-probe pre-registration}: the \textbf{five executed probes}
(morphology, metrology, phonology, cross-script image alignment, and
contamination) each carry their own pre-committed $\alpha$, and every probe
is reported regardless of outcome. \textbf{Error control is defined within
each pre-registered probe; we claim no global family-wise guarantee across
probes.} Nothing in the battery is read as a decipherment.

\subsection{Limitations}

The nulls are honestly bounded, not decisive. A dependence-adjusted
trial-count correction of the morphology z-test is deferred; the deferral
would only strengthen the null (Appendix~A.1). No held-out
\emph{decipherment} has
been achieved, and the study reports \textbf{no positive decipherment
claim} --- the deliverable is calibrated absence, not a reading. Nothing
here claims a rejected reading is \emph{false} (\S3.4): every REJECT /
NULL in this study bounds \emph{evidential support}
on held-out data, not truth. The one probe that recovers an above-chance
number, cross-script sign-image alignment, is a font-robust but circular
\emph{engineering demonstration} (a font-swap control refutes the
typeface-artifact worry; \S7.3, Appendix~B.3), truth-capped at
$\le 0.75$, never a reading. No
phonetic claim is made. The philological points in \S5 rest on the
published literature cited there rather than on independent
Aegean-epigraphic
review, which we disclose as a limitation; because those points are cited
rather than original, the argument does not depend on any new sign reading
of our own. Chronological cross-validation is largely foreclosed by the
single-horizon corpus (\S4; Appendix~A.1). And the corpus itself
--- small, partly lossy, with no bilingual yet found
\citep[\S8]{salgarella2025} --- sets the ceiling.

\subsection{Conclusion and reproducibility}

The framework withholds validation when a proposed reading does not clear
the pre-registered gates (\S1, \S3); the contribution is the
framework and the calibrated null itself --- the instrument and the honest
limits of what the corpus can support --- not a reading, and the calibrated
null is a safeguard against unsupported claims.
\ifanon
The code, pre-registrations, and all findings are open under the MIT
licence and available as an anonymized artifact (code and pre-registrations
released on acceptance).
\else
The code, pre-registrations, and all findings are open under the MIT
licence and archived on Zenodo (concept DOI 10.5281/zenodo.21087572;
repository \url{https://github.com/papadopouloskyriakos/logos}).
\fi
The licensed epigraphic corpora are not redistributed; each source's
licence and retrieval path is documented (data-provenance; Software \&
data note), so the harness is reproducible against a
locally obtained corpus. This work is independent and not affiliated with,
or endorsed by, any of the cited projects or their authors.

\bibliographystyle{acl_natbib}
\nocite{*}
\bibliography{refs}

\appendix

\section{Morphology (replication detail)}

\paragraph{A.1 Morphology: rates, floors, word-length, stratification,
sensitivity.} The probe is pre-registered, unsupervised morphology and
segmentation induction over the full Linear A corpus
(Table~\ref{tab:denoms}). Target claims are the recurring affixes and
positional grammar of the Davis/Thomas/Duhoux--Val\'{e}rio tradition:
prefixes and suffixes (\texttt{i-}, \texttt{a-}, \texttt{-te/-ti}, and
Davis's \texttt{i-*301} verb stem). Pre-specification anchoring (\S3.1):
pre-specifications were frozen in-repository before each run (commit
history); the external immutable deposit
\ifanon
([artifact DOI withheld for review])
\else
(Zenodo concept DOI 10.5281/zenodo.21087572)
\fi
certifies the current state of the registrations, not temporal priority.
Each pre-registered affix
(drawn from \citet{thomas2020} and Duhoux/Val\'{e}rio) must recur on
$\ge$2 distinct stems across independent inscriptions, above a
within-form permutation null, and survive multiplicity correction at a
target of \emph{p} $< 0.01$ corrected. Every candidate was graded against two
null floors: a within-word sign-order permutation (shuffle) floor, and an
\texttt{L\_fake} floor built from a first-order sign-bigram (Markov) corpus
with zero lexical overlap with the real inscriptions.

\emph{Pooled run.} Of the 16 pre-registered affixes, 4 beat the within-form
permutation null and the deflated multiplicity bar (\texttt{i-*301},
\texttt{a-}, \texttt{i-}, \texttt{-te/-ti}), clearing the shuffle floor
decisively (rates and Wilson intervals in Table~\ref{tab:morphrates}).
The apparent standout \texttt{i-*301} sits far above the permutation null
--- but that figure is computed against the weaker within-form permutation
null only, and it does not survive the operative \texttt{L\_fake} bigram
floor, under which the whole set fails: the fabricated Markov corpus
confirms the same affixes at a \emph{higher} rate than the real corpus
(Table~\ref{tab:morphrates}). The outcome is classified NO POWER, not
evidence against morphology; \texttt{has\_morphology\_power} is False and
the module reports the null (the Linear-B positive control that validates
the detector is reported below).

\emph{Word-length distribution.} Linear A words are short
(Table~\ref{tab:morphrates}), and on such words ``morphology'' and
``bigram order'' are statistically indistinguishable. This reconciles
with Fuls's 3.3-sign figure once the denominator is matched:
\citet[57, \S6.5]{fuls2015} computes the mean over a word \emph{list} of
``554 complete sign sequences, where duplicates are reduced to one''
(GORILA; \citealp{godart1985}) --- distinct word \emph{types}, not
tokens; restricting our corpus to distinct multi-sign types likewise
recovers 3.07. Fuls reads the 3.3-sign length as evidence that Minoan is
\emph{polysynthetic}; our direct test finds the apparent affixation
bigram-reproducible with no power, and we report the null. The direct
positive control --- the identical bigram-floor test on Mycenaean Greek
(Linear B), whose longer words carry uncontested inflectional morphology
--- was RUN on the D\=AMOS corpus (parsed by the same
\texttt{\_damos\_wordforms} extractor as \S7.3) against a pre-registered
16-affix Mycenaean panel \citep{ventris1953}. It fires
(Table~\ref{tab:morphrates}), carried by the securely-grammatical
suffixes -jo, -o-jo (gen sg -oio), -de (allative), -qe (enclitic
``and''), -o-i (dat pl), so \texttt{has\_morphology\_power = True}. The
identical test that returns NO POWER on Linear A's short words recovers
Mycenaean morphology on longer ones (a sensitivity check: an ad-hoc
parser yielding mean 2.19 signs/word did \emph{not} fire --- power tracks
word length exactly as the short-word argument predicts), which rules out
a universally insensitive detector and supports the interpretation that
Linear A's shorter forms contribute materially to its lack of power
(\texttt{scripts/\allowbreak comparison/\allowbreak linb\_morphology\_control.py}).

\emph{Segmentation positive.} Under leave-one-site-out cross-validation (52
sites; not k-fold, given formulaic dependence), a DP-unigram
(Goldwater-style) segmenter recovers the scribe's own word divisions above
a random-boundary baseline; a site-clustered bootstrap (resample the 52
sites with replacement, 4,000 draws) puts the real-minus-random gap above
zero, and the paired gap, not the wide marginal per-segmenter CIs, is the
operative test (Table~\ref{tab:morphrates}). Morfessor over-segments the
short syllabic corpus and is reported as not usable rather than hidden
(Table~\ref{tab:morphrates}).

\emph{Stratified re-run.} The strata reconcile exactly to the
1,341-inscription corpus total (Table~\ref{tab:strata}). The
99-inscription ``libation'' \emph{genre} stratum is not the same object
as the libation-\emph{formula} corpus (\S A.3). Following
\citet{salgarellajudson2024}, chronological cross-validation was explicitly
declined (988/1,341 $\approx$ 74\% single-horizon LM IB), and site --- not
scribal hand, which is not individuated for Linear A
\citep{salgarella2019} --- is the independence unit, gated by presence at
$\ge$2 distinct sites. The result is \texttt{has\_validated\_morphology} =
False (\S9.1). Stratification also showed
that even \texttt{i-*301} is bigram-reproducible in the repetitive libation
register, where the Markov corpus manufactures the apparent verb
morphology; Davis's slot grammar is not refuted, merely not separable from
sign bigrams here. Where formula-bearing sites drive a signal, power is
intrinsically low: a leave-one-site-out CV over so few formula findspots
(an effective $\sim$5-fold site partition) cannot resolve a modest effect.

\emph{Sensitivity, disclosed.} The honest 2$\times$2 over genre bucketing
$\times$ $\alpha$ threshold is Table~\ref{tab:sens}; the doubly-relaxed
corner's \texttt{i-} (locative) and \texttt{-te/-ti} (ablative) are
precisely the Duhoux/Val\'{e}rio H3 nominal affixes that
\citet{salgarella2025} \S8 independently endorses --- the named borderline
candidates, fragile to \emph{both} the $\alpha$ threshold and a defensible
genre choice (the $\alpha = 0.01$ null is itself robust to the bucketing).
The present analytical approximation uses attempted-candidate
count rather than a dependence-adjusted effective trial count --- a
recorded, deferred limitation whose correction would only strengthen the
null.

\paragraph{A.2 \texttt{S\_morph} detector calibration (\S3.2).} After the productivity-gate fix (\S3.2 --- an affix must attach to
$\ge$2 distinct residual stems), \texttt{S\_morph}'s measured false-positive
rate on within-form-shuffled corpora is consistent with the nominal
one-sided $z \ge 2$ rate within sampling error, and its planted-affix
power is a saturated point, not a full sensitivity curve
(Table~\ref{tab:calib}; contrast the phonology power curve, \S7). This is
an \textbf{indicative calibration check, not the operative gate}: 6/200
is too coarse to certify an $\alpha = 0.01$ threshold on its own, and the
gate applies the order-statistic bar over the far larger permutation and
candidate counts the deflation machinery tracks (\S3.4).

\paragraph{A.3 Held-out design: the libation-formula corpus (\S9.3).} A
leave-one-site-out split over so few formula sites (\S1) has
low statistical power --- and site-level CV is in any case necessary, not
sufficient, since within-site scribal-hand correlation must also be controlled
\citep{salgarellajudson2024}. This formula corpus should not be conflated with
the 99 stone-vessel inscriptions of the ``libation'' genre stratum (\S6): they
are different objects.

\section{Metrology, phonology, and cross-script (methods and controls)}

\paragraph{B.1 Metrology.} Direction D parses HT (LM I) documents into
balance units \texttt{[recipient][commodity logogram][integer][fraction
signs]} and solves the fraction-sign values by RANSAC over exact-\texttt{Fraction}
Gaussian elimination (robust to the unbalanced-tablet outliers that defeat
least-squares), requiring line items to balance the preserved KU-RO totals;
validation is grouped k-fold by document (no tablet straddles train and
test) against a permutation null. Results are in Table~\ref{tab:metro}: of
$\sim$7 fraction-bearing constraints only one is satisfied by the solved
system, the sole determinable value is J = 1/2, held-out fraction balance
is not separated from the permutation null, and the synthetic positive
control (planted values, all recovered) separates. The null traces to sparse
automated parse coverage (multi-commodity blocks, illegible number glyphs,
KI-RO deficit sub-lists, editorial restorations); \citet[p.~2]{corazza2021}
abandoned the balance-to-totals method for the same reason. The surviving J
= 1/2 (\S7.1) gains Schrijver's
\citeyearpar{schrijver2014} three further tablets (HT Zd 156, PE 1, HT
123a): at least four corroborating documents.

\paragraph{B.2 Distributional phonology.} Track 2 asks whether a sign's
distributional context --- which signs sit beside it --- carries
information about its phonetic value, measured on the known \texttt{AB}
signs whose Linear B sound values can be borrowed. Features are
PPMI-weighted, L2-normalised left/right co-occurrence counts; labels are
the Linear-B-transferred (C, V) parsed from the GORILA name, so features
never see the labels. The test is leave-one-out 1-NN against a
label-permutation null, \v{S}id\'{a}k-deflated over the two tests; 50 of
$\sim$55 AB signs met the $\ge$3-occurrence threshold.

Per-test results are in Table~\ref{tab:phon}. The honest verdict is
\emph{data-limited},
not \emph{no bridge exists}, because the positive control disciplines the
reading: the positive-control power curve over planted strengths shows
the test firing at strength $\ge$ 2 but not at 1. The test
works but is not infinitely sensitive at n = 50 / 13 classes, so we cannot
distinguish ``no distribution$\rightarrow$phonetics bridge'' from ``too few
signs to detect a weak one'' (\S7.2, \S B.3).

\paragraph{B.3 Cross-script image alignment.} The palaeographic probe
renders 88 Linear-A and 74 Linear-B glyphs and tests held-out recovery of
shared A$\leftrightarrow$B anchors. Classical descriptors (HOG + Hu moments
+ shape-context) sit far above image chance, and a from-scratch I-JEPA is
excluded from the claims as over-parameterized for $\sim$90 glyphs
(results and controls in Table~\ref{tab:xscript}). Two controls decide how
to read the image-vs-sequence asymmetry (\S7.3), and together they
\emph{downgrade} it from a positive to a demonstration.

\emph{First, a font-swap control.} Rendering both scripts in one typeface
(Aegean, George Douros) risks making the alignment a single-designer
house-style artifact rather than a fact about the signs. We re-render Linear
A in Noto Sans Linear A and Linear B in Noto Sans Linear B --- two
independent type designers, each font covering only its own Unicode block
(Noto-A carries zero Linear B glyphs and vice-versa), so no single hand
draws both scripts. The recovery survives and is if anything stronger
(Table~\ref{tab:xscript}; both legs $\gg$ chance): the font-artifact
hypothesis is \emph{refuted}, and ``distances reflect shape, not font'' is
now demonstrated rather than asserted
(\texttt{scripts/\allowbreak palaeo/\allowbreak font\_control.py}).

\emph{Second, and decisively, the circularity the font control cannot
remove.} Linear B was historically adapted from Linear A, and the
A$\leftrightarrow$B transcription \emph{values} that define the anchor set
were assigned by twentieth-century epigraphers largely on the basis of
sign-shape similarity, so the recovery re-derives a shape-defined
correspondence (\S7.3): it confirms a known palaeographic fact (borrowed
signs look alike) and validates a font-robust descriptor pipeline, but it
reads no Minoan. It remains truth-layer-capped in any case --- the
\citet{ferrara2021} ``Archanes Formula'' hard negative (\texttt{CH 095}
resembles \texttt{AB 60}/\emph{ra} but is not its phonetic counterpart)
shows a plausible shape match can be phonetically wrong, so a shape match is
a $\le 0.75$ signal (invariant 5), never a reading. The honest asymmetry is
thus: the image path re-derives a shape-defined mapping (expected, circular,
capped), while the sequence and cognate paths \textbf{remain at chance even
after ingesting the full D\=AMOS Mycenaean corpus} --- 13,562 wordforms,
$\approx$15$\times$ the earlier cognate file, 56 anchors
(Table~\ref{tab:xscript}) --- while the \emph{positive control
recovers}, so the harness works and the null is a real null: the
A$\leftrightarrow$B phonetic correspondence is simply not carried by
cross-script sign co-occurrence, even at full Linear-B scale
(\S B.2). And the A-only signs on which
positive-decipherment claims would pivot have no cross-script
anchor at all (\S5), so images cannot impute them even in principle.

\section{Contamination and sufficiency (detail)}

\paragraph{C.1 LLM-ablation reproduction gradient.} The ablation
(\S3.3) is the remaining gate on the contamination number, matching a
model's cognate glosses against the indexed lexical claims. We ran it
against \texttt{qwen2.5:72b} --- a large open-weight model, a far
stronger memoriser than a local 12B --- on a rented H100, with an
NW-Semitic candidate and a Hebrew skeleton lexicon for the model-free arm
(scale and quantities in Table~\ref{tab:canary}) (\artref). The
first-pass contamination rate is an \textbf{artifact, reported as such}:
the mechanical arm never shares a literature-matching correspondence with
the LLM (their value vocabularies barely overlap), so the rate is forced
to 1.0 by construction and is not reported as a contamination result
(Table~\ref{tab:canary}; \artref).

\emph{Generation settings (replication).} Checkpoint \texttt{qwen2.5:72b}
as served by Ollama (Ollama's default \texttt{q4\_K\_M} quantization;
Ollama version not recorded), on an ephemeral vast.ai H100 SXM 80\,GB
instance (torn down after the run). One \texttt{/api/generate} call per
seed, stream off, timeout 180\,s; temperature 0.8; sampling seed = the
replicate seed; top-p, top-k, maximum output length, and stop conditions
left at server defaults (not recorded). Each
seed's 40 held-out forms are drawn deterministically
(\texttt{numpy.default\_rng(seed)}) and rendered
as canonical sign names in a fixed JSON-request prompt
(\texttt{\_PROMPT\_TEMPLATE} in
\texttt{scripts/\allowbreak comparison/\allowbreak ablation.py}). Parsing:
the JSON \texttt{partial\_map} is extracted; keys not naming a sign shown
to the model are dropped (the proposer cannot invent signs); a dead or
garbled call contributes an empty proposal without crashing the batch
(fail-closed). The gate-2 rerun
prepends the 7 indexed probe forms to every seed. Per-call logs:
\texttt{runtime/ablation/llm\_calls.jsonl}; harness commits
\texttt{0582966} (ablation) and \texttt{9cd6d04} (gate-2).

The gate-2 rerun, reconciled to a common consonant-skeleton vocabulary,
yields the honest deliverable: the per-item reproduction gradient over 40
seeds (Table~\ref{tab:repro}; \artref). An
external red-team correction split \texttt{KU-RO} between the
Semitic-root route (\emph{kull}) and the general-knowledge meaning
``total'' (confirmed by tablet arithmetic), so \texttt{KU-RO} = ``total''
is \textbf{excluded} as a witness and \texttt{KU-RO} = \emph{kull} is
downweighted (Table~\ref{tab:repro}; \artref). The surviving gradient ---
famous readings reproduced, obscure ones not --- is what the ablation
detects (\S8.1). The gate-2 rate itself remains confounded by the baseline's
representational ceiling and is not reported as a contamination measure
(Table~\ref{tab:canary}).

\paragraph{C.2 \texttt{L\_not\_indexed} abstention asymmetry.} The
\texttt{L\_not\_indexed} probe shows a sign in real context (\S8.2); the
not-indexed signs are absent from the indexed
sources, hence not memorisable from them (never ``impossible to have seen''
in any absolute sense) (\artref). Across 40 seeds, \texttt{qwen2.5:72b}
abstains $\sim$97\% of
the time on the untouched signs while confidently valuing the memorisable
ones (Table~\ref{tab:abstain}). The apparent known$-$not-indexed
consistency gap is coverage-confounded and was flagged, not asserted: the
matched-\emph{n} robustness check returned \texttt{no\_power}, so the
rigorous comparison could not run (Table~\ref{tab:abstain}).
The interpretable signal is the abstention asymmetry itself, read under
the standing rule of \S8.2 (\artref).

\paragraph{C.3 Information-sufficiency CSA sweep (design and deposit).}
The sweep downsamples known-answer benchmarks --- Linear B /
Mycenaean Greek as the primary syllabic analog, sourced through D\=AMOS,
and further syllabic corpora --- and measures held-out cognate-ID
recovery against per-size permutation floors, using the published
cognate-ID matcher of \S2.3 \citep{tamburini2025}, validated in the
artifact against its published Table-3 values (\artref). The reading
rules and the reported curve are in the body (Registered extension); the
full per-size report and per-cell artifacts are deposited in the
repository (\texttt{results/csa/}).
The phonological / sound path is no longer data-deferred: even the full
ingested D\=AMOS corpus leaves the cross-script distributional alignment
at chance (\S7.3, \S B.3).

\paragraph{C.4 Information floor: corpus denominators and the unicity
computation (\S4).} The corpus envelope, the four
sign-count denominators, the two occurrence channels, and the measured
unicity range are tabulated in Table~\ref{tab:denoms}
\citep[Schoep 2002;][\S1--3, \S6]{salgarella2025}; the envelope spans ca.
1800--1450 BCE (Middle Minoan II--Late Minoan IB; \S4). Each denominator
must be labelled where it is used (Table~\ref{tab:denoms}). The channel through
which structure is observed is a \emph{sign-occurrence} count, not a
document count or a sign inventory (invariant 7): the broader $\sim$7,400
occurrences count every sign type, while the narrower
$N \approx 5{,}792$ parsed syllable-signs --- after logograms and
numerals are stripped --- is what the unicity computation consumes. Two
further properties tighten the budget: Linear A words are short and
formulaic (the six-position libation-formula template;
\citealp[\S7]{salgarella2025}), and a single language/horizon dominates.

On the logos-normalized corpus, Shannon's unicity-distance criterion
yields $U \approx 204$--$415$ signs across the whole plausible $(V, P)$
parameter space, everywhere far below $N$ (Table~\ref{tab:denoms}) ---
retained only as a toy-model precondition diagnostic (\S4). Unicity
distance presupposes a \emph{known plaintext-language oracle}; because we
estimate the redundancy $D$ from the ciphertext itself, that oracle
collapses into Linear A up to relabelling, so the computation
demonstrates only measurable recurrent structure, not that any reading is
identifiable; and a lossy, non-injective sign$\rightarrow$sound channel
may possess no unicity distance at all (\S9.4).

\paragraph{C.5 Literature index and \texttt{L\_fake} canary (\S3.3).} \emph{Index seed and adversarial population.}
The seed is deliberately conservative --- a sign wrongly omitted is wrongly
granted \texttt{L\_not\_indexed} status, the dangerous direction --- and
every verifiably-published disputed entry is indexed to \emph{catch}
regurgitation, never as an
accepted reading. Population is adversarial: of the seeded West-Semitic
lexical proposals, five are Gordon's equations and a-sa-sa-ra-me = ``oh
Asherah!'' is \citet{best1981}, not Gordon; provenance-unverifiable
candidates --- \texttt{qa-pa} = kappu and a Semitic \texttt{ki-ro}, whose
attribution to Gordon is not supported by the primary account
\citep{rendsburg1996} --- were excluded rather than indexed. Exposure
indexing is independent of evidential validity --- a false public claim
can still contaminate. Generated from the
seed (invariant 12), the current census (Table~\ref{tab:canary}) is a
deliberately conservative, non-exhaustive seed, not yet the full \S C.1
literature census.
\emph{\texttt{L\_fake} calibration.} The canary is calibrated to a real
candidate language's phoneme frequencies, root-template structure, and
lexicon size (\S3.3). Its own divergence from the calibration
targets is reported, never minimised; with a Semitic root-template sampler
and an external Hebrew (ETCBC/bhsa) reject set, the root-template
total-variation distance fell markedly (Table~\ref{tab:canary}), and any
residual real-lexicon collision is measured and published rather than
asserted to be zero.

\section{The order-statistic graduation bar: formula, assumptions,
calibration}

The operative deflation clause (\S3.4) is a normal-order-statistic
\textbf{approximation} to the expected maximum of the multiplicity-null over
the attempted candidates. For a null of mean $\mu_0$ and standard deviation
$\sigma_0$ and $n$ attempted candidates,
\[
  \resizebox{0.95\columnwidth}{!}{$\displaystyle
  E[\max] \approx \mu_0 + \sigma_0\bigl[(1-\gamma)\,\Phi^{-1}(1-\tfrac{1}{n})
  + \gamma\,\Phi^{-1}(1-\tfrac{1}{ne})\bigr]$}
\]
with $\gamma = 0.5772\ldots$ the Euler--Mascheroni constant --- the
blended expected-maximum expression used by \citet{bailey2014}; the
general order-statistics treatment is \citet[\emph{Order Statistics}, 3rd
ed.]{david2003}. The
pipeline exposes the bar as
\texttt{expected\_max\_order\_stat}($\mu_0, \sigma_0, N$); the
deduplicating \texttt{SearchLog}
(\S3.2, invariant 12) reports the exact distinct-candidate count
\texttt{n\_trials} a scanner produces. The bar fails closed on a degenerate null
($\sigma_0 \le 0 \Rightarrow E[\max] = \mu_0 \Rightarrow$ no graduation).
The \texttt{generalizes\_to\_not\_indexed} clause (\S3.4) requires recovery
above a pre-committed threshold and on a minimum number of distinct
not-indexed signs, and fails closed when no threshold or no supported
not-indexed sign is present. The full clause conjunction: the candidate
must separate from the \texttt{L\_fake} negative-control distribution,
its corrected margin must be cleared, and the observed statistic must
beat the order-statistic bar over the Packard/random-lexeme multiplicity
null.

\textbf{Assumptions and the direction of their error.} The approximation is
exact only under (i) an approximately Gaussian null and (ii) independence
among the $n$ attempted candidates. Neither holds exactly. (i) The
permutation / \texttt{L\_fake} nulls are not perfectly Gaussian; the blended
Euler-$\gamma$ form is a first-order correction, not a guarantee. (ii) The
candidates a real search tries are typically \textbf{positively correlated}
(overlapping sign-value assignments). Under a common-variance Gaussian
comparison model, positive dependence lowers the expected maximum relative
to independence (a Slepian-type comparison; \citealp{slepian1962}), so
using the
raw attempted-candidate count $n$ in place of a smaller effective
independent-trial count places the bar \emph{above} the correlated truth;
this does \textbf{not} establish conservatism for arbitrary heterogeneous
or non-Gaussian dependence --- the empirical best-of-100 calibration
therefore carries the operational error-control claim. We use the
raw deduplicated candidate count; a dependence-adjusted effective count is a
recorded, deferred refinement (\S9.6) whose correction would only strengthen
the nulls.

\textbf{Calibration and its interval.} Because the analytic bar rests on
assumptions that do not hold exactly, the error-control claim is carried not
by the formula but by the direct empirical calibration (\S3.4): a
Monte-Carlo of the full gate on a best-of-100 random-map null-generating
process (\texttt{scripts/\allowbreak gate\_null\_calibration.py}). The
interval on that calibration is the \textbf{one-sided exact
Clopper--Pearson} bound (the Beta-quantile inversion of the binomial), which
needs no normal approximation to the calibration count --- for 3 graduations
in 500 replicates, a 95\% upper bound of $\approx 1.54\%$; the verdict
partition of the 500 replicates is in Table~\ref{tab:calib} (no
INCOMPLETE, the search being instrumented).
\citet{bailey2014} supply the selective-inference lineage (\S2.4) and the
blended expression above; the operational error-control claim nonetheless
rests on the empirical calibration, not on the formula.

\textbf{Cognate-ID baseline (\S3.4).} The match-rate the deflation is
applied to is supplied by the published cognate-ID matcher identified in
\S2.3 (reproduced in the artifact), so the gate does not depend on any
single citation.

\begin{center}\rule{0.5\columnwidth}{0.4pt}\end{center}

\noindent\emph{Software \& data.}
\ifanon
The artifact. Code, pre-registrations, and findings are released as an
anonymized artifact on acceptance. Licensed corpora (GORILA-, SigLA-,
lineara.xyz-, D\=AMOS-derived) are not redistributed; a data-provenance note
documents each source's licence and retrieval path.
\else
logos (v0.1.0). Zenodo. Concept DOI 10.5281/zenodo.21087572; repository
\url{https://github.com/papadopouloskyriakos/logos}. Licensed corpora
(GORILA-, SigLA-, lineara.xyz-, D\=AMOS-derived) are not redistributed; see
\texttt{docs/data-provenance.md}.
\fi

\section{Complementary tables}

\begin{table}[h]
\centering\small
\resizebox{\columnwidth}{!}{%
\begin{tabular}{@{}lrrrl@{}}
\toprule
\textbf{Stratum} & \textbf{Inscr.} & \textbf{Words} & \textbf{Sites} &
\textbf{Role} \\
\midrule
admin & 290 & 1,486 & 14 & induction \\
libation (stone vessels) & 99 & 370 & 15 & induction \\
other & 129 & 187 & 40 & induction \\
seal & 740 & --- & --- & excluded \\
emptied by abbrev.-strip & 83 & --- & --- & excluded \\
\midrule
total & 1,341 & & & \\
\bottomrule
\end{tabular}}
\caption{Morphology corpus strata (Appendix~A.1). Seals are structurally
excluded as an abbreviation channel; the strata reconcile exactly to the
1,341-inscription corpus total.}
\label{tab:strata}
\end{table}

\begin{table}[h]
\centering\small
\resizebox{\columnwidth}{!}{%
\begin{tabular}{@{}lll@{}}
\toprule
\textbf{Bucketing} & $\alpha=0.01$ & $\alpha=0.05$ \\
\midrule
strict (stone vessels only) & NULL & NULL \\
loose (stone objects + inked) & NULL & \texttt{i-}, \texttt{-te/-ti}
validate \\
\bottomrule
\end{tabular}}
\caption{Morphology sensitivity $2\times2$ (Appendix~A.1): genre bucketing
$\times$ significance threshold. Only the doubly-relaxed corner validates
the H3 nominal affixes \texttt{i-} (locative) and \texttt{-te/-ti}
(ablative).}
\label{tab:sens}
\end{table}

\begin{table}[h]
\centering\small
\resizebox{\columnwidth}{!}{%
\begin{tabular}{@{}ll@{}}
\toprule
\textbf{Quantity} & \textbf{Value} \\
\midrule
Linear A confirm rate (16 affixes) & 0.250 (4/16; [0.10, 0.50]) \\
\texttt{L\_fake} bigram-floor rate & 0.375 (6/16; [0.18, 0.61]) \\
\texttt{i-*301} vs.\ within-form null & $z \approx 8.0$ (6 environments) \\
\midrule
Linear B (D\=AMOS) confirm rate & 0.562 (9/16) \\
Linear B within-word-shuffle floor & $\sim$0.03 \\
Linear B \texttt{L\_fake} bigram floor & $\sim$0.30 (5 draws) \\
\midrule
word-token mean length & 1.84 signs (median/mode 1) \\
word tokens $\le$2 signs / 1 sign & 76.1\% / 56.5\% \\
multi-sign \emph{type} mean length & 3.07 (Fuls 3.3, 554-type list) \\
Linear B mean length & 3.23 (13,562 wordforms) \\
\midrule
DP-unigram segmenter micro-F1 & 0.436 \\
random-boundary baseline & 0.389 (base rate 0.406) \\
real$-$random gap, 95\% CI & [0.021, 0.099]; P(gap${>}$0) = 0.998 \\
Morfessor & 0.156 (over-segments; not usable) \\
\bottomrule
\end{tabular}}
\caption{Morphology-induction replication quantities (Appendix~A.1):
confirm rates with Wilson 95\% intervals, word-length denominators, and
the segmentation positive with its site-clustered bootstrap gap (4,000
draws over 52 sites).}
\label{tab:morphrates}
\end{table}

\begin{table}[h]
\centering\small
\resizebox{\columnwidth}{!}{%
\begin{tabular}{@{}ll@{}}
\toprule
\textbf{Quantity} & \textbf{Value} \\
\midrule
documents / balance units & 32 / 35 \\
integer-level balance & 7/28 (25\%) \\
determinable fraction value & J = 1/2 (HT104: $2 \cdot J = 1$) \\
held-out fraction balance (real) & 0.0 \\
permutation null & mean 0.113 (max 0.143); $p = 1.0$ \\
planted positive control & 0.9 vs.\ null 0.22; $p = 0.003$ \\
\bottomrule
\end{tabular}}
\caption{Metrology (Direction D) replication quantities (Appendix~B.1).
The real corpus is not separated from the permutation null; the planted
control separates, so the null is not a machinery bug.}
\label{tab:metro}
\end{table}

\begin{table}[h]
\centering\small
\resizebox{\columnwidth}{!}{%
\begin{tabular}{@{}lrrrr@{}}
\toprule
\textbf{Test} & \textbf{Classes} & \textbf{Acc.} & \textbf{Null} & $p$ \\
\midrule
consonant class & 13 & 0.080 & 0.065 & 0.40 \\
vowel class & 5 & 0.140 & 0.201 & 0.86 \\
\bottomrule
\end{tabular}}
\caption{Distributional-phonology per-test results (Appendix~B.2):
leave-one-out 1-NN vs.\ a label-permutation null, \v{S}id\'{a}k-deflated
over the two tests. The positive-control power curve fires ($p<.05$) at
planted strength $\ge 2$ but not at 1.}
\label{tab:phon}
\end{table}

\begin{table}[h]
\centering\small
\resizebox{\columnwidth}{!}{%
\begin{tabular}{@{}lll@{}}
\toprule
\textbf{Leg} & \textbf{Accuracy} & \textbf{Floor} \\
\midrule
classical, direct-NN & 0.410 [0.18, 0.64] & 0.0135 \\
classical, Procrustes & 0.185 [0.00, 0.36] & 0.0135 \\
I-JEPA (3 seeds; excluded) & $0.324 \pm 0.025$ & 0.0135 \\
font-swap, direct-NN & $0.367 \rightarrow 0.416$ & $\gg$ chance \\
font-swap, Procrustes & $0.156 \rightarrow 0.195$ & $\gg$ chance \\
sequence/cognate (D\=AMOS) & 0.025 & 0.011 \\
\quad positive control & 0.947 & --- \\
\bottomrule
\end{tabular}}
\caption{Cross-script sign-image alignment and controls (Appendix~B.3).
Font-swap rows give same-font $\rightarrow$ independent-fonts accuracy on
the control's matched anchor set; the sequence/cognate row is the best of
five alignment methods over 56 anchors
(\texttt{clearly\_above\_chance} = False).}
\label{tab:xscript}
\end{table}

\begin{table}[h]
\centering\small
\resizebox{\columnwidth}{!}{%
\begin{tabular}{@{}llrr@{}}
\toprule
\textbf{Item} & \textbf{Source} & \textbf{LLM} & \textbf{Baseline} \\
\midrule
\texttt{JA-NE} = \emph{yn} ``wine'' & Gordon & 14/40 & 0/40 \\
\texttt{A-SA-SA-RA-ME} & \citet{best1981} & 2/40 & 0/40 \\
\texttt{KA-RO-PA} & Gordon & 0/40 & 0/40 \\
\texttt{SU-PA-RA} & Gordon & 0/40 & 0/40 \\
\texttt{KU-RO} = \emph{kull} (downwt.) & Gordon & 8/40 & --- \\
\texttt{KU-RO} = ``total'' (excl.) & gen.\ knowledge & 7/40 & --- \\
\bottomrule
\end{tabular}}
\caption{LLM-ablation per-item reproduction over 40 seeds,
\texttt{qwen2.5:72b} vs.\ the model-free edit-distance baseline
(Appendix~C.1). The two \texttt{KU-RO} routes overlap on one seed; three
single-word items are among Gordon's five West-Semitic equations, while
\texttt{A-SA-SA-RA-ME} is Best's.}
\label{tab:repro}
\end{table}

\begin{table}[h]
\centering\small
\resizebox{\columnwidth}{!}{%
\begin{tabular}{@{}lrr@{}}
\toprule
\textbf{Sign set} & \textbf{Signs} & \textbf{Valued (of 40)} \\
\midrule
known AB (indexed) & 20 & $\sim$20 \\
not-indexed *-series & 11 & $\sim$1 ($\sim$97\% abstention) \\
\bottomrule
\end{tabular}}
\caption{\texttt{L\_not\_indexed} abstention asymmetry (Appendix~C.2).
The apparent consistency gap 0.601 (permutation $p = 0.0005$) is
coverage-confounded; the matched-$n$ check returned \texttt{no\_power}
(0 of 11 not-indexed signs reached $n \ge 10$).}
\label{tab:abstain}
\end{table}

\begin{table}[h]
\centering\small
\resizebox{\columnwidth}{!}{%
\begin{tabular}{@{}ll@{}}
\toprule
\textbf{Denominator / channel} & \textbf{Value} \\
\midrule
token-frequent syllabograms & 97 \\
full syllabary estimate & $\sim$150 \\
working repertoire ($V$; conflates syll./subscr./logograms) & 259 \\
incl.\ ideograms (GORILA: 390) & $\sim$350 \\
\midrule
sign occurrences, all types & $\sim$7,400 \\
parsed syllable-signs (unicity input) & $N \approx 5{,}792$ \\
\midrule
corpus envelope & $\sim$2,500 finds; $\sim$1,400 readable \\
findplaces (logos-normalized) & 56 ($\rightarrow$ 52 sites) \\
readable inscriptions & 1,341 \\
distinct word-forms & 539 \\
\midrule
unicity distance over $(V,P)$ space & $U \approx 204$--$415$ signs \\
\bottomrule
\end{tabular}}
\caption{Corpus denominators, occurrence channels, and the unicity range
(Appendix~C.4). The $\sim$7,400 all-type occurrence channel never enters
the unicity arithmetic, which consumes only the parsed syllable-signs.}
\label{tab:denoms}
\end{table}

\begin{table}[h]
\centering\small
\resizebox{\columnwidth}{!}{%
\begin{tabular}{@{}ll@{}}
\toprule
\textbf{Calibration quantity} & \textbf{Value} \\
\midrule
\texttt{S\_morph} FPR (within-form shuffle) & 3.0\% (6/200; [1.4, 6.4]\%) \\
nominal one-sided $z \ge 2$ rate & $\approx$2.3\% \\
\texttt{S\_morph} power (planted affix) & 100\% (50/50; saturated) \\
\midrule
gate false-graduation (best-of-100 null) & 0.6\% (3/500; CP $\le$ 1.54\%) \\
gate verdict partition (500 replicates) & 3 GRAD / 10 NULL / 487 REJ \\
\bottomrule
\end{tabular}}
\caption{Detector and gate calibration quantities (Appendices~A.2
and~D): Wilson interval for the \texttt{S\_morph} false-positive rate;
one-sided exact Clopper--Pearson 95\% upper bound for the gate's
false-graduation rate under the tested best-of-100 random-map null.}
\label{tab:calib}
\end{table}

\begin{table}[h]
\centering\small
\resizebox{\columnwidth}{!}{%
\begin{tabular}{@{}ll@{}}
\toprule
\textbf{Quantity} & \textbf{Value} \\
\midrule
literature-index census & 63 claims / 62 entries (55 single) \\
root-template TV distance & $\sim$0.84 $\rightarrow$ $\sim$0.33
  ($\sim$0.23 standalone) \\
\midrule
ablation scale & 40 seeds $\times$ 40 held-out forms \\
model-free lexicon & 7,682 Hebrew skeletons \\
LLM errors & 0 (40/40 seeds) \\
value-vocabulary overlap & 10 of 134 \\
first-pass contamination rate & 1.000 (33/33 defined; artifact) \\
gate-2 rate & mean 0.646 (confounded; not reported) \\
\bottomrule
\end{tabular}}
\caption{Literature-index, \texttt{L\_fake}-canary, and LLM-ablation
quantities (Appendices~C.1 and~C.5). The TV distance is measured
end-to-end on the Ugaritic$\leftrightarrow$Hebrew gold; the first-pass
contamination rate is forced to 1.0 by construction and reported as an
artifact.}
\label{tab:canary}
\end{table}

\end{document}
.

% ---- identity macros (switch on \ifanon) ----
\ifanon
  \newcommand{\artref}{the artifact}
  \newcommand{\artobj}{the artifact}
\else
  \newcommand{\artref}{logos v0.1.0}
  \newcommand{\artobj}{the logos artifact (v0.1.0)}
\fi

% ---- PDF metadata: named for arXiv, stripped for the anonymous TACL PDF ----
\ifanon
  \hypersetup{pdfauthor={},pdftitle={},pdfsubject={},pdfkeywords={},pdfcreator={},pdfproducer={}}
\else
  \hypersetup{pdfauthor={Kyriakos Papadopoulos},
    pdftitle={Falsification-First Decipherment: A Decontaminated Inference
Framework for Testing Undeciphered-Script Claims, with Linear A as a
Worked Null},
    pdfcreator={},pdfproducer={}}
\fi

\title{Falsification-First Decipherment: A Decontaminated Inference
Framework for Testing Undeciphered-Script Claims, with Linear A as a
Worked Null}

\ifanon
  \author{}
\else
  \author{
    Kyriakos Papadopoulos \\
    Independent Researcher \\
    Netherlands \\
    ORCID 0009-0003-0995-2518
  }
\fi

\date{}

\begin{document}
\maketitle
\ifanon\else
\begin{center}\small\emph{Preprint --- under review; not peer
reviewed.}\end{center}
\fi

\begin{abstract}
Decipherment research has a false-positive problem and no shared
instrument: a small corpus, an unknown language, and many candidate
readings let an analyst ``translate'' almost anything. The flagship
survey lists fifteen challenges of script and data; none addresses the
orthogonal inference-control layer --- overfitting, multiple testing,
contamination --- a candidate sixteenth challenge. We build a
falsification-first framework for it: a fabricated-language
negative-control family, multiple-comparison deflation, a literature
index separating published from not-indexed signs, and pre-registration
with a mechanically-graded gate. This is a methods paper; we assert no
decipherment. The machinery calibrates: on a tested null the gate's
false-graduation rate is 0.6\% (3/500); the identical morphology test
finding no power on Linear A recovers Mycenaean (Linear B) inflection
(9/16 affixes above the bigram floor), supporting the interpretation
that Linear A's short forms contribute materially to its lack of power.
Across probes it returns differentiated, calibrated outcomes --- null,
no-power, supported-structural, circular, and contamination-sensitive.
Results are consistent with identifiability, rather than corpus size,
being the more fundamental present constraint. The June-2026 *301
``cracked'' claim is a qualitative illustration of the failure modes the
gate is designed to detect.
\end{abstract}

\section{Introduction}

Decipherment has a false-positive problem, and the field has no agreed
instrument for measuring it. The history of undeciphered-script research
is, in large part, a history of uncontrolled model selection: a tiny
corpus, an unknown language, and a large space of candidate readings let
a determined analyst ``translate'' almost anything, with internal
consistency on the cherry-picked set standing in for confirmation. The
mechanism is not any single bad match but uncorrected multiplicity: with
enough candidate roots and a small enough corpus, a pile of
individually-plausible matches is the \emph{expected} outcome under pure
chance, so each individual comparison looks suggestive while the
conclusion is nonsense. A standard demonstration of spurious
comparison: one can ``prove'' \emph{English is a Semitic language} by
matching a handful of English words to Proto-Semitic roots while quietly
ignoring the far larger number that do not match. The same move, applied
to a real undeciphered script, has produced a long line of decipherment
claims that no one can either confirm or refute.

That the discipline lacks a shared standard here is not our editorializing
but the field's own assessment. The flagship computational-decipherment
survey --- \citeauthor{braovic2024}'s \citeyearpar{braovic2024}
systematic review of Bronze Age Aegean and Cypriot scripts ---
enumerates fifteen ``major challenges'' of script, language, and data
(\S2.1) \citep{braovic2024}. Those fifteen are
primarily properties of the \emph{object and its data} --- the script,
the language, the surviving corpus, and the resource/comparison-data
conditions. What
governs whether a \emph{claimed reading} is believable is a different
kind of thing: a property of the inference methodology and its
epistemics --- overfitting, multiple testing, false-positive control,
and contamination. We therefore frame this axis, once and modestly, as a
candidate \textbf{sixteenth challenge} --- a deliberate rhetorical foil
to the survey's fifteen, not a claim to a new taxonomic slot or to sole
discovery. It is \emph{under-instrumented} rather than unrecognised
(\S2.4): our contribution is to
\emph{consolidate and operationalize} that axis for decipherment, not to
invent it. It is the methodological axis the object-level taxonomy was
never meant to cover, and it is distinct from the survey's own
acknowledged evaluation gap (\S2.1) --- an \emph{evaluation-target}
problem, separable from the false-positive /
multiplicity problem; it is the latter this paper instruments.

\textbf{Contribution.} We build the machinery for that axis: a
falsification-first, decontaminated inference framework for testing
undeciphered-script decipherment claims \citep{braovic2024}. The harness
applies fail-closed validation to every positive and enforces four
concrete defences: a fabricated-language canary that supplies a
\textbf{generator-conditional negative-control distribution} (a
false-positive floor conditional on the fabricated-language generator ---
here a first-order sign-bigram Markov model; a broader control ensemble
spanning unigram, root-template, and real-language generators is
disclosed as future work, not built here); multiple-comparison deflation
for the number of signs, roots, and families a search was free to try; a
literature index that partitions published (\texttt{L\_indexed}) from
literature-unseen (\texttt{L\_not\_indexed}) signs so that reproducing a
known value cannot masquerade as discovery; and pre-registration with a
held-out gate whose verdict is computed mechanically, never by the model
that proposed the reading. The evaluation gap gets a concrete, if
power-limited, answer: the recurring libation \emph{formula} --- roughly
50 inscribed objects spread across peak sanctuaries and caves at more
than five findspots --- supports a natural leave-one-site-out
cross-validation graded from held-out data --- explicitly a
low-statistical-power test, reported as such rather than as decisive
(\S2.1; Appendix~A.3) \citep{braovic2024}.

\textbf{This is a methods paper, and Linear A is its worked null.} We do
not claim a decipherment. The framework is designed to withhold
validation when a proposed reading does not clear pre-registered negative
controls, multiplicity adjustment, contamination checks, and held-out
evaluation; when nothing survives, the calibrated null --- what the
corpus can and cannot support --- is itself the deliverable. The working
conclusion on current data is that the results are \textbf{consistent
with identifiability} being the more fundamental present constraint ---
not a demonstrated corpus-size threshold: there is no known related
language to map onto, the multiple-testing burden of candidate readings
goes uncorrected, and the sign inventory is mixed and partly lossy. We do
not claim additional data would be irrelevant --- a null can arise from
identifiability, low power, or corpus size alike. A lossy, non-injective
sign-to-sound channel need not admit a decipherment at any corpus size,
so we make no claim that ``there is enough text to pin a reading''; the
rigorous demonstration of \emph{that} constraint is the result the field
can use. The June-2026 AI-assisted ``Linear A cracked'' claim --- reading
the *301-anchored libation verb as an extinct Semitic form --- serves
throughout as a qualitative illustration of the failure modes the gate is
designed to detect, analysed from its public, archived presentation as a
claim, not a person, and neither reproducible nor held-out validated.
The artifact is open and archived (\S9.7).

\textbf{Roadmap.} Section~2 situates logos against the Braovi\'{c} et al.
taxonomy and the computational-decipherment lineage, including the
cognate-ID baseline. Section~3 specifies the harness --- canary,
deflation, literature index, and the pre-registration/held-out gate.
Section~4 develops the information-floor diagnostic and states its
identifiability limits. Section~5 works the June-2026 *301 claim as a
case study of the failure modes. Sections~6--8 report the pre-registered
nulls: morphology (\S6); metrology, phonology, and cross-script signals
(\S7); and contamination and sufficiency probes (\S8). Section~9 draws
the discussion, limitations, and conclusion.

\section{Related work}

\subsection{Position relative to the survey taxonomy}

The reference frame for computational work on the Bronze Age Aegean
scripts is \citet[\emph{Computational Linguistics}
50(2):725--779]{braovic2024}, the field's first survey focused
specifically on Aegean and Cypriot scripts; its only comparable
predecessor is \citet{ferrara2022}, and the broader
machine-learning-for-ancient-languages landscape is surveyed by
\citet[\emph{Computational Linguistics} 49(3)]{sommerschield2023}.
Braovi\'{c} et al.'s Section~4 enumerates fifteen ``major challenges'' ---
from unknown writing system, reading direction and language, through
small dataset, allographs and incomplete vocabulary, to unavailable
parallel data and hardware --- \textbf{none} of which concerns
overfitting, multiple testing, false-positive control, contamination, or
distinguishing a \emph{fitted} decipherment from a \emph{real} one. The
same gap runs through Ferrara and Tamburini's
\citeyearpar[\emph{Lingue e Linguaggio} XXI.2:239--259]{ferrara2022}
\S3.5, which reviews the cognate-matching lineage logos builds on ---
\citet{snyder2010} on Ugaritic$\leftrightarrow$Hebrew (the 29/30-sign
benchmark our canary calibrates against), \citet{bergkirkpatrick2011},
and \citet{luo2019, luo2021} --- and frames Linear A, via Gelb and
Whiting's typology \citep{gelb1975}, as a Type III script (unknown script
\emph{and} language) whose ``necessary validation'' is ``altogether
impossible,'' yet nowhere treats how a match is to be \emph{deflated} for
the multiplicity it was selected from. Two surveys, the same unfilled
axis. We do not score this as an oversight: their challenges are
properties of the object and its data, whereas the axis logos foregrounds
is a property of inference \emph{methodology} --- named-but-unfilled
rather than absent. logos supplies the machinery to police it (\S3).

Braovi\'{c} et al. do gesture at a related but separable difficulty:
evaluation. Because computational analyses of an undeciphered script have
``nothing to compare them against,'' they conclude (Sec.~6.3) that
``their evaluation remains an open problem.'' That is a distinct problem
from overfitting and multiplicity control, and logos addresses it
separately. Its held-out design is a concrete response --- the libation
formula recurs across a small set of peak-sanctuary and cave findspots,
permitting a leave-one-site-out cross-validation adjudicated mechanically
by \texttt{verdict.py} --- with the low-power caveat flagged at the
outset (\S1). We locate our components in the survey's own
method taxonomy (Sec.~7.1) and, where prior art exists, defer to it
rather than claim novelty.

\subsection{The current philological synthesis}

\citet{salgarella2025}, \emph{Writing in Bronze Age Crete: ``Minoan''
Linear A} (Cambridge Elements), by the author of the SigLA database, is
the most current authoritative one-author overview. It concludes that
Minoan is a language \textbf{isolate} --- it ``does not belong in any
known language family'' (the families used for comparison being
Indo-European, Semitic and Afro-Asiatic) --- typologically agglutinative
with a provisional Verb--Subject--Object order, and it records that ``no
bilingual text has so far been found'' in a corpus of roughly 2,500
inscriptions. Crucially, Salgarella ranks the etymological/comparative
method below combinatorial-contextual analysis and names
\textbf{digital/statistical analysis as ``the next desideratum in the
field.''} logos positions itself as exactly that desideratum --- but
under discipline, treating the isolate verdict as a central
\emph{identifiability} constraint rather than a target to be argued away.

\subsection{Prior computational attempts and the search trap}

The seminal automatic-decipherment result is \citet{luo2019}, which
recovers a character mapping from a lost script to a \emph{known related}
language via neural minimum-cost flow. It auto-deciphered
Ugaritic$\rightarrow$Hebrew ($+5.5\%$ over prior state of the art) and
translated 67.3\% of Linear B$\rightarrow$Greek cognates, but did not
address Linear A and would fail there by construction: with no known
cognate plaintext, the objective has nothing to optimise toward. This
cognate-language precondition is the empirical form of our
identifiability argument. The cognate-ID baseline logos actually runs is
\emph{not} a reimplementation of Luo's neural flow: it is a published
matcher, \citeauthor{tamburini2025}'s \citeyearpar{tamburini2025}
CSA\_OptMatcher (coupled simulated annealing), which reports the same
\citet{luo2019} accuracy metric. We adopt it as the non-neural baseline
against which any Linear A cognate claim must clear.

Against that disciplined baseline stand the in-domain instances of the
search trap: the dictionary-fishing programmes of \citet{eu2019} and
\citet{loh2020}, the latter explicitly a ``brute force attack'' on
dictionaries. These are the concrete Linear-A cases the multiplicity
correction is built to deflate. On the visual side, cross-script
comparison is \emph{not} our novelty: \citet{daggumati2019, daggumati2023}
already applied CNN+SVM cross-script similarity, and \citet{corazza2022}
(``Sign2Vec'') clustered signs unsupervised; our defensible contribution
is only the controlled image-versus-sequence \emph{asymmetry} under one
harness.

\subsection{The discipline lineage}

Two older strands supply the honesty machinery. \citet{packard1974},
\emph{Minoan Linear A}, constructed nine deliberately fictitious
decipherments as a null control --- and the genuine Ventris-value transfer
beat them only about 2:1 --- anticipating our permutation-null canary. The
statistical lineage is selective inference and multiple-comparison
correction \citep{bailey2014}, the correction for the number of
hypotheses $\times$ signs $\times$ roots $\times$ families tried that
turns a raw match-rate into a defensible one; its transfer from the
finance/backtest setting to script decipherment is specified in
Section~3. logos's empty cell is the intersection none of these occupy: a
falsification-first, decontaminated inference framework that couples a
published cognate-ID matcher \citep[reporting the Luo 2019
metric]{tamburini2025} with preregistration, multiplicity deflation and
mechanical held-out verdicts, released openly (\S9.7), with Linear A
worked as the honest null rather than a claimed crack.

\section{Methods: the discipline harness}

The harness is a comparison-layer pipeline that a reader can reimplement.
A candidate reading enters as a \emph{proposal} and passes through four
ordered stages --- pre-registration, deflation, decontamination, held-out
graduation --- under a cross-cutting persistence invariant: every decision
is computed from stored artifacts, never from a model's narration. Model-derived signals --- LLM theses, learned embeddings,
cognate-matcher assignments --- enter capped at confidence $\le 0.75$,
structurally unable to \emph{decide} a verdict; only
mechanically-verified held-out reads rank higher (\S1--2).

\subsection{Pre-registration}

Before any model is built or run, the predictions are frozen. The
morphology pre-registration fixes, for each hypothesis, the claimed affix,
its source, the falsifiable test, and the acceptance threshold, frozen
in-repository before each run so the model cannot be tuned to them after
the fact (the anchoring mechanics and the per-affix acceptance criteria
are given in Appendix~A.1). Outcomes are enumerated in advance --- CONFIRM, EXTEND,
REFUTE/NULL, or NO POWER --- so that a negative is a publishable result,
not a failed experiment. Any change is a new dated pre-registration,
never an edit.

\subsection{Deflation}

Internal fit on the derivation set is treated as no evidence. Every match
rate is deflated for everything tried: \texttt{n\_hypotheses} $\times$
\texttt{n\_signs} $\times$ \texttt{n\_roots}. The nulls are pipeline-level
--- a frequency-banded permutation null (after \citealp{packard1974}) and
the order-statistic bar of \S3.4 (implementation detail in Appendix~D).
Multiplicity is instrumented, not hand-passed, by a deduplicating
\texttt{SearchLog} (invariant 12; Appendix~D). The strong structural statistic,
\texttt{S\_morph}, earns credit only through a productivity gate --- an
affix must attach to $\ge$2 distinct residual stems, so a single formula
word copied across inscriptions scores zero (calibration in
Appendix~A.2, Table~\ref{tab:calib}).

\subsection{Decontamination}

A large open-weight model may have encountered longstanding published
readings by Gordon and Best; reproduction of indexed historical proposals
is therefore treated as potential shared-source contamination rather than
independent evidence. The June-2026 Di Mino claim is analysed separately
as a contemporary public case and is not assumed to be present in the
model's pretraining data --- it is indexed so that future claims and
models cannot reproduce it as discovery. Three mechanisms address this.

\textbf{The literature index} records published Linear A sign-value and
correspondence claims; any proposal matching the index is tagged
\texttt{literature\_match} and quarantined --- it may be tested but can
never count as discovery. The index partitions the sign inventory into
\texttt{L\_indexed} (referenced by $\ge$1 published proposal) and
\texttt{L\_not\_indexed} (untouched). Discovery claims may rest only on
\texttt{L\_not\_indexed} held-out success, since regurgitation can only
return what is in the literature. (The conservative seed, the adversarial
population, and the census are given in Appendix~C.5,
Table~\ref{tab:canary}.)

\textbf{The \texttt{L\_fake} canary} is a phonotactically-plausible
invented lexicon, calibrated to a real candidate language
(Appendix~C.5), generated after the model's release and never published
by us; exact-string collisions with real lexica are measured
(Appendix~C.5). Running the full pipeline
against it yields a \textbf{generator-conditional} false-positive floor:
any apparent lexical support it produces is spurious with respect to the
target historical-language hypothesis, and memorisation cannot
be invoked because there is nothing to memorise. A real candidate's match
must outperform this \emph{one calibrated class of negative controls} by
the corrected margin --- support means the candidate beat this
generator-conditional control family (here a first-order sign-bigram /
root-template generator), not that it defeated every possible false-language
process.

\textbf{LLM-ablation} compares the pipeline with and without the model's
proposals to isolate contamination (procedure in Appendix~C.1).

\subsection{The held-out graduation gate}

A reading graduates only mechanically, from held-out data, as a
conjunction of pre-registered clauses computed by \texttt{scripts/verdict.py}
--- never by the model that proposed the reading. The operative deflation
clause is an \textbf{order-statistic bar}: the observed held-out statistic
must exceed the \emph{expected maximum} of the multiplicity-null taken over
the number of trials the search was free to run (the E[max] over
\texttt{n\_trials} deflation), a \textbf{normal-order-statistic
approximation} in which \textbf{$n =$ \texttt{n\_trials} is the \emph{count of
candidates the search tried}, not an estimated number of independent
hypotheses}. The explicit formula and its Gaussian/independence assumptions
are given in Appendix~D; the error-control claim is carried not by the formula
but by the direct empirical calibration below. This is a direct
correction for selection across the hypothesis
space, and it fails closed on a degenerate null (Appendix~D). The framework yields
four primary evidential outcomes: \textbf{GRADUATE} (every clause holds and the
search multiplicity was instrumented), \textbf{REJECT} (a clause failed, or
the statistic sits below the corrected bar), \textbf{NULL\_PUBLISHED} (the
statistic is indistinguishable from the \texttt{L\_fake} null --- a
calibrated non-result, publishable as such), and \textbf{NO POWER} ---
the last assigned \emph{upstream}, by the pre-registered module outcome
when the experiment cannot adjudicate (\S3.1), before the gate sequence
runs; the gate itself additionally emits a fail-closed \textbf{INCOMPLETE}
status when search multiplicity is un-instrumented. These outcomes concern \emph{evidential
support, not truth}: because Linear A has no ground truth, REJECT / NULL /
NO POWER mean a reading is unsupported or insufficiently validated on
held-out data, \textbf{never that the underlying linguistic proposition is
false}. The verdict follows a fixed decision sequence
(\texttt{verdict.py}): all pre-registered clauses hold on an above-null
match $\rightarrow$ \textbf{GRADUATE}; else a match clearing the
substantive clauses but on an un-instrumented search $\rightarrow$
\textbf{INCOMPLETE}; else an observed statistic not exceeding
the null mean $\rightarrow$ \textbf{NULL\_PUBLISHED}; else $\rightarrow$
\textbf{REJECT}.
We calibrate the refusal rule directly: under the
tested best-of-100 random-map null --- an analyst trying 100 random
candidate maps and submitting the best --- 3 of 500 replicates graduated,
an observed false-graduation rate of 0.6\% (one-sided exact
Clopper--Pearson 95\% upper bound $\approx 1.54\%$), with 487/500 rejected
outright and 10 published as calibrated nulls (Appendix~D); this is an
empirical
error-control result conditional on that null-generating process, not a
universal guarantee. A second
pre-registered clause, \textbf{\texttt{generalizes\_to\_not\_indexed}},
requires the reading to recover literature-unseen (\texttt{L\_not\_indexed})
signs --- so reproducing already-published values cannot by itself
graduate a claim; the clause fails closed (its pre-committed thresholds
are in Appendix~D). Graduation is defined as validated \emph{novel}
discovery; independent held-out support for an already-published reading
is recorded as a distinct outcome and does not graduate. The remaining
pre-registered clauses (\texttt{L\_fake} separation, corrected margin,
order-statistic bar) are conjoined as specified in Appendix~D. These are
\emph{gate parameters} --- pre-registered
choices for this corpus, not universal constants. Validation is
leave-one-site-out, not ordinary \emph{k}-fold (Appendix~A.1). The gate is
AND-monotone --- clauses can only make it stricter --- so hardening never
lets a previously-failing reading graduate.

\subsection{Grade from persisted artifacts}

The invariant across every gate: verdicts are computed from persisted
rows, never from a model's account of its own outcome. The verdict writer
is grep-clean of any model call and is the sole writer of the verdicts
table.

\begin{table*}[t]
\centering
\begin{tabular}{@{}p{2.6cm}p{3.0cm}p{2.4cm}p{3.0cm}p{3.3cm}@{}}
\toprule
\textbf{Probe} & \textbf{Control / comparator} & \textbf{Real outcome} &
\textbf{Null / control floor} & \textbf{Verdict} \\
\midrule
Graduation gate (\S3.4) & --- & --- & best-of-100 random-map null &
3/500 false graduations (0.6\%; calibrated under the tested null) \\
Morphology (\S6) & Linear B 9/16 & Linear A 4/16 & \texttt{L\_fake}
bigram 6/16 & NO POWER \\
Segmentation (\S6) & --- & micro-F1 0.436 & random 0.389 & supported (gap
CI excludes 0) \\
Metrology (\S7.1) & planted 0.9 & 0.000 & permutation 0.113 & null
(agrees with the field) \\
Phonology (\S7.2) & planted (fires $\ge$ 2) & 0.080 / 0.140 & 0.065 /
0.201 & null (data-limited) \\
Image alignment (\S7.3) & font-swap survives & 0.410 & image chance
0.0135 & circular demonstration (capped) \\
Contamination (\S8.1) & model-free 0/40 & 14/40 (JA-NE) & 0/40 &
exploratory (public-exposure gradient) \\
\bottomrule
\end{tabular}
\caption{One harness, calibrated epistemic outcomes across probes.}
\label{tab:harness}
\end{table*}

The same falsification-first harness returns \emph{different} calibrated
verdicts across probes (Table~\ref{tab:harness}) --- the paper's
contribution in one view. Each number is detailed in the cited section.

\section{The information floor}

Before any decipherment method is graded, the corpus must be sized against
what it can identify. Two quantities matter: how much text exists, and how
many free parameters a reading spends against it. Confusing the first for
the second is the origin of the ``too little text'' folk claim --- and, we
argue, of its equally unearned inverse, the optimism that a big-enough
corpus makes a reading identifiable.

\textbf{The corpus.} The current authoritative synthesis
\citep{salgarella2025} gives an envelope of roughly 2,500 total Linear A
finds, of which about 1,400 are readable inscriptions, overwhelmingly
concentrated in the single LM IB destruction horizon (findplace and
dating detail in Appendix~C.4). That near-single horizon is why a
site-held-out split, not a chronological one, is the only realistic
partition. No bilingual text has ever been found
\citep[\S8]{salgarella2025}.

\textbf{Corpus sizing and the unicity computation (numeric detail in
Appendix~C.4).} The recurring sign-count confusion resolves into several
distinct sign-count denominators and two occurrence channels that must be
kept distinct (Table~\ref{tab:denoms}; Appendix~C.4). On the parsed
syllable-sign budget, Shannon's unicity distance sits everywhere far
below the corpus size (Table~\ref{tab:denoms};
\citealp[\S1--3]{salgarella2025}); we retain
this only as a toy-model \emph{precondition} diagnostic and explicitly
withdraw it as any sufficiency or refutation claim. Qualitatively, the
figure cannot settle identifiability either way (Appendix~C.4): unicity
distance presupposes a known plaintext-language oracle we do not have.
Corpus size, in short, cannot decide the question.

\textbf{What actually binds.} The results are consistent with
\emph{identifiability} being the more fundamental present constraint rather
than corpus size \citep[\S8]{salgarella2025}; we do not establish that more
data would be irrelevant. It has three parts: (i) the \textbf{absence of a
known cognate language} --- the map has a value-space but no anchored
target (Luo 2019's limit); (ii) the \textbf{multiple-testing / search
trap} --- that a unique consistent map may exist in principle neither
guarantees an algorithm finds it nor stops a cherry-picked false positive
across the parameters tried (invariant 8); and (iii) a \textbf{mixed
inventory} of syllabograms, logograms, and numerals that a real reading
must first separate. These, not the count of surviving signs, are what a
decipherment must overcome. Accordingly, when a candidate lexicon grows
without bound on a fixed occurrence budget --- free parameters exceeding
the floor --- the reading is underdetermined, and the framework must say
so.

\section{A qualitative audit of the *301 claim}

The framework of Section~3 is abstract until it meets a concrete claim. In
June 2026 a widely circulated ``Linear A cracked'' announcement
\citep{dimino2026} argued that Linear A records an extinct Semitic
language, reviving Gordon's hypothesis \citep{gordon1957, gordon1966}. Its
load-bearing anchor is a single reading: the recurring peak-sanctuary
libation formula word A-TA-I-*301-WA-JA, in which the Linear-A-only sign
*301 is \emph{itself} assigned a /na/-initial value, so the verb root
reads as West-Semitic N-W-Y ``to dwell'' (\emph{nawaya}) and is thereafter
matched to later Hebrew liturgy. We analyse this \textbf{as a claim, on its
public evidence} --- not the person --- from an immutable archived copy
carried in \artobj. It is a clean, self-contained instance of the failure
modes the gate is built to catch, and the three-part refutation below rests
wherever possible on the claim's \textbf{own published table} and on
primary sources \citep{davis2013, duhoux1992, thomas2020, salgarella2025,
steele2017} --- none advanced in logos's own voice, and the paper
introduces no new sign reading. We flag the meta-point once: this is a
literature-and-primary-source argument, not itself a logos held-out result
--- the graduation gate and \texttt{L\_fake} canary are validated as
instruments (unit tests; the known-answer Ugaritic$\leftrightarrow$Hebrew
surrogate) but have \textbf{not been run on this claim}. Scoring Gordon's
equations and the Di Mino lexicon against the \texttt{L\_fake} floor and
the order-statistic bar is directly runnable with the components of \S3 and
is the priority of the planned revision; that the refutation has not been
independently vetted by Aegean epigraphers is a disclosed limitation
(\S9.6), not a standard the work claims to meet.

\textbf{(a) The segmentation is not novel.} Root i-*301 plus affixation is
the reading of \citet{davis2013}; an independent segmentation by
\citet{thomas2020} reaches the same core and is tested head-to-head against
it (the current synthesis, \citealp{salgarella2025}, does not itself adopt
Thomas). On the Davis and Thomas segmentations, the affix-stacking is
evidenced by multi-attestation forms --- ta-na-i-*301-u-ti-nu (IO Za 6) and
a-ta-i-*301-de-ka (ZA Zb 3) --- that vary the material around a stable
i-*301 core. (These sigla are confirmed against Davis 2013: \texttt{TL Za
1}, \texttt{IO Za 6}, and \texttt{ZA Zb 3}.) The claim re-segments a
structure already established in the literature. Under the decontamination
layer this matters procedurally, not just historically: a correspondence
that reproduces a published segmentation is treated as potential
shared-source contamination,
tagged \texttt{literature\_match} and barred from counting as discovery
(\artref, \S C.1).

\textbf{(b) The gloss is contested and weaker.} The mainstream value of the
verb is ``give/dedicate'' --- Duhoux's \citeyearpar{duhoux1992} gloss,
adopted by \citet{davis2013} as a \emph{structural} reading that assigns
*301 \textbf{no phonetic value at all}, and endorsed by the current
synthesis \citep[\S7.2]{salgarella2025}. ``Dwell'' is a semantically
distant, unmotivated substitution. The point is not that ``give'' is
certainly right, but that a \emph{structural} reading needing no
sound-value for *301 already accounts for the data the phonetic claim is
invoked to explain.

\textbf{(c) The primary descriptions of this form's morphology do not
support the Semitic label.} The claim's own table segments the verb as
prefix-stacking and agglutinative --- [prefix A ``I''] [stem-morpheme TA]
[stem-vowel I] [root *301]. We deliberately do \textbf{not} rest on an
\emph{a-priori} typological contradiction here, and we retract the earlier
draft's claim that a prefix-stacked verb is ``internally incoherent'' with
a Semitic reading: Semitic verbal morphology includes productive prefix
conjugations, so a prefix on the verb is not by itself un-Semitic. The
point is empirical, not typological, and it comes from the primary
literature. \citet[p.~38 n.~13]{davis2013} characterises i-*301's
morphology as ``neither IE nor Afroasiatic'' --- and Semitic \emph{is}
Afroasiatic, so a decade earlier the primary description already found
\textbf{no Afroasiatic morphological signature} in this very form. (We rest
on what the footnote states --- the absence of an Afroasiatic signature ---
not on a stronger ``ruled out'' than a hedged footnote can bear.)
Independently, \citet[\S8]{salgarella2025} concludes Minoan is
agglutinative (after Duhoux 1978) and a language isolate. The claim's own
segmentation is thus consistent with the isolate typology, while the
specific Semitic assignment finds no morphological support in the primary
descriptions of the form.

\textbf{The free-parameter receipt.} The decisive point for identifiability
is that *301 has no cross-script anchor: in the claim's own table its
AB-number column is ``---'' and its Linear B correspondence ``(none)''.
Even the standard practice of reading Linear A homomorphs with their Linear
B values is a backward projection to be handled with care
\citep[already cautioned by Duhoux 1989, apud Davis 2013, p.~38
n.~11]{steele2017} --- and an A-only sign such as *301 has no homomorph to
project from at all, so nothing cross-script constrains its value. The
single sign on which the entire decipherment pivots is thus a \textbf{free
parameter} --- its value was assigned to fit the desired root, not
constrained by any independent evidence. This is precisely the case the
minimum-description-length / identifiability \emph{diagnostic} flags: a
value chosen to make one word read out carries no description-length saving
that survives the multiplicity of alternatives it was selected from (cf.
Section~3; the diagnostic is reported next to every claim, and the
operative graduation clause --- the order-statistic bar over exactly that
multiplicity --- is one a hand-fitted value cannot clear). For this reading
the obstacle is identifiability rather than corpus size --- no known cognate
language to constrain the assignment, uncorrected multiplicity, and a sign
whose value is chosen rather than anchored (a free parameter is unanchored
at any corpus size).

Read against Section~3, the claim instantiates two distinct failure modes
at once. It is \textbf{conclusion-driven model specification} --- the
table's \emph{structure} says agglutinative and prefixing, but the
\emph{label} says Semitic, and the label wins because the desired
conclusion (a Semitic prayer) is a fixed prior rather than an emergent
reading. And it is a \textbf{free-parameter / identifiability} failure ---
an unanchored sign-value tuned to the target, on a corpus with no known
cognate language to constrain it. logos indexes the reading (*301 = /na/)
precisely so a model cannot regurgitate it as discovery, and quarantines
rather than dismisses it. The lesson generalises: a match rate on a
hand-picked word, absent a committed prediction and multiplicity
accounting, is not evidence. The *301 claim is a qualitative illustration
of the failure modes the gate is designed to detect.

\section{Results I: morphology induction}

We test whether Linear A's pre-registered affixes and word divisions
reflect genuine morphology. The apparent affixation does not clear the
operative \texttt{L\_fake} bigram floor, so the verdict is NO POWER and
\texttt{has\_validated\_morphology} = False. The identical test
\textbf{validates on Mycenaean Greek (Linear B)}, whose longer words carry
uncontested inflection: it confirms affixation at rate 0.562 --- above the
same bigram floor ($\sim$0.30) --- with \texttt{has\_morphology\_power} =
True, so the control rules out a universally insensitive detector and
supports the reading that Linear A's shorter forms contribute materially to
its lack of power (Table~\ref{tab:morphrates}; Appendix~A.1). One held-out positive stands and survives
its own uncertainty: a DP-unigram segmenter recovers the scribe's own word
divisions at micro-F1 0.436 versus a random-boundary baseline of 0.389, a
gap whose site-clustered 95\% bootstrap CI [0.021, 0.099] excludes zero
(P(gap $>$ 0) = 0.998; Appendix~A.1). This
separates real recoverable structure from fitted morphology --- the paper's
core discipline.

\section{Results II: metrology, phonology, and the cross-script asymmetry}

Three further pre-registered probes exercise the harness against distinct
field claims. Each returns a calibrated null or a truth-capped positive;
together they map where the Linear A corpus does and does not carry
recoverable signal, and they are consistent with identifiability rather
than effort being the more fundamental present constraint.

\subsection{Metrology: a parse-coverage null that agrees with the field}

Direction D poses the accounting tablets as a constraint-optimisation
over the fraction-sign values (Appendix~B.1; Table~\ref{tab:metro}).
Held-out fraction balance
does not beat the shuffle null
($p = 1.0$), so the verdict is NULL; the only value the method determines,
the fraction sign J = 1/2, is the long-standing field consensus
\citep{bennett1950}, not a logos discovery. This is a null the field itself
reached \citep{corazza2021} --- the harness crediting a surviving value to
its originator rather than re-branding it.

\subsection{Distributional phonology: a data-limited null bounded by a
power curve}

Does a sign's distributional context (which signs sit beside it) carry
information about its phonetic value? Tested on the known AB signs, neither
the consonant class nor the vowel class clears chance (permutation
$p = 0.40$ / $0.86$). The verdict is a data-limited NULL: the sound path is
not recoverable from Linear A alone, and no sound value is imputed for any
sign (Table~\ref{tab:phon}; Appendix~B.2).

\subsection{The cross-script asymmetry: image alignment recovers shared
anchors, but the recovery is circular by construction}

Do image embeddings recover the cross-script A$\leftrightarrow$B
correspondence that sequence and cognate co-occurrence cannot? Classical
shape descriptors reach direct-NN \textbf{0.410} on held-out
A$\leftrightarrow$B anchors while sequence and cognate co-occurrence stay
at chance. But the recovery is circular by construction: the
A$\leftrightarrow$B anchor values were themselves assigned by
twentieth-century epigraphers on sign-shape similarity, so this is a
shape-genealogy engineering demonstration, truth-layer-capped $\le$
\textbf{0.75}, not a positive decipherment claim (\S9.6;
Table~\ref{tab:xscript}).

\section{Results III: contamination probes}

One failure mode is invisible to the false-positive / multiplicity axis of
\S7: a decipherment that \emph{looks} novel but silently reproduces
published readings a language model has memorised. This section reports the
two decontamination probes that target it --- both run as method
demonstrations, returning completed (if deliberately bounded) findings.
Throughout, where a clean number was not obtainable we report the
partial-null and the confound rather than a headline.

\subsection{LLM-ablation: does a large open-weight model regurgitate the
literature?}

In the LLM-ablation, \texttt{qwen2.5:72b} reproduces famous
published readings --- Gordon's JA-NE = ``wine'' in 14/40 seeds --- that
the model-free edit-distance baseline never reaches (0/40), yet never
reproduces obscure vessel equations (0/40; Table~\ref{tab:repro}).
Verdict: a reproduction
gradient consistent with differential public exposure --- though the
experiment does not isolate memorisation from correlated item-level
differences --- the memorisation confound the discipline layer must gate
out.

\subsection{\texttt{L\_not\_indexed}: indexed versus not-indexed
abstention}

Does the model's behaviour on \emph{published} signs generalise to signs
unseen in the indexed literature? \texttt{qwen2.5:72b} valued the
memorisable known signs in $\sim$20/40 seeds but abstained $\sim$97\% of
the time on the literature-unseen *-series signs ($\sim$1/40;
Table~\ref{tab:abstain}) --- an
abstention asymmetry consistent with differential public exposure (the
\S8.1 caveat applies). With no ground truth on the not-indexed signs,
\textbf{consistency is not correctness}: it bounds how the model behaves,
never whether it is right.

\section*{Registered extension: a known-answer sufficiency curve}
\addcontentsline{toc}{section}{Registered extension: a known-answer
sufficiency curve}

Is a Linear-A-scale corpus large enough to recover a decipherment if a
known cognate existed? Because every benchmark is handed a real cognate
signal Linear A lacks, any recovery is an \textbf{optimistic benchmark,
conditional on a known cognate relationship}: at-floor at Linear-A scale
$\Rightarrow$ information-insufficient even given a cognate; above-chance
$\Rightarrow$ identifiability (cognate-absence, uncorrected multiplicity)
rather than corpus size would be implicated; neither branch is a
decipherability claim. The completed 2{,}000-step curve (near-converged:
Luvian$\to$Hittite 0.462 vs.\ published 0.475, Phoenician$\to$Ugaritic
0.760 vs.\ 0.805; every recovery a \emph{lower bound}) splits the
branches at the sweep's Linear-A-scale locator (600--700 distinct
word-forms, syllabic normalization): Linear~B$\to$Greek, the primary
analog, stays at its chance floor ($0.044\pm0.023$, $n{=}276$; floor
$\approx0.004$), while Cypriot$\to$Greek sits marginally above
($0.066$, $n{=}346$). Under-convergence keeps the at-floor branch
non-definitive; full per-size and per-cell artifacts are deposited
(Appendix~C.3).

\section{Discussion, limitations, and conclusion}

\subsection{The discipline is the contribution}

The deliverable of this work is not a reading of Linear A. It is a
\emph{harness} --- a literature index with an indexed versus
literature-unseen partition, a fabricated-language control corpus,
model-ablation contamination checks, multiple-comparison correction, and
pre-registration, all graded mechanically from persisted artifacts ---
together with a body of pre-registered, calibrated nulls, each with its
borderline case named. The stratified morphology re-run is the template: at
the pre-registered $\alpha = 0.01$, no pre-registered affix is cross-stratum
stable above the bigram floor, robust to the genre-bucketing choice
(Tables~\ref{tab:strata}--\ref{tab:sens}; Appendix~A.1). A null that
names its threshold, its sensitivity, and its
near-miss is a stronger referee object than a flat negative.

The axis this harness targets is the \textbf{false-positive / multiplicity
axis} --- the mechanism by which a determined analyst ``translates'' almost
anything on a small corpus --- a property of inference methodology, not of
the object or its data (\S1--2).

\subsection{Two named failure modes for the methods literature}

\textbf{(i) Narrative capture.} Structure says X; the label says the
desired Y; the label wins. In the June-2026 *301 claim the published table
segments an agglutinative, prefix-stacking verb while the assigned family
label is Semitic, and the desired output operates as a fixed prior, so
the conclusion conforms to the prior rather than emerging from
constrained readings (\S5). This
is distinct from statistical overfitting and must be named separately.
Notably, the mismatch is empirical --- between the label and the primary
descriptions of the form (\S5) --- not an a-priori typological
impossibility.

\textbf{(ii) Free parameters exceeding the information floor.} A lexicon
reported as ``408 terms and counting'' on a fixed corpus has more degrees
of freedom than the data can identify; a decipherment whose lexicon grows
weekly is underdetermined by construction. The minimum-description-length /
unicity \emph{diagnostic} (\texttt{k $\le$ U\_floor}) is built to flag
exactly this (it is reported next to every claim rather than enforced as a
graduation clause; \S4): the single sign the whole
reading pivots on is the textbook free-parameter case (\S5), its value
assigned to fit the desired root; the refutation rests on the claim's own
published table, archived in \artobj.

\subsection{The post-hoc evaluation sequence}

These failure modes trace to sequencing. The *301 claim publicly commits to
a post-hoc evaluation sequence: declare ``officially deciphered,'' then vet
later. logos's documented inverse is \textbf{pre-register $\rightarrow$
held-out $\rightarrow$ locked artifact} (\S3.1, \S3.4; the held-out
libation-formula corpus and its low-power caveat are in Appendix~A.3).

\subsection{Identifiability appears the more fundamental present constraint}

We take care not to write a defeatist framing, and equally not to
over-claim from corpus size: the toy-model unicity distance
($U \approx 204$--$415$ signs; Table~\ref{tab:denoms}, \S4) is a
\emph{precondition} diagnostic
only, explicitly withdrawn as any sufficiency or refutation claim ---
``$U \ll N$'' does \textbf{not} license the statement that there is enough
text to pin a decipherment (\S4; Appendix~C.4). The
results are consistent with \emph{identifiability} being the more
fundamental present constraint --- the absence of a known cognate language
to map to, an uncorrected multiple-testing search burden, and a mixed sign
inventory --- rather than a demonstrated corpus-size threshold; a null can
arise from identifiability, low power, or size alike, and we do not
establish that more data would be irrelevant.
\citet[\S8]{salgarella2025}
reaches the field-level version of this by a different route (\S2.2).

\subsection{Study-wide error budget}

Multiplicity discipline operates at two levels, and we are explicit about
both. Within a single probe, matches are deflated for everything
searched (\S3.2--3.4). Across the study, the unit of multiplicity is the
\textbf{per-probe pre-registration}: the \textbf{five executed probes}
(morphology, metrology, phonology, cross-script image alignment, and
contamination) each carry their own pre-committed $\alpha$, and every probe
is reported regardless of outcome. \textbf{Error control is defined within
each pre-registered probe; we claim no global family-wise guarantee across
probes.} Nothing in the battery is read as a decipherment.

\subsection{Limitations}

The nulls are honestly bounded, not decisive. A dependence-adjusted
trial-count correction of the morphology z-test is deferred; the deferral
would only strengthen the null (Appendix~A.1). No held-out
\emph{decipherment} has
been achieved, and the study reports \textbf{no positive decipherment
claim} --- the deliverable is calibrated absence, not a reading. Nothing
here claims a rejected reading is \emph{false} (\S3.4): every REJECT /
NULL in this study bounds \emph{evidential support}
on held-out data, not truth. The one probe that recovers an above-chance
number, cross-script sign-image alignment, is a font-robust but circular
\emph{engineering demonstration} (a font-swap control refutes the
typeface-artifact worry; \S7.3, Appendix~B.3), truth-capped at
$\le 0.75$, never a reading. No
phonetic claim is made. The philological points in \S5 rest on the
published literature cited there rather than on independent
Aegean-epigraphic
review, which we disclose as a limitation; because those points are cited
rather than original, the argument does not depend on any new sign reading
of our own. Chronological cross-validation is largely foreclosed by the
single-horizon corpus (\S4; Appendix~A.1). And the corpus itself
--- small, partly lossy, with no bilingual yet found
\citep[\S8]{salgarella2025} --- sets the ceiling.

\subsection{Conclusion and reproducibility}

The framework withholds validation when a proposed reading does not clear
the pre-registered gates (\S1, \S3); the contribution is the
framework and the calibrated null itself --- the instrument and the honest
limits of what the corpus can support --- not a reading, and the calibrated
null is a safeguard against unsupported claims.
\ifanon
The code, pre-registrations, and all findings are open under the MIT
licence and available as an anonymized artifact (code and pre-registrations
released on acceptance).
\else
The code, pre-registrations, and all findings are open under the MIT
licence and archived on Zenodo (concept DOI 10.5281/zenodo.21087572;
repository \url{https://github.com/papadopouloskyriakos/logos}).
\fi
The licensed epigraphic corpora are not redistributed; each source's
licence and retrieval path is documented (data-provenance; Software \&
data note), so the harness is reproducible against a
locally obtained corpus. This work is independent and not affiliated with,
or endorsed by, any of the cited projects or their authors.

\bibliographystyle{acl_natbib}
\nocite{*}
\bibliography{refs}

\appendix

\section{Morphology (replication detail)}

\paragraph{A.1 Morphology: rates, floors, word-length, stratification,
sensitivity.} The probe is pre-registered, unsupervised morphology and
segmentation induction over the full Linear A corpus
(Table~\ref{tab:denoms}). Target claims are the recurring affixes and
positional grammar of the Davis/Thomas/Duhoux--Val\'{e}rio tradition:
prefixes and suffixes (\texttt{i-}, \texttt{a-}, \texttt{-te/-ti}, and
Davis's \texttt{i-*301} verb stem). Pre-specification anchoring (\S3.1):
pre-specifications were frozen in-repository before each run (commit
history); the external immutable deposit
\ifanon
([artifact DOI withheld for review])
\else
(Zenodo concept DOI 10.5281/zenodo.21087572)
\fi
certifies the current state of the registrations, not temporal priority.
Each pre-registered affix
(drawn from \citet{thomas2020} and Duhoux/Val\'{e}rio) must recur on
$\ge$2 distinct stems across independent inscriptions, above a
within-form permutation null, and survive multiplicity correction at a
target of \emph{p} $< 0.01$ corrected. Every candidate was graded against two
null floors: a within-word sign-order permutation (shuffle) floor, and an
\texttt{L\_fake} floor built from a first-order sign-bigram (Markov) corpus
with zero lexical overlap with the real inscriptions.

\emph{Pooled run.} Of the 16 pre-registered affixes, 4 beat the within-form
permutation null and the deflated multiplicity bar (\texttt{i-*301},
\texttt{a-}, \texttt{i-}, \texttt{-te/-ti}), clearing the shuffle floor
decisively (rates and Wilson intervals in Table~\ref{tab:morphrates}).
The apparent standout \texttt{i-*301} sits far above the permutation null
--- but that figure is computed against the weaker within-form permutation
null only, and it does not survive the operative \texttt{L\_fake} bigram
floor, under which the whole set fails: the fabricated Markov corpus
confirms the same affixes at a \emph{higher} rate than the real corpus
(Table~\ref{tab:morphrates}). The outcome is classified NO POWER, not
evidence against morphology; \texttt{has\_morphology\_power} is False and
the module reports the null (the Linear-B positive control that validates
the detector is reported below).

\emph{Word-length distribution.} Linear A words are short
(Table~\ref{tab:morphrates}), and on such words ``morphology'' and
``bigram order'' are statistically indistinguishable. This reconciles
with Fuls's 3.3-sign figure once the denominator is matched:
\citet[57, \S6.5]{fuls2015} computes the mean over a word \emph{list} of
``554 complete sign sequences, where duplicates are reduced to one''
(GORILA; \citealp{godart1985}) --- distinct word \emph{types}, not
tokens; restricting our corpus to distinct multi-sign types likewise
recovers 3.07. Fuls reads the 3.3-sign length as evidence that Minoan is
\emph{polysynthetic}; our direct test finds the apparent affixation
bigram-reproducible with no power, and we report the null. The direct
positive control --- the identical bigram-floor test on Mycenaean Greek
(Linear B), whose longer words carry uncontested inflectional morphology
--- was RUN on the D\=AMOS corpus (parsed by the same
\texttt{\_damos\_wordforms} extractor as \S7.3) against a pre-registered
16-affix Mycenaean panel \citep{ventris1953}. It fires
(Table~\ref{tab:morphrates}), carried by the securely-grammatical
suffixes -jo, -o-jo (gen sg -oio), -de (allative), -qe (enclitic
``and''), -o-i (dat pl), so \texttt{has\_morphology\_power = True}. The
identical test that returns NO POWER on Linear A's short words recovers
Mycenaean morphology on longer ones (a sensitivity check: an ad-hoc
parser yielding mean 2.19 signs/word did \emph{not} fire --- power tracks
word length exactly as the short-word argument predicts), which rules out
a universally insensitive detector and supports the interpretation that
Linear A's shorter forms contribute materially to its lack of power
(\texttt{scripts/\allowbreak comparison/\allowbreak linb\_morphology\_control.py}).

\emph{Segmentation positive.} Under leave-one-site-out cross-validation (52
sites; not k-fold, given formulaic dependence), a DP-unigram
(Goldwater-style) segmenter recovers the scribe's own word divisions above
a random-boundary baseline; a site-clustered bootstrap (resample the 52
sites with replacement, 4,000 draws) puts the real-minus-random gap above
zero, and the paired gap, not the wide marginal per-segmenter CIs, is the
operative test (Table~\ref{tab:morphrates}). Morfessor over-segments the
short syllabic corpus and is reported as not usable rather than hidden
(Table~\ref{tab:morphrates}).

\emph{Stratified re-run.} The strata reconcile exactly to the
1,341-inscription corpus total (Table~\ref{tab:strata}). The
99-inscription ``libation'' \emph{genre} stratum is not the same object
as the libation-\emph{formula} corpus (\S A.3). Following
\citet{salgarellajudson2024}, chronological cross-validation was explicitly
declined (988/1,341 $\approx$ 74\% single-horizon LM IB), and site --- not
scribal hand, which is not individuated for Linear A
\citep{salgarella2019} --- is the independence unit, gated by presence at
$\ge$2 distinct sites. The result is \texttt{has\_validated\_morphology} =
False (\S9.1). Stratification also showed
that even \texttt{i-*301} is bigram-reproducible in the repetitive libation
register, where the Markov corpus manufactures the apparent verb
morphology; Davis's slot grammar is not refuted, merely not separable from
sign bigrams here. Where formula-bearing sites drive a signal, power is
intrinsically low: a leave-one-site-out CV over so few formula findspots
(an effective $\sim$5-fold site partition) cannot resolve a modest effect.

\emph{Sensitivity, disclosed.} The honest 2$\times$2 over genre bucketing
$\times$ $\alpha$ threshold is Table~\ref{tab:sens}; the doubly-relaxed
corner's \texttt{i-} (locative) and \texttt{-te/-ti} (ablative) are
precisely the Duhoux/Val\'{e}rio H3 nominal affixes that
\citet{salgarella2025} \S8 independently endorses --- the named borderline
candidates, fragile to \emph{both} the $\alpha$ threshold and a defensible
genre choice (the $\alpha = 0.01$ null is itself robust to the bucketing).
The present analytical approximation uses attempted-candidate
count rather than a dependence-adjusted effective trial count --- a
recorded, deferred limitation whose correction would only strengthen the
null.

\paragraph{A.2 \texttt{S\_morph} detector calibration (\S3.2).} After the productivity-gate fix (\S3.2 --- an affix must attach to
$\ge$2 distinct residual stems), \texttt{S\_morph}'s measured false-positive
rate on within-form-shuffled corpora is consistent with the nominal
one-sided $z \ge 2$ rate within sampling error, and its planted-affix
power is a saturated point, not a full sensitivity curve
(Table~\ref{tab:calib}; contrast the phonology power curve, \S7). This is
an \textbf{indicative calibration check, not the operative gate}: 6/200
is too coarse to certify an $\alpha = 0.01$ threshold on its own, and the
gate applies the order-statistic bar over the far larger permutation and
candidate counts the deflation machinery tracks (\S3.4).

\paragraph{A.3 Held-out design: the libation-formula corpus (\S9.3).} A
leave-one-site-out split over so few formula sites (\S1) has
low statistical power --- and site-level CV is in any case necessary, not
sufficient, since within-site scribal-hand correlation must also be controlled
\citep{salgarellajudson2024}. This formula corpus should not be conflated with
the 99 stone-vessel inscriptions of the ``libation'' genre stratum (\S6): they
are different objects.

\section{Metrology, phonology, and cross-script (methods and controls)}

\paragraph{B.1 Metrology.} Direction D parses HT (LM I) documents into
balance units \texttt{[recipient][commodity logogram][integer][fraction
signs]} and solves the fraction-sign values by RANSAC over exact-\texttt{Fraction}
Gaussian elimination (robust to the unbalanced-tablet outliers that defeat
least-squares), requiring line items to balance the preserved KU-RO totals;
validation is grouped k-fold by document (no tablet straddles train and
test) against a permutation null. Results are in Table~\ref{tab:metro}: of
$\sim$7 fraction-bearing constraints only one is satisfied by the solved
system, the sole determinable value is J = 1/2, held-out fraction balance
is not separated from the permutation null, and the synthetic positive
control (planted values, all recovered) separates. The null traces to sparse
automated parse coverage (multi-commodity blocks, illegible number glyphs,
KI-RO deficit sub-lists, editorial restorations); \citet[p.~2]{corazza2021}
abandoned the balance-to-totals method for the same reason. The surviving J
= 1/2 (\S7.1) gains Schrijver's
\citeyearpar{schrijver2014} three further tablets (HT Zd 156, PE 1, HT
123a): at least four corroborating documents.

\paragraph{B.2 Distributional phonology.} Track 2 asks whether a sign's
distributional context --- which signs sit beside it --- carries
information about its phonetic value, measured on the known \texttt{AB}
signs whose Linear B sound values can be borrowed. Features are
PPMI-weighted, L2-normalised left/right co-occurrence counts; labels are
the Linear-B-transferred (C, V) parsed from the GORILA name, so features
never see the labels. The test is leave-one-out 1-NN against a
label-permutation null, \v{S}id\'{a}k-deflated over the two tests; 50 of
$\sim$55 AB signs met the $\ge$3-occurrence threshold.

Per-test results are in Table~\ref{tab:phon}. The honest verdict is
\emph{data-limited},
not \emph{no bridge exists}, because the positive control disciplines the
reading: the positive-control power curve over planted strengths shows
the test firing at strength $\ge$ 2 but not at 1. The test
works but is not infinitely sensitive at n = 50 / 13 classes, so we cannot
distinguish ``no distribution$\rightarrow$phonetics bridge'' from ``too few
signs to detect a weak one'' (\S7.2, \S B.3).

\paragraph{B.3 Cross-script image alignment.} The palaeographic probe
renders 88 Linear-A and 74 Linear-B glyphs and tests held-out recovery of
shared A$\leftrightarrow$B anchors. Classical descriptors (HOG + Hu moments
+ shape-context) sit far above image chance, and a from-scratch I-JEPA is
excluded from the claims as over-parameterized for $\sim$90 glyphs
(results and controls in Table~\ref{tab:xscript}). Two controls decide how
to read the image-vs-sequence asymmetry (\S7.3), and together they
\emph{downgrade} it from a positive to a demonstration.

\emph{First, a font-swap control.} Rendering both scripts in one typeface
(Aegean, George Douros) risks making the alignment a single-designer
house-style artifact rather than a fact about the signs. We re-render Linear
A in Noto Sans Linear A and Linear B in Noto Sans Linear B --- two
independent type designers, each font covering only its own Unicode block
(Noto-A carries zero Linear B glyphs and vice-versa), so no single hand
draws both scripts. The recovery survives and is if anything stronger
(Table~\ref{tab:xscript}; both legs $\gg$ chance): the font-artifact
hypothesis is \emph{refuted}, and ``distances reflect shape, not font'' is
now demonstrated rather than asserted
(\texttt{scripts/\allowbreak palaeo/\allowbreak font\_control.py}).

\emph{Second, and decisively, the circularity the font control cannot
remove.} Linear B was historically adapted from Linear A, and the
A$\leftrightarrow$B transcription \emph{values} that define the anchor set
were assigned by twentieth-century epigraphers largely on the basis of
sign-shape similarity, so the recovery re-derives a shape-defined
correspondence (\S7.3): it confirms a known palaeographic fact (borrowed
signs look alike) and validates a font-robust descriptor pipeline, but it
reads no Minoan. It remains truth-layer-capped in any case --- the
\citet{ferrara2021} ``Archanes Formula'' hard negative (\texttt{CH 095}
resembles \texttt{AB 60}/\emph{ra} but is not its phonetic counterpart)
shows a plausible shape match can be phonetically wrong, so a shape match is
a $\le 0.75$ signal (invariant 5), never a reading. The honest asymmetry is
thus: the image path re-derives a shape-defined mapping (expected, circular,
capped), while the sequence and cognate paths \textbf{remain at chance even
after ingesting the full D\=AMOS Mycenaean corpus} --- 13,562 wordforms,
$\approx$15$\times$ the earlier cognate file, 56 anchors
(Table~\ref{tab:xscript}) --- while the \emph{positive control
recovers}, so the harness works and the null is a real null: the
A$\leftrightarrow$B phonetic correspondence is simply not carried by
cross-script sign co-occurrence, even at full Linear-B scale
(\S B.2). And the A-only signs on which
positive-decipherment claims would pivot have no cross-script
anchor at all (\S5), so images cannot impute them even in principle.

\section{Contamination and sufficiency (detail)}

\paragraph{C.1 LLM-ablation reproduction gradient.} The ablation
(\S3.3) is the remaining gate on the contamination number, matching a
model's cognate glosses against the indexed lexical claims. We ran it
against \texttt{qwen2.5:72b} --- a large open-weight model, a far
stronger memoriser than a local 12B --- on a rented H100, with an
NW-Semitic candidate and a Hebrew skeleton lexicon for the model-free arm
(scale and quantities in Table~\ref{tab:canary}) (\artref). The
first-pass contamination rate is an \textbf{artifact, reported as such}:
the mechanical arm never shares a literature-matching correspondence with
the LLM (their value vocabularies barely overlap), so the rate is forced
to 1.0 by construction and is not reported as a contamination result
(Table~\ref{tab:canary}; \artref).

\emph{Generation settings (replication).} Checkpoint \texttt{qwen2.5:72b}
as served by Ollama (Ollama's default \texttt{q4\_K\_M} quantization;
Ollama version not recorded), on an ephemeral vast.ai H100 SXM 80\,GB
instance (torn down after the run). One \texttt{/api/generate} call per
seed, stream off, timeout 180\,s; temperature 0.8; sampling seed = the
replicate seed; top-p, top-k, maximum output length, and stop conditions
left at server defaults (not recorded). Each
seed's 40 held-out forms are drawn deterministically
(\texttt{numpy.default\_rng(seed)}) and rendered
as canonical sign names in a fixed JSON-request prompt
(\texttt{\_PROMPT\_TEMPLATE} in
\texttt{scripts/\allowbreak comparison/\allowbreak ablation.py}). Parsing:
the JSON \texttt{partial\_map} is extracted; keys not naming a sign shown
to the model are dropped (the proposer cannot invent signs); a dead or
garbled call contributes an empty proposal without crashing the batch
(fail-closed). The gate-2 rerun
prepends the 7 indexed probe forms to every seed. Per-call logs:
\texttt{runtime/ablation/llm\_calls.jsonl}; harness commits
\texttt{0582966} (ablation) and \texttt{9cd6d04} (gate-2).

The gate-2 rerun, reconciled to a common consonant-skeleton vocabulary,
yields the honest deliverable: the per-item reproduction gradient over 40
seeds (Table~\ref{tab:repro}; \artref). An
external red-team correction split \texttt{KU-RO} between the
Semitic-root route (\emph{kull}) and the general-knowledge meaning
``total'' (confirmed by tablet arithmetic), so \texttt{KU-RO} = ``total''
is \textbf{excluded} as a witness and \texttt{KU-RO} = \emph{kull} is
downweighted (Table~\ref{tab:repro}; \artref). The surviving gradient ---
famous readings reproduced, obscure ones not --- is what the ablation
detects (\S8.1). The gate-2 rate itself remains confounded by the baseline's
representational ceiling and is not reported as a contamination measure
(Table~\ref{tab:canary}).

\paragraph{C.2 \texttt{L\_not\_indexed} abstention asymmetry.} The
\texttt{L\_not\_indexed} probe shows a sign in real context (\S8.2); the
not-indexed signs are absent from the indexed
sources, hence not memorisable from them (never ``impossible to have seen''
in any absolute sense) (\artref). Across 40 seeds, \texttt{qwen2.5:72b}
abstains $\sim$97\% of
the time on the untouched signs while confidently valuing the memorisable
ones (Table~\ref{tab:abstain}). The apparent known$-$not-indexed
consistency gap is coverage-confounded and was flagged, not asserted: the
matched-\emph{n} robustness check returned \texttt{no\_power}, so the
rigorous comparison could not run (Table~\ref{tab:abstain}).
The interpretable signal is the abstention asymmetry itself, read under
the standing rule of \S8.2 (\artref).

\paragraph{C.3 Information-sufficiency CSA sweep (design and deposit).}
The sweep downsamples known-answer benchmarks --- Linear B /
Mycenaean Greek as the primary syllabic analog, sourced through D\=AMOS,
and further syllabic corpora --- and measures held-out cognate-ID
recovery against per-size permutation floors, using the published
cognate-ID matcher of \S2.3 \citep{tamburini2025}, validated in the
artifact against its published Table-3 values (\artref). The reading
rules and the reported curve are in the body (Registered extension); the
full per-size report and per-cell artifacts are deposited in the
repository (\texttt{results/csa/}).
The phonological / sound path is no longer data-deferred: even the full
ingested D\=AMOS corpus leaves the cross-script distributional alignment
at chance (\S7.3, \S B.3).

\paragraph{C.4 Information floor: corpus denominators and the unicity
computation (\S4).} The corpus envelope, the four
sign-count denominators, the two occurrence channels, and the measured
unicity range are tabulated in Table~\ref{tab:denoms}
\citep[Schoep 2002;][\S1--3, \S6]{salgarella2025}; the envelope spans ca.
1800--1450 BCE (Middle Minoan II--Late Minoan IB; \S4). Each denominator
must be labelled where it is used (Table~\ref{tab:denoms}). The channel through
which structure is observed is a \emph{sign-occurrence} count, not a
document count or a sign inventory (invariant 7): the broader $\sim$7,400
occurrences count every sign type, while the narrower
$N \approx 5{,}792$ parsed syllable-signs --- after logograms and
numerals are stripped --- is what the unicity computation consumes. Two
further properties tighten the budget: Linear A words are short and
formulaic (the six-position libation-formula template;
\citealp[\S7]{salgarella2025}), and a single language/horizon dominates.

On the logos-normalized corpus, Shannon's unicity-distance criterion
yields $U \approx 204$--$415$ signs across the whole plausible $(V, P)$
parameter space, everywhere far below $N$ (Table~\ref{tab:denoms}) ---
retained only as a toy-model precondition diagnostic (\S4). Unicity
distance presupposes a \emph{known plaintext-language oracle}; because we
estimate the redundancy $D$ from the ciphertext itself, that oracle
collapses into Linear A up to relabelling, so the computation
demonstrates only measurable recurrent structure, not that any reading is
identifiable; and a lossy, non-injective sign$\rightarrow$sound channel
may possess no unicity distance at all (\S9.4).

\paragraph{C.5 Literature index and \texttt{L\_fake} canary (\S3.3).} \emph{Index seed and adversarial population.}
The seed is deliberately conservative --- a sign wrongly omitted is wrongly
granted \texttt{L\_not\_indexed} status, the dangerous direction --- and
every verifiably-published disputed entry is indexed to \emph{catch}
regurgitation, never as an
accepted reading. Population is adversarial: of the seeded West-Semitic
lexical proposals, five are Gordon's equations and a-sa-sa-ra-me = ``oh
Asherah!'' is \citet{best1981}, not Gordon; provenance-unverifiable
candidates --- \texttt{qa-pa} = kappu and a Semitic \texttt{ki-ro}, whose
attribution to Gordon is not supported by the primary account
\citep{rendsburg1996} --- were excluded rather than indexed. Exposure
indexing is independent of evidential validity --- a false public claim
can still contaminate. Generated from the
seed (invariant 12), the current census (Table~\ref{tab:canary}) is a
deliberately conservative, non-exhaustive seed, not yet the full \S C.1
literature census.
\emph{\texttt{L\_fake} calibration.} The canary is calibrated to a real
candidate language's phoneme frequencies, root-template structure, and
lexicon size (\S3.3). Its own divergence from the calibration
targets is reported, never minimised; with a Semitic root-template sampler
and an external Hebrew (ETCBC/bhsa) reject set, the root-template
total-variation distance fell markedly (Table~\ref{tab:canary}), and any
residual real-lexicon collision is measured and published rather than
asserted to be zero.

\section{The order-statistic graduation bar: formula, assumptions,
calibration}

The operative deflation clause (\S3.4) is a normal-order-statistic
\textbf{approximation} to the expected maximum of the multiplicity-null over
the attempted candidates. For a null of mean $\mu_0$ and standard deviation
$\sigma_0$ and $n$ attempted candidates,
\[
  \resizebox{0.95\columnwidth}{!}{$\displaystyle
  E[\max] \approx \mu_0 + \sigma_0\bigl[(1-\gamma)\,\Phi^{-1}(1-\tfrac{1}{n})
  + \gamma\,\Phi^{-1}(1-\tfrac{1}{ne})\bigr]$}
\]
with $\gamma = 0.5772\ldots$ the Euler--Mascheroni constant --- the
blended expected-maximum expression used by \citet{bailey2014}; the
general order-statistics treatment is \citet[\emph{Order Statistics}, 3rd
ed.]{david2003}. The
pipeline exposes the bar as
\texttt{expected\_max\_order\_stat}($\mu_0, \sigma_0, N$); the
deduplicating \texttt{SearchLog}
(\S3.2, invariant 12) reports the exact distinct-candidate count
\texttt{n\_trials} a scanner produces. The bar fails closed on a degenerate null
($\sigma_0 \le 0 \Rightarrow E[\max] = \mu_0 \Rightarrow$ no graduation).
The \texttt{generalizes\_to\_not\_indexed} clause (\S3.4) requires recovery
above a pre-committed threshold and on a minimum number of distinct
not-indexed signs, and fails closed when no threshold or no supported
not-indexed sign is present. The full clause conjunction: the candidate
must separate from the \texttt{L\_fake} negative-control distribution,
its corrected margin must be cleared, and the observed statistic must
beat the order-statistic bar over the Packard/random-lexeme multiplicity
null.

\textbf{Assumptions and the direction of their error.} The approximation is
exact only under (i) an approximately Gaussian null and (ii) independence
among the $n$ attempted candidates. Neither holds exactly. (i) The
permutation / \texttt{L\_fake} nulls are not perfectly Gaussian; the blended
Euler-$\gamma$ form is a first-order correction, not a guarantee. (ii) The
candidates a real search tries are typically \textbf{positively correlated}
(overlapping sign-value assignments). Under a common-variance Gaussian
comparison model, positive dependence lowers the expected maximum relative
to independence (a Slepian-type comparison; \citealp{slepian1962}), so
using the
raw attempted-candidate count $n$ in place of a smaller effective
independent-trial count places the bar \emph{above} the correlated truth;
this does \textbf{not} establish conservatism for arbitrary heterogeneous
or non-Gaussian dependence --- the empirical best-of-100 calibration
therefore carries the operational error-control claim. We use the
raw deduplicated candidate count; a dependence-adjusted effective count is a
recorded, deferred refinement (\S9.6) whose correction would only strengthen
the nulls.

\textbf{Calibration and its interval.} Because the analytic bar rests on
assumptions that do not hold exactly, the error-control claim is carried not
by the formula but by the direct empirical calibration (\S3.4): a
Monte-Carlo of the full gate on a best-of-100 random-map null-generating
process (\texttt{scripts/\allowbreak gate\_null\_calibration.py}). The
interval on that calibration is the \textbf{one-sided exact
Clopper--Pearson} bound (the Beta-quantile inversion of the binomial), which
needs no normal approximation to the calibration count --- for 3 graduations
in 500 replicates, a 95\% upper bound of $\approx 1.54\%$; the verdict
partition of the 500 replicates is in Table~\ref{tab:calib} (no
INCOMPLETE, the search being instrumented).
\citet{bailey2014} supply the selective-inference lineage (\S2.4) and the
blended expression above; the operational error-control claim nonetheless
rests on the empirical calibration, not on the formula.

\textbf{Cognate-ID baseline (\S3.4).} The match-rate the deflation is
applied to is supplied by the published cognate-ID matcher identified in
\S2.3 (reproduced in the artifact), so the gate does not depend on any
single citation.

\begin{center}\rule{0.5\columnwidth}{0.4pt}\end{center}

\noindent\emph{Software \& data.}
\ifanon
The artifact. Code, pre-registrations, and findings are released as an
anonymized artifact on acceptance. Licensed corpora (GORILA-, SigLA-,
lineara.xyz-, D\=AMOS-derived) are not redistributed; a data-provenance note
documents each source's licence and retrieval path.
\else
logos (v0.1.0). Zenodo. Concept DOI 10.5281/zenodo.21087572; repository
\url{https://github.com/papadopouloskyriakos/logos}. Licensed corpora
(GORILA-, SigLA-, lineara.xyz-, D\=AMOS-derived) are not redistributed; see
\texttt{docs/data-provenance.md}.
\fi

\section{Complementary tables}

\begin{table}[h]
\centering\small
\resizebox{\columnwidth}{!}{%
\begin{tabular}{@{}lrrrl@{}}
\toprule
\textbf{Stratum} & \textbf{Inscr.} & \textbf{Words} & \textbf{Sites} &
\textbf{Role} \\
\midrule
admin & 290 & 1,486 & 14 & induction \\
libation (stone vessels) & 99 & 370 & 15 & induction \\
other & 129 & 187 & 40 & induction \\
seal & 740 & --- & --- & excluded \\
emptied by abbrev.-strip & 83 & --- & --- & excluded \\
\midrule
total & 1,341 & & & \\
\bottomrule
\end{tabular}}
\caption{Morphology corpus strata (Appendix~A.1). Seals are structurally
excluded as an abbreviation channel; the strata reconcile exactly to the
1,341-inscription corpus total.}
\label{tab:strata}
\end{table}

\begin{table}[h]
\centering\small
\resizebox{\columnwidth}{!}{%
\begin{tabular}{@{}lll@{}}
\toprule
\textbf{Bucketing} & $\alpha=0.01$ & $\alpha=0.05$ \\
\midrule
strict (stone vessels only) & NULL & NULL \\
loose (stone objects + inked) & NULL & \texttt{i-}, \texttt{-te/-ti}
validate \\
\bottomrule
\end{tabular}}
\caption{Morphology sensitivity $2\times2$ (Appendix~A.1): genre bucketing
$\times$ significance threshold. Only the doubly-relaxed corner validates
the H3 nominal affixes \texttt{i-} (locative) and \texttt{-te/-ti}
(ablative).}
\label{tab:sens}
\end{table}

\begin{table}[h]
\centering\small
\resizebox{\columnwidth}{!}{%
\begin{tabular}{@{}ll@{}}
\toprule
\textbf{Quantity} & \textbf{Value} \\
\midrule
Linear A confirm rate (16 affixes) & 0.250 (4/16; [0.10, 0.50]) \\
\texttt{L\_fake} bigram-floor rate & 0.375 (6/16; [0.18, 0.61]) \\
\texttt{i-*301} vs.\ within-form null & $z \approx 8.0$ (6 environments) \\
\midrule
Linear B (D\=AMOS) confirm rate & 0.562 (9/16) \\
Linear B within-word-shuffle floor & $\sim$0.03 \\
Linear B \texttt{L\_fake} bigram floor & $\sim$0.30 (5 draws) \\
\midrule
word-token mean length & 1.84 signs (median/mode 1) \\
word tokens $\le$2 signs / 1 sign & 76.1\% / 56.5\% \\
multi-sign \emph{type} mean length & 3.07 (Fuls 3.3, 554-type list) \\
Linear B mean length & 3.23 (13,562 wordforms) \\
\midrule
DP-unigram segmenter micro-F1 & 0.436 \\
random-boundary baseline & 0.389 (base rate 0.406) \\
real$-$random gap, 95\% CI & [0.021, 0.099]; P(gap${>}$0) = 0.998 \\
Morfessor & 0.156 (over-segments; not usable) \\
\bottomrule
\end{tabular}}
\caption{Morphology-induction replication quantities (Appendix~A.1):
confirm rates with Wilson 95\% intervals, word-length denominators, and
the segmentation positive with its site-clustered bootstrap gap (4,000
draws over 52 sites).}
\label{tab:morphrates}
\end{table}

\begin{table}[h]
\centering\small
\resizebox{\columnwidth}{!}{%
\begin{tabular}{@{}ll@{}}
\toprule
\textbf{Quantity} & \textbf{Value} \\
\midrule
documents / balance units & 32 / 35 \\
integer-level balance & 7/28 (25\%) \\
determinable fraction value & J = 1/2 (HT104: $2 \cdot J = 1$) \\
held-out fraction balance (real) & 0.0 \\
permutation null & mean 0.113 (max 0.143); $p = 1.0$ \\
planted positive control & 0.9 vs.\ null 0.22; $p = 0.003$ \\
\bottomrule
\end{tabular}}
\caption{Metrology (Direction D) replication quantities (Appendix~B.1).
The real corpus is not separated from the permutation null; the planted
control separates, so the null is not a machinery bug.}
\label{tab:metro}
\end{table}

\begin{table}[h]
\centering\small
\resizebox{\columnwidth}{!}{%
\begin{tabular}{@{}lrrrr@{}}
\toprule
\textbf{Test} & \textbf{Classes} & \textbf{Acc.} & \textbf{Null} & $p$ \\
\midrule
consonant class & 13 & 0.080 & 0.065 & 0.40 \\
vowel class & 5 & 0.140 & 0.201 & 0.86 \\
\bottomrule
\end{tabular}}
\caption{Distributional-phonology per-test results (Appendix~B.2):
leave-one-out 1-NN vs.\ a label-permutation null, \v{S}id\'{a}k-deflated
over the two tests. The positive-control power curve fires ($p<.05$) at
planted strength $\ge 2$ but not at 1.}
\label{tab:phon}
\end{table}

\begin{table}[h]
\centering\small
\resizebox{\columnwidth}{!}{%
\begin{tabular}{@{}lll@{}}
\toprule
\textbf{Leg} & \textbf{Accuracy} & \textbf{Floor} \\
\midrule
classical, direct-NN & 0.410 [0.18, 0.64] & 0.0135 \\
classical, Procrustes & 0.185 [0.00, 0.36] & 0.0135 \\
I-JEPA (3 seeds; excluded) & $0.324 \pm 0.025$ & 0.0135 \\
font-swap, direct-NN & $0.367 \rightarrow 0.416$ & $\gg$ chance \\
font-swap, Procrustes & $0.156 \rightarrow 0.195$ & $\gg$ chance \\
sequence/cognate (D\=AMOS) & 0.025 & 0.011 \\
\quad positive control & 0.947 & --- \\
\bottomrule
\end{tabular}}
\caption{Cross-script sign-image alignment and controls (Appendix~B.3).
Font-swap rows give same-font $\rightarrow$ independent-fonts accuracy on
the control's matched anchor set; the sequence/cognate row is the best of
five alignment methods over 56 anchors
(\texttt{clearly\_above\_chance} = False).}
\label{tab:xscript}
\end{table}

\begin{table}[h]
\centering\small
\resizebox{\columnwidth}{!}{%
\begin{tabular}{@{}llrr@{}}
\toprule
\textbf{Item} & \textbf{Source} & \textbf{LLM} & \textbf{Baseline} \\
\midrule
\texttt{JA-NE} = \emph{yn} ``wine'' & Gordon & 14/40 & 0/40 \\
\texttt{A-SA-SA-RA-ME} & \citet{best1981} & 2/40 & 0/40 \\
\texttt{KA-RO-PA} & Gordon & 0/40 & 0/40 \\
\texttt{SU-PA-RA} & Gordon & 0/40 & 0/40 \\
\texttt{KU-RO} = \emph{kull} (downwt.) & Gordon & 8/40 & --- \\
\texttt{KU-RO} = ``total'' (excl.) & gen.\ knowledge & 7/40 & --- \\
\bottomrule
\end{tabular}}
\caption{LLM-ablation per-item reproduction over 40 seeds,
\texttt{qwen2.5:72b} vs.\ the model-free edit-distance baseline
(Appendix~C.1). The two \texttt{KU-RO} routes overlap on one seed; three
single-word items are among Gordon's five West-Semitic equations, while
\texttt{A-SA-SA-RA-ME} is Best's.}
\label{tab:repro}
\end{table}

\begin{table}[h]
\centering\small
\resizebox{\columnwidth}{!}{%
\begin{tabular}{@{}lrr@{}}
\toprule
\textbf{Sign set} & \textbf{Signs} & \textbf{Valued (of 40)} \\
\midrule
known AB (indexed) & 20 & $\sim$20 \\
not-indexed *-series & 11 & $\sim$1 ($\sim$97\% abstention) \\
\bottomrule
\end{tabular}}
\caption{\texttt{L\_not\_indexed} abstention asymmetry (Appendix~C.2).
The apparent consistency gap 0.601 (permutation $p = 0.0005$) is
coverage-confounded; the matched-$n$ check returned \texttt{no\_power}
(0 of 11 not-indexed signs reached $n \ge 10$).}
\label{tab:abstain}
\end{table}

\begin{table}[h]
\centering\small
\resizebox{\columnwidth}{!}{%
\begin{tabular}{@{}ll@{}}
\toprule
\textbf{Denominator / channel} & \textbf{Value} \\
\midrule
token-frequent syllabograms & 97 \\
full syllabary estimate & $\sim$150 \\
working repertoire ($V$; conflates syll./subscr./logograms) & 259 \\
incl.\ ideograms (GORILA: 390) & $\sim$350 \\
\midrule
sign occurrences, all types & $\sim$7,400 \\
parsed syllable-signs (unicity input) & $N \approx 5{,}792$ \\
\midrule
corpus envelope & $\sim$2,500 finds; $\sim$1,400 readable \\
findplaces (logos-normalized) & 56 ($\rightarrow$ 52 sites) \\
readable inscriptions & 1,341 \\
distinct word-forms & 539 \\
\midrule
unicity distance over $(V,P)$ space & $U \approx 204$--$415$ signs \\
\bottomrule
\end{tabular}}
\caption{Corpus denominators, occurrence channels, and the unicity range
(Appendix~C.4). The $\sim$7,400 all-type occurrence channel never enters
the unicity arithmetic, which consumes only the parsed syllable-signs.}
\label{tab:denoms}
\end{table}

\begin{table}[h]
\centering\small
\resizebox{\columnwidth}{!}{%
\begin{tabular}{@{}ll@{}}
\toprule
\textbf{Calibration quantity} & \textbf{Value} \\
\midrule
\texttt{S\_morph} FPR (within-form shuffle) & 3.0\% (6/200; [1.4, 6.4]\%) \\
nominal one-sided $z \ge 2$ rate & $\approx$2.3\% \\
\texttt{S\_morph} power (planted affix) & 100\% (50/50; saturated) \\
\midrule
gate false-graduation (best-of-100 null) & 0.6\% (3/500; CP $\le$ 1.54\%) \\
gate verdict partition (500 replicates) & 3 GRAD / 10 NULL / 487 REJ \\
\bottomrule
\end{tabular}}
\caption{Detector and gate calibration quantities (Appendices~A.2
and~D): Wilson interval for the \texttt{S\_morph} false-positive rate;
one-sided exact Clopper--Pearson 95\% upper bound for the gate's
false-graduation rate under the tested best-of-100 random-map null.}
\label{tab:calib}
\end{table}

\begin{table}[h]
\centering\small
\resizebox{\columnwidth}{!}{%
\begin{tabular}{@{}ll@{}}
\toprule
\textbf{Quantity} & \textbf{Value} \\
\midrule
literature-index census & 63 claims / 62 entries (55 single) \\
root-template TV distance & $\sim$0.84 $\rightarrow$ $\sim$0.33
  ($\sim$0.23 standalone) \\
\midrule
ablation scale & 40 seeds $\times$ 40 held-out forms \\
model-free lexicon & 7,682 Hebrew skeletons \\
LLM errors & 0 (40/40 seeds) \\
value-vocabulary overlap & 10 of 134 \\
first-pass contamination rate & 1.000 (33/33 defined; artifact) \\
gate-2 rate & mean 0.646 (confounded; not reported) \\
\bottomrule
\end{tabular}}
\caption{Literature-index, \texttt{L\_fake}-canary, and LLM-ablation
quantities (Appendices~C.1 and~C.5). The TV distance is measured
end-to-end on the Ugaritic$\leftrightarrow$Hebrew gold; the first-pass
contamination rate is forced to 1.0 by construction and reported as an
artifact.}
\label{tab:canary}
\end{table}

\end{document}
.

% ---- identity macros (switch on \ifanon) ----
\ifanon
  \newcommand{\artref}{the artifact}
  \newcommand{\artobj}{the artifact}
\else
  \newcommand{\artref}{logos v0.1.0}
  \newcommand{\artobj}{the logos artifact (v0.1.0)}
\fi

% ---- PDF metadata: named for arXiv, stripped for the anonymous TACL PDF ----
\ifanon
  \hypersetup{pdfauthor={},pdftitle={},pdfsubject={},pdfkeywords={},pdfcreator={},pdfproducer={}}
\else
  \hypersetup{pdfauthor={Kyriakos Papadopoulos},
    pdftitle={Falsification-First Decipherment: A Decontaminated Inference
Framework for Testing Undeciphered-Script Claims, with Linear A as a
Worked Null},
    pdfcreator={},pdfproducer={}}
\fi

\title{Falsification-First Decipherment: A Decontaminated Inference
Framework for Testing Undeciphered-Script Claims, with Linear A as a
Worked Null}

\ifanon
  \author{}
\else
  \author{
    Kyriakos Papadopoulos \\
    Independent Researcher \\
    Netherlands \\
    ORCID 0009-0003-0995-2518
  }
\fi

\date{}

\begin{document}
\maketitle
\ifanon\else
\begin{center}\small\emph{Preprint --- under review; not peer
reviewed.}\end{center}
\fi

\begin{abstract}
Decipherment research has a false-positive problem and no shared
instrument: a small corpus, an unknown language, and many candidate
readings let an analyst ``translate'' almost anything. The flagship
survey lists fifteen challenges of script and data; none addresses the
orthogonal inference-control layer --- overfitting, multiple testing,
contamination --- a candidate sixteenth challenge. We build a
falsification-first framework for it: a fabricated-language
negative-control family, multiple-comparison deflation, a literature
index separating published from not-indexed signs, and pre-registration
with a mechanically-graded gate. This is a methods paper; we assert no
decipherment. The machinery calibrates: on a tested null the gate's
false-graduation rate is 0.6\% (3/500); the identical morphology test
finding no power on Linear A recovers Mycenaean (Linear B) inflection
(9/16 affixes above the bigram floor), supporting the interpretation
that Linear A's short forms contribute materially to its lack of power.
Across probes it returns differentiated, calibrated outcomes --- null,
no-power, supported-structural, circular, and contamination-sensitive.
Results are consistent with identifiability, rather than corpus size,
being the more fundamental present constraint. The June-2026 *301
``cracked'' claim is a qualitative illustration of the failure modes the
gate is designed to detect.
\end{abstract}

\section{Introduction}

Decipherment has a false-positive problem, and the field has no agreed
instrument for measuring it. The history of undeciphered-script research
is, in large part, a history of uncontrolled model selection: a tiny
corpus, an unknown language, and a large space of candidate readings let
a determined analyst ``translate'' almost anything, with internal
consistency on the cherry-picked set standing in for confirmation. The
mechanism is not any single bad match but uncorrected multiplicity: with
enough candidate roots and a small enough corpus, a pile of
individually-plausible matches is the \emph{expected} outcome under pure
chance, so each individual comparison looks suggestive while the
conclusion is nonsense. A standard demonstration of spurious
comparison: one can ``prove'' \emph{English is a Semitic language} by
matching a handful of English words to Proto-Semitic roots while quietly
ignoring the far larger number that do not match. The same move, applied
to a real undeciphered script, has produced a long line of decipherment
claims that no one can either confirm or refute.

That the discipline lacks a shared standard here is not our editorializing
but the field's own assessment. The flagship computational-decipherment
survey --- \citeauthor{braovic2024}'s \citeyearpar{braovic2024}
systematic review of Bronze Age Aegean and Cypriot scripts ---
enumerates fifteen ``major challenges'' of script, language, and data
(\S2.1) \citep{braovic2024}. Those fifteen are
primarily properties of the \emph{object and its data} --- the script,
the language, the surviving corpus, and the resource/comparison-data
conditions. What
governs whether a \emph{claimed reading} is believable is a different
kind of thing: a property of the inference methodology and its
epistemics --- overfitting, multiple testing, false-positive control,
and contamination. We therefore frame this axis, once and modestly, as a
candidate \textbf{sixteenth challenge} --- a deliberate rhetorical foil
to the survey's fifteen, not a claim to a new taxonomic slot or to sole
discovery. It is \emph{under-instrumented} rather than unrecognised
(\S2.4): our contribution is to
\emph{consolidate and operationalize} that axis for decipherment, not to
invent it. It is the methodological axis the object-level taxonomy was
never meant to cover, and it is distinct from the survey's own
acknowledged evaluation gap (\S2.1) --- an \emph{evaluation-target}
problem, separable from the false-positive /
multiplicity problem; it is the latter this paper instruments.

\textbf{Contribution.} We build the machinery for that axis: a
falsification-first, decontaminated inference framework for testing
undeciphered-script decipherment claims \citep{braovic2024}. The harness
applies fail-closed validation to every positive and enforces four
concrete defences: a fabricated-language canary that supplies a
\textbf{generator-conditional negative-control distribution} (a
false-positive floor conditional on the fabricated-language generator ---
here a first-order sign-bigram Markov model; a broader control ensemble
spanning unigram, root-template, and real-language generators is
disclosed as future work, not built here); multiple-comparison deflation
for the number of signs, roots, and families a search was free to try; a
literature index that partitions published (\texttt{L\_indexed}) from
literature-unseen (\texttt{L\_not\_indexed}) signs so that reproducing a
known value cannot masquerade as discovery; and pre-registration with a
held-out gate whose verdict is computed mechanically, never by the model
that proposed the reading. The evaluation gap gets a concrete, if
power-limited, answer: the recurring libation \emph{formula} --- roughly
50 inscribed objects spread across peak sanctuaries and caves at more
than five findspots --- supports a natural leave-one-site-out
cross-validation graded from held-out data --- explicitly a
low-statistical-power test, reported as such rather than as decisive
(\S2.1; Appendix~A.3) \citep{braovic2024}.

\textbf{This is a methods paper, and Linear A is its worked null.} We do
not claim a decipherment. The framework is designed to withhold
validation when a proposed reading does not clear pre-registered negative
controls, multiplicity adjustment, contamination checks, and held-out
evaluation; when nothing survives, the calibrated null --- what the
corpus can and cannot support --- is itself the deliverable. The working
conclusion on current data is that the results are \textbf{consistent
with identifiability} being the more fundamental present constraint ---
not a demonstrated corpus-size threshold: there is no known related
language to map onto, the multiple-testing burden of candidate readings
goes uncorrected, and the sign inventory is mixed and partly lossy. We do
not claim additional data would be irrelevant --- a null can arise from
identifiability, low power, or corpus size alike. A lossy, non-injective
sign-to-sound channel need not admit a decipherment at any corpus size,
so we make no claim that ``there is enough text to pin a reading''; the
rigorous demonstration of \emph{that} constraint is the result the field
can use. The June-2026 AI-assisted ``Linear A cracked'' claim --- reading
the *301-anchored libation verb as an extinct Semitic form --- serves
throughout as a qualitative illustration of the failure modes the gate is
designed to detect, analysed from its public, archived presentation as a
claim, not a person, and neither reproducible nor held-out validated.
The artifact is open and archived (\S9.7).

\textbf{Roadmap.} Section~2 situates logos against the Braovi\'{c} et al.
taxonomy and the computational-decipherment lineage, including the
cognate-ID baseline. Section~3 specifies the harness --- canary,
deflation, literature index, and the pre-registration/held-out gate.
Section~4 develops the information-floor diagnostic and states its
identifiability limits. Section~5 works the June-2026 *301 claim as a
case study of the failure modes. Sections~6--8 report the pre-registered
nulls: morphology (\S6); metrology, phonology, and cross-script signals
(\S7); and contamination and sufficiency probes (\S8). Section~9 draws
the discussion, limitations, and conclusion.

\section{Related work}

\subsection{Position relative to the survey taxonomy}

The reference frame for computational work on the Bronze Age Aegean
scripts is \citet[\emph{Computational Linguistics}
50(2):725--779]{braovic2024}, the field's first survey focused
specifically on Aegean and Cypriot scripts; its only comparable
predecessor is \citet{ferrara2022}, and the broader
machine-learning-for-ancient-languages landscape is surveyed by
\citet[\emph{Computational Linguistics} 49(3)]{sommerschield2023}.
Braovi\'{c} et al.'s Section~4 enumerates fifteen ``major challenges'' ---
from unknown writing system, reading direction and language, through
small dataset, allographs and incomplete vocabulary, to unavailable
parallel data and hardware --- \textbf{none} of which concerns
overfitting, multiple testing, false-positive control, contamination, or
distinguishing a \emph{fitted} decipherment from a \emph{real} one. The
same gap runs through Ferrara and Tamburini's
\citeyearpar[\emph{Lingue e Linguaggio} XXI.2:239--259]{ferrara2022}
\S3.5, which reviews the cognate-matching lineage logos builds on ---
\citet{snyder2010} on Ugaritic$\leftrightarrow$Hebrew (the 29/30-sign
benchmark our canary calibrates against), \citet{bergkirkpatrick2011},
and \citet{luo2019, luo2021} --- and frames Linear A, via Gelb and
Whiting's typology \citep{gelb1975}, as a Type III script (unknown script
\emph{and} language) whose ``necessary validation'' is ``altogether
impossible,'' yet nowhere treats how a match is to be \emph{deflated} for
the multiplicity it was selected from. Two surveys, the same unfilled
axis. We do not score this as an oversight: their challenges are
properties of the object and its data, whereas the axis logos foregrounds
is a property of inference \emph{methodology} --- named-but-unfilled
rather than absent. logos supplies the machinery to police it (\S3).

Braovi\'{c} et al. do gesture at a related but separable difficulty:
evaluation. Because computational analyses of an undeciphered script have
``nothing to compare them against,'' they conclude (Sec.~6.3) that
``their evaluation remains an open problem.'' That is a distinct problem
from overfitting and multiplicity control, and logos addresses it
separately. Its held-out design is a concrete response --- the libation
formula recurs across a small set of peak-sanctuary and cave findspots,
permitting a leave-one-site-out cross-validation adjudicated mechanically
by \texttt{verdict.py} --- with the low-power caveat flagged at the
outset (\S1). We locate our components in the survey's own
method taxonomy (Sec.~7.1) and, where prior art exists, defer to it
rather than claim novelty.

\subsection{The current philological synthesis}

\citet{salgarella2025}, \emph{Writing in Bronze Age Crete: ``Minoan''
Linear A} (Cambridge Elements), by the author of the SigLA database, is
the most current authoritative one-author overview. It concludes that
Minoan is a language \textbf{isolate} --- it ``does not belong in any
known language family'' (the families used for comparison being
Indo-European, Semitic and Afro-Asiatic) --- typologically agglutinative
with a provisional Verb--Subject--Object order, and it records that ``no
bilingual text has so far been found'' in a corpus of roughly 2,500
inscriptions. Crucially, Salgarella ranks the etymological/comparative
method below combinatorial-contextual analysis and names
\textbf{digital/statistical analysis as ``the next desideratum in the
field.''} logos positions itself as exactly that desideratum --- but
under discipline, treating the isolate verdict as a central
\emph{identifiability} constraint rather than a target to be argued away.

\subsection{Prior computational attempts and the search trap}

The seminal automatic-decipherment result is \citet{luo2019}, which
recovers a character mapping from a lost script to a \emph{known related}
language via neural minimum-cost flow. It auto-deciphered
Ugaritic$\rightarrow$Hebrew ($+5.5\%$ over prior state of the art) and
translated 67.3\% of Linear B$\rightarrow$Greek cognates, but did not
address Linear A and would fail there by construction: with no known
cognate plaintext, the objective has nothing to optimise toward. This
cognate-language precondition is the empirical form of our
identifiability argument. The cognate-ID baseline logos actually runs is
\emph{not} a reimplementation of Luo's neural flow: it is a published
matcher, \citeauthor{tamburini2025}'s \citeyearpar{tamburini2025}
CSA\_OptMatcher (coupled simulated annealing), which reports the same
\citet{luo2019} accuracy metric. We adopt it as the non-neural baseline
against which any Linear A cognate claim must clear.

Against that disciplined baseline stand the in-domain instances of the
search trap: the dictionary-fishing programmes of \citet{eu2019} and
\citet{loh2020}, the latter explicitly a ``brute force attack'' on
dictionaries. These are the concrete Linear-A cases the multiplicity
correction is built to deflate. On the visual side, cross-script
comparison is \emph{not} our novelty: \citet{daggumati2019, daggumati2023}
already applied CNN+SVM cross-script similarity, and \citet{corazza2022}
(``Sign2Vec'') clustered signs unsupervised; our defensible contribution
is only the controlled image-versus-sequence \emph{asymmetry} under one
harness.

\subsection{The discipline lineage}

Two older strands supply the honesty machinery. \citet{packard1974},
\emph{Minoan Linear A}, constructed nine deliberately fictitious
decipherments as a null control --- and the genuine Ventris-value transfer
beat them only about 2:1 --- anticipating our permutation-null canary. The
statistical lineage is selective inference and multiple-comparison
correction \citep{bailey2014}, the correction for the number of
hypotheses $\times$ signs $\times$ roots $\times$ families tried that
turns a raw match-rate into a defensible one; its transfer from the
finance/backtest setting to script decipherment is specified in
Section~3. logos's empty cell is the intersection none of these occupy: a
falsification-first, decontaminated inference framework that couples a
published cognate-ID matcher \citep[reporting the Luo 2019
metric]{tamburini2025} with preregistration, multiplicity deflation and
mechanical held-out verdicts, released openly (\S9.7), with Linear A
worked as the honest null rather than a claimed crack.

\section{Methods: the discipline harness}

The harness is a comparison-layer pipeline that a reader can reimplement.
A candidate reading enters as a \emph{proposal} and passes through four
ordered stages --- pre-registration, deflation, decontamination, held-out
graduation --- under a cross-cutting persistence invariant: every decision
is computed from stored artifacts, never from a model's narration. Model-derived signals --- LLM theses, learned embeddings,
cognate-matcher assignments --- enter capped at confidence $\le 0.75$,
structurally unable to \emph{decide} a verdict; only
mechanically-verified held-out reads rank higher (\S1--2).

\subsection{Pre-registration}

Before any model is built or run, the predictions are frozen. The
morphology pre-registration fixes, for each hypothesis, the claimed affix,
its source, the falsifiable test, and the acceptance threshold, frozen
in-repository before each run so the model cannot be tuned to them after
the fact (the anchoring mechanics and the per-affix acceptance criteria
are given in Appendix~A.1). Outcomes are enumerated in advance --- CONFIRM, EXTEND,
REFUTE/NULL, or NO POWER --- so that a negative is a publishable result,
not a failed experiment. Any change is a new dated pre-registration,
never an edit.

\subsection{Deflation}

Internal fit on the derivation set is treated as no evidence. Every match
rate is deflated for everything tried: \texttt{n\_hypotheses} $\times$
\texttt{n\_signs} $\times$ \texttt{n\_roots}. The nulls are pipeline-level
--- a frequency-banded permutation null (after \citealp{packard1974}) and
the order-statistic bar of \S3.4 (implementation detail in Appendix~D).
Multiplicity is instrumented, not hand-passed, by a deduplicating
\texttt{SearchLog} (invariant 12; Appendix~D). The strong structural statistic,
\texttt{S\_morph}, earns credit only through a productivity gate --- an
affix must attach to $\ge$2 distinct residual stems, so a single formula
word copied across inscriptions scores zero (calibration in
Appendix~A.2, Table~\ref{tab:calib}).

\subsection{Decontamination}

A large open-weight model may have encountered longstanding published
readings by Gordon and Best; reproduction of indexed historical proposals
is therefore treated as potential shared-source contamination rather than
independent evidence. The June-2026 Di Mino claim is analysed separately
as a contemporary public case and is not assumed to be present in the
model's pretraining data --- it is indexed so that future claims and
models cannot reproduce it as discovery. Three mechanisms address this.

\textbf{The literature index} records published Linear A sign-value and
correspondence claims; any proposal matching the index is tagged
\texttt{literature\_match} and quarantined --- it may be tested but can
never count as discovery. The index partitions the sign inventory into
\texttt{L\_indexed} (referenced by $\ge$1 published proposal) and
\texttt{L\_not\_indexed} (untouched). Discovery claims may rest only on
\texttt{L\_not\_indexed} held-out success, since regurgitation can only
return what is in the literature. (The conservative seed, the adversarial
population, and the census are given in Appendix~C.5,
Table~\ref{tab:canary}.)

\textbf{The \texttt{L\_fake} canary} is a phonotactically-plausible
invented lexicon, calibrated to a real candidate language
(Appendix~C.5), generated after the model's release and never published
by us; exact-string collisions with real lexica are measured
(Appendix~C.5). Running the full pipeline
against it yields a \textbf{generator-conditional} false-positive floor:
any apparent lexical support it produces is spurious with respect to the
target historical-language hypothesis, and memorisation cannot
be invoked because there is nothing to memorise. A real candidate's match
must outperform this \emph{one calibrated class of negative controls} by
the corrected margin --- support means the candidate beat this
generator-conditional control family (here a first-order sign-bigram /
root-template generator), not that it defeated every possible false-language
process.

\textbf{LLM-ablation} compares the pipeline with and without the model's
proposals to isolate contamination (procedure in Appendix~C.1).

\subsection{The held-out graduation gate}

A reading graduates only mechanically, from held-out data, as a
conjunction of pre-registered clauses computed by \texttt{scripts/verdict.py}
--- never by the model that proposed the reading. The operative deflation
clause is an \textbf{order-statistic bar}: the observed held-out statistic
must exceed the \emph{expected maximum} of the multiplicity-null taken over
the number of trials the search was free to run (the E[max] over
\texttt{n\_trials} deflation), a \textbf{normal-order-statistic
approximation} in which \textbf{$n =$ \texttt{n\_trials} is the \emph{count of
candidates the search tried}, not an estimated number of independent
hypotheses}. The explicit formula and its Gaussian/independence assumptions
are given in Appendix~D; the error-control claim is carried not by the formula
but by the direct empirical calibration below. This is a direct
correction for selection across the hypothesis
space, and it fails closed on a degenerate null (Appendix~D). The framework yields
four primary evidential outcomes: \textbf{GRADUATE} (every clause holds and the
search multiplicity was instrumented), \textbf{REJECT} (a clause failed, or
the statistic sits below the corrected bar), \textbf{NULL\_PUBLISHED} (the
statistic is indistinguishable from the \texttt{L\_fake} null --- a
calibrated non-result, publishable as such), and \textbf{NO POWER} ---
the last assigned \emph{upstream}, by the pre-registered module outcome
when the experiment cannot adjudicate (\S3.1), before the gate sequence
runs; the gate itself additionally emits a fail-closed \textbf{INCOMPLETE}
status when search multiplicity is un-instrumented. These outcomes concern \emph{evidential
support, not truth}: because Linear A has no ground truth, REJECT / NULL /
NO POWER mean a reading is unsupported or insufficiently validated on
held-out data, \textbf{never that the underlying linguistic proposition is
false}. The verdict follows a fixed decision sequence
(\texttt{verdict.py}): all pre-registered clauses hold on an above-null
match $\rightarrow$ \textbf{GRADUATE}; else a match clearing the
substantive clauses but on an un-instrumented search $\rightarrow$
\textbf{INCOMPLETE}; else an observed statistic not exceeding
the null mean $\rightarrow$ \textbf{NULL\_PUBLISHED}; else $\rightarrow$
\textbf{REJECT}.
We calibrate the refusal rule directly: under the
tested best-of-100 random-map null --- an analyst trying 100 random
candidate maps and submitting the best --- 3 of 500 replicates graduated,
an observed false-graduation rate of 0.6\% (one-sided exact
Clopper--Pearson 95\% upper bound $\approx 1.54\%$), with 487/500 rejected
outright and 10 published as calibrated nulls (Appendix~D); this is an
empirical
error-control result conditional on that null-generating process, not a
universal guarantee. A second
pre-registered clause, \textbf{\texttt{generalizes\_to\_not\_indexed}},
requires the reading to recover literature-unseen (\texttt{L\_not\_indexed})
signs --- so reproducing already-published values cannot by itself
graduate a claim; the clause fails closed (its pre-committed thresholds
are in Appendix~D). Graduation is defined as validated \emph{novel}
discovery; independent held-out support for an already-published reading
is recorded as a distinct outcome and does not graduate. The remaining
pre-registered clauses (\texttt{L\_fake} separation, corrected margin,
order-statistic bar) are conjoined as specified in Appendix~D. These are
\emph{gate parameters} --- pre-registered
choices for this corpus, not universal constants. Validation is
leave-one-site-out, not ordinary \emph{k}-fold (Appendix~A.1). The gate is
AND-monotone --- clauses can only make it stricter --- so hardening never
lets a previously-failing reading graduate.

\subsection{Grade from persisted artifacts}

The invariant across every gate: verdicts are computed from persisted
rows, never from a model's account of its own outcome. The verdict writer
is grep-clean of any model call and is the sole writer of the verdicts
table.

\begin{table*}[t]
\centering
\begin{tabular}{@{}p{2.6cm}p{3.0cm}p{2.4cm}p{3.0cm}p{3.3cm}@{}}
\toprule
\textbf{Probe} & \textbf{Control / comparator} & \textbf{Real outcome} &
\textbf{Null / control floor} & \textbf{Verdict} \\
\midrule
Graduation gate (\S3.4) & --- & --- & best-of-100 random-map null &
3/500 false graduations (0.6\%; calibrated under the tested null) \\
Morphology (\S6) & Linear B 9/16 & Linear A 4/16 & \texttt{L\_fake}
bigram 6/16 & NO POWER \\
Segmentation (\S6) & --- & micro-F1 0.436 & random 0.389 & supported (gap
CI excludes 0) \\
Metrology (\S7.1) & planted 0.9 & 0.000 & permutation 0.113 & null
(agrees with the field) \\
Phonology (\S7.2) & planted (fires $\ge$ 2) & 0.080 / 0.140 & 0.065 /
0.201 & null (data-limited) \\
Image alignment (\S7.3) & font-swap survives & 0.410 & image chance
0.0135 & circular demonstration (capped) \\
Contamination (\S8.1) & model-free 0/40 & 14/40 (JA-NE) & 0/40 &
exploratory (public-exposure gradient) \\
\bottomrule
\end{tabular}
\caption{One harness, calibrated epistemic outcomes across probes.}
\label{tab:harness}
\end{table*}

The same falsification-first harness returns \emph{different} calibrated
verdicts across probes (Table~\ref{tab:harness}) --- the paper's
contribution in one view. Each number is detailed in the cited section.

\section{The information floor}

Before any decipherment method is graded, the corpus must be sized against
what it can identify. Two quantities matter: how much text exists, and how
many free parameters a reading spends against it. Confusing the first for
the second is the origin of the ``too little text'' folk claim --- and, we
argue, of its equally unearned inverse, the optimism that a big-enough
corpus makes a reading identifiable.

\textbf{The corpus.} The current authoritative synthesis
\citep{salgarella2025} gives an envelope of roughly 2,500 total Linear A
finds, of which about 1,400 are readable inscriptions, overwhelmingly
concentrated in the single LM IB destruction horizon (findplace and
dating detail in Appendix~C.4). That near-single horizon is why a
site-held-out split, not a chronological one, is the only realistic
partition. No bilingual text has ever been found
\citep[\S8]{salgarella2025}.

\textbf{Corpus sizing and the unicity computation (numeric detail in
Appendix~C.4).} The recurring sign-count confusion resolves into several
distinct sign-count denominators and two occurrence channels that must be
kept distinct (Table~\ref{tab:denoms}; Appendix~C.4). On the parsed
syllable-sign budget, Shannon's unicity distance sits everywhere far
below the corpus size (Table~\ref{tab:denoms};
\citealp[\S1--3]{salgarella2025}); we retain
this only as a toy-model \emph{precondition} diagnostic and explicitly
withdraw it as any sufficiency or refutation claim. Qualitatively, the
figure cannot settle identifiability either way (Appendix~C.4): unicity
distance presupposes a known plaintext-language oracle we do not have.
Corpus size, in short, cannot decide the question.

\textbf{What actually binds.} The results are consistent with
\emph{identifiability} being the more fundamental present constraint rather
than corpus size \citep[\S8]{salgarella2025}; we do not establish that more
data would be irrelevant. It has three parts: (i) the \textbf{absence of a
known cognate language} --- the map has a value-space but no anchored
target (Luo 2019's limit); (ii) the \textbf{multiple-testing / search
trap} --- that a unique consistent map may exist in principle neither
guarantees an algorithm finds it nor stops a cherry-picked false positive
across the parameters tried (invariant 8); and (iii) a \textbf{mixed
inventory} of syllabograms, logograms, and numerals that a real reading
must first separate. These, not the count of surviving signs, are what a
decipherment must overcome. Accordingly, when a candidate lexicon grows
without bound on a fixed occurrence budget --- free parameters exceeding
the floor --- the reading is underdetermined, and the framework must say
so.

\section{A qualitative audit of the *301 claim}

The framework of Section~3 is abstract until it meets a concrete claim. In
June 2026 a widely circulated ``Linear A cracked'' announcement
\citep{dimino2026} argued that Linear A records an extinct Semitic
language, reviving Gordon's hypothesis \citep{gordon1957, gordon1966}. Its
load-bearing anchor is a single reading: the recurring peak-sanctuary
libation formula word A-TA-I-*301-WA-JA, in which the Linear-A-only sign
*301 is \emph{itself} assigned a /na/-initial value, so the verb root
reads as West-Semitic N-W-Y ``to dwell'' (\emph{nawaya}) and is thereafter
matched to later Hebrew liturgy. We analyse this \textbf{as a claim, on its
public evidence} --- not the person --- from an immutable archived copy
carried in \artobj. It is a clean, self-contained instance of the failure
modes the gate is built to catch, and the three-part refutation below rests
wherever possible on the claim's \textbf{own published table} and on
primary sources \citep{davis2013, duhoux1992, thomas2020, salgarella2025,
steele2017} --- none advanced in logos's own voice, and the paper
introduces no new sign reading. We flag the meta-point once: this is a
literature-and-primary-source argument, not itself a logos held-out result
--- the graduation gate and \texttt{L\_fake} canary are validated as
instruments (unit tests; the known-answer Ugaritic$\leftrightarrow$Hebrew
surrogate) but have \textbf{not been run on this claim}. Scoring Gordon's
equations and the Di Mino lexicon against the \texttt{L\_fake} floor and
the order-statistic bar is directly runnable with the components of \S3 and
is the priority of the planned revision; that the refutation has not been
independently vetted by Aegean epigraphers is a disclosed limitation
(\S9.6), not a standard the work claims to meet.

\textbf{(a) The segmentation is not novel.} Root i-*301 plus affixation is
the reading of \citet{davis2013}; an independent segmentation by
\citet{thomas2020} reaches the same core and is tested head-to-head against
it (the current synthesis, \citealp{salgarella2025}, does not itself adopt
Thomas). On the Davis and Thomas segmentations, the affix-stacking is
evidenced by multi-attestation forms --- ta-na-i-*301-u-ti-nu (IO Za 6) and
a-ta-i-*301-de-ka (ZA Zb 3) --- that vary the material around a stable
i-*301 core. (These sigla are confirmed against Davis 2013: \texttt{TL Za
1}, \texttt{IO Za 6}, and \texttt{ZA Zb 3}.) The claim re-segments a
structure already established in the literature. Under the decontamination
layer this matters procedurally, not just historically: a correspondence
that reproduces a published segmentation is treated as potential
shared-source contamination,
tagged \texttt{literature\_match} and barred from counting as discovery
(\artref, \S C.1).

\textbf{(b) The gloss is contested and weaker.} The mainstream value of the
verb is ``give/dedicate'' --- Duhoux's \citeyearpar{duhoux1992} gloss,
adopted by \citet{davis2013} as a \emph{structural} reading that assigns
*301 \textbf{no phonetic value at all}, and endorsed by the current
synthesis \citep[\S7.2]{salgarella2025}. ``Dwell'' is a semantically
distant, unmotivated substitution. The point is not that ``give'' is
certainly right, but that a \emph{structural} reading needing no
sound-value for *301 already accounts for the data the phonetic claim is
invoked to explain.

\textbf{(c) The primary descriptions of this form's morphology do not
support the Semitic label.} The claim's own table segments the verb as
prefix-stacking and agglutinative --- [prefix A ``I''] [stem-morpheme TA]
[stem-vowel I] [root *301]. We deliberately do \textbf{not} rest on an
\emph{a-priori} typological contradiction here, and we retract the earlier
draft's claim that a prefix-stacked verb is ``internally incoherent'' with
a Semitic reading: Semitic verbal morphology includes productive prefix
conjugations, so a prefix on the verb is not by itself un-Semitic. The
point is empirical, not typological, and it comes from the primary
literature. \citet[p.~38 n.~13]{davis2013} characterises i-*301's
morphology as ``neither IE nor Afroasiatic'' --- and Semitic \emph{is}
Afroasiatic, so a decade earlier the primary description already found
\textbf{no Afroasiatic morphological signature} in this very form. (We rest
on what the footnote states --- the absence of an Afroasiatic signature ---
not on a stronger ``ruled out'' than a hedged footnote can bear.)
Independently, \citet[\S8]{salgarella2025} concludes Minoan is
agglutinative (after Duhoux 1978) and a language isolate. The claim's own
segmentation is thus consistent with the isolate typology, while the
specific Semitic assignment finds no morphological support in the primary
descriptions of the form.

\textbf{The free-parameter receipt.} The decisive point for identifiability
is that *301 has no cross-script anchor: in the claim's own table its
AB-number column is ``---'' and its Linear B correspondence ``(none)''.
Even the standard practice of reading Linear A homomorphs with their Linear
B values is a backward projection to be handled with care
\citep[already cautioned by Duhoux 1989, apud Davis 2013, p.~38
n.~11]{steele2017} --- and an A-only sign such as *301 has no homomorph to
project from at all, so nothing cross-script constrains its value. The
single sign on which the entire decipherment pivots is thus a \textbf{free
parameter} --- its value was assigned to fit the desired root, not
constrained by any independent evidence. This is precisely the case the
minimum-description-length / identifiability \emph{diagnostic} flags: a
value chosen to make one word read out carries no description-length saving
that survives the multiplicity of alternatives it was selected from (cf.
Section~3; the diagnostic is reported next to every claim, and the
operative graduation clause --- the order-statistic bar over exactly that
multiplicity --- is one a hand-fitted value cannot clear). For this reading
the obstacle is identifiability rather than corpus size --- no known cognate
language to constrain the assignment, uncorrected multiplicity, and a sign
whose value is chosen rather than anchored (a free parameter is unanchored
at any corpus size).

Read against Section~3, the claim instantiates two distinct failure modes
at once. It is \textbf{conclusion-driven model specification} --- the
table's \emph{structure} says agglutinative and prefixing, but the
\emph{label} says Semitic, and the label wins because the desired
conclusion (a Semitic prayer) is a fixed prior rather than an emergent
reading. And it is a \textbf{free-parameter / identifiability} failure ---
an unanchored sign-value tuned to the target, on a corpus with no known
cognate language to constrain it. logos indexes the reading (*301 = /na/)
precisely so a model cannot regurgitate it as discovery, and quarantines
rather than dismisses it. The lesson generalises: a match rate on a
hand-picked word, absent a committed prediction and multiplicity
accounting, is not evidence. The *301 claim is a qualitative illustration
of the failure modes the gate is designed to detect.

\section{Results I: morphology induction}

We test whether Linear A's pre-registered affixes and word divisions
reflect genuine morphology. The apparent affixation does not clear the
operative \texttt{L\_fake} bigram floor, so the verdict is NO POWER and
\texttt{has\_validated\_morphology} = False. The identical test
\textbf{validates on Mycenaean Greek (Linear B)}, whose longer words carry
uncontested inflection: it confirms affixation at rate 0.562 --- above the
same bigram floor ($\sim$0.30) --- with \texttt{has\_morphology\_power} =
True, so the control rules out a universally insensitive detector and
supports the reading that Linear A's shorter forms contribute materially to
its lack of power (Table~\ref{tab:morphrates}; Appendix~A.1). One held-out positive stands and survives
its own uncertainty: a DP-unigram segmenter recovers the scribe's own word
divisions at micro-F1 0.436 versus a random-boundary baseline of 0.389, a
gap whose site-clustered 95\% bootstrap CI [0.021, 0.099] excludes zero
(P(gap $>$ 0) = 0.998; Appendix~A.1). This
separates real recoverable structure from fitted morphology --- the paper's
core discipline.

\section{Results II: metrology, phonology, and the cross-script asymmetry}

Three further pre-registered probes exercise the harness against distinct
field claims. Each returns a calibrated null or a truth-capped positive;
together they map where the Linear A corpus does and does not carry
recoverable signal, and they are consistent with identifiability rather
than effort being the more fundamental present constraint.

\subsection{Metrology: a parse-coverage null that agrees with the field}

Direction D poses the accounting tablets as a constraint-optimisation
over the fraction-sign values (Appendix~B.1; Table~\ref{tab:metro}).
Held-out fraction balance
does not beat the shuffle null
($p = 1.0$), so the verdict is NULL; the only value the method determines,
the fraction sign J = 1/2, is the long-standing field consensus
\citep{bennett1950}, not a logos discovery. This is a null the field itself
reached \citep{corazza2021} --- the harness crediting a surviving value to
its originator rather than re-branding it.

\subsection{Distributional phonology: a data-limited null bounded by a
power curve}

Does a sign's distributional context (which signs sit beside it) carry
information about its phonetic value? Tested on the known AB signs, neither
the consonant class nor the vowel class clears chance (permutation
$p = 0.40$ / $0.86$). The verdict is a data-limited NULL: the sound path is
not recoverable from Linear A alone, and no sound value is imputed for any
sign (Table~\ref{tab:phon}; Appendix~B.2).

\subsection{The cross-script asymmetry: image alignment recovers shared
anchors, but the recovery is circular by construction}

Do image embeddings recover the cross-script A$\leftrightarrow$B
correspondence that sequence and cognate co-occurrence cannot? Classical
shape descriptors reach direct-NN \textbf{0.410} on held-out
A$\leftrightarrow$B anchors while sequence and cognate co-occurrence stay
at chance. But the recovery is circular by construction: the
A$\leftrightarrow$B anchor values were themselves assigned by
twentieth-century epigraphers on sign-shape similarity, so this is a
shape-genealogy engineering demonstration, truth-layer-capped $\le$
\textbf{0.75}, not a positive decipherment claim (\S9.6;
Table~\ref{tab:xscript}).

\section{Results III: contamination probes}

One failure mode is invisible to the false-positive / multiplicity axis of
\S7: a decipherment that \emph{looks} novel but silently reproduces
published readings a language model has memorised. This section reports the
two decontamination probes that target it --- both run as method
demonstrations, returning completed (if deliberately bounded) findings.
Throughout, where a clean number was not obtainable we report the
partial-null and the confound rather than a headline.

\subsection{LLM-ablation: does a large open-weight model regurgitate the
literature?}

In the LLM-ablation, \texttt{qwen2.5:72b} reproduces famous
published readings --- Gordon's JA-NE = ``wine'' in 14/40 seeds --- that
the model-free edit-distance baseline never reaches (0/40), yet never
reproduces obscure vessel equations (0/40; Table~\ref{tab:repro}).
Verdict: a reproduction
gradient consistent with differential public exposure --- though the
experiment does not isolate memorisation from correlated item-level
differences --- the memorisation confound the discipline layer must gate
out.

\subsection{\texttt{L\_not\_indexed}: indexed versus not-indexed
abstention}

Does the model's behaviour on \emph{published} signs generalise to signs
unseen in the indexed literature? \texttt{qwen2.5:72b} valued the
memorisable known signs in $\sim$20/40 seeds but abstained $\sim$97\% of
the time on the literature-unseen *-series signs ($\sim$1/40;
Table~\ref{tab:abstain}) --- an
abstention asymmetry consistent with differential public exposure (the
\S8.1 caveat applies). With no ground truth on the not-indexed signs,
\textbf{consistency is not correctness}: it bounds how the model behaves,
never whether it is right.

\section*{Registered extension: a known-answer sufficiency curve}
\addcontentsline{toc}{section}{Registered extension: a known-answer
sufficiency curve}

Is a Linear-A-scale corpus large enough to recover a decipherment if a
known cognate existed? Because every benchmark is handed a real cognate
signal Linear A lacks, any recovery is an \textbf{optimistic benchmark,
conditional on a known cognate relationship}: at-floor at Linear-A scale
$\Rightarrow$ information-insufficient even given a cognate; above-chance
$\Rightarrow$ identifiability (cognate-absence, uncorrected multiplicity)
rather than corpus size would be implicated; neither branch is a
decipherability claim. The completed 2{,}000-step curve (near-converged:
Luvian$\to$Hittite 0.462 vs.\ published 0.475, Phoenician$\to$Ugaritic
0.760 vs.\ 0.805; every recovery a \emph{lower bound}) splits the
branches at the sweep's Linear-A-scale locator (600--700 distinct
word-forms, syllabic normalization): Linear~B$\to$Greek, the primary
analog, stays at its chance floor ($0.044\pm0.023$, $n{=}276$; floor
$\approx0.004$), while Cypriot$\to$Greek sits marginally above
($0.066$, $n{=}346$). Under-convergence keeps the at-floor branch
non-definitive; full per-size and per-cell artifacts are deposited
(Appendix~C.3).

\section{Discussion, limitations, and conclusion}

\subsection{The discipline is the contribution}

The deliverable of this work is not a reading of Linear A. It is a
\emph{harness} --- a literature index with an indexed versus
literature-unseen partition, a fabricated-language control corpus,
model-ablation contamination checks, multiple-comparison correction, and
pre-registration, all graded mechanically from persisted artifacts ---
together with a body of pre-registered, calibrated nulls, each with its
borderline case named. The stratified morphology re-run is the template: at
the pre-registered $\alpha = 0.01$, no pre-registered affix is cross-stratum
stable above the bigram floor, robust to the genre-bucketing choice
(Tables~\ref{tab:strata}--\ref{tab:sens}; Appendix~A.1). A null that
names its threshold, its sensitivity, and its
near-miss is a stronger referee object than a flat negative.

The axis this harness targets is the \textbf{false-positive / multiplicity
axis} --- the mechanism by which a determined analyst ``translates'' almost
anything on a small corpus --- a property of inference methodology, not of
the object or its data (\S1--2).

\subsection{Two named failure modes for the methods literature}

\textbf{(i) Narrative capture.} Structure says X; the label says the
desired Y; the label wins. In the June-2026 *301 claim the published table
segments an agglutinative, prefix-stacking verb while the assigned family
label is Semitic, and the desired output operates as a fixed prior, so
the conclusion conforms to the prior rather than emerging from
constrained readings (\S5). This
is distinct from statistical overfitting and must be named separately.
Notably, the mismatch is empirical --- between the label and the primary
descriptions of the form (\S5) --- not an a-priori typological
impossibility.

\textbf{(ii) Free parameters exceeding the information floor.} A lexicon
reported as ``408 terms and counting'' on a fixed corpus has more degrees
of freedom than the data can identify; a decipherment whose lexicon grows
weekly is underdetermined by construction. The minimum-description-length /
unicity \emph{diagnostic} (\texttt{k $\le$ U\_floor}) is built to flag
exactly this (it is reported next to every claim rather than enforced as a
graduation clause; \S4): the single sign the whole
reading pivots on is the textbook free-parameter case (\S5), its value
assigned to fit the desired root; the refutation rests on the claim's own
published table, archived in \artobj.

\subsection{The post-hoc evaluation sequence}

These failure modes trace to sequencing. The *301 claim publicly commits to
a post-hoc evaluation sequence: declare ``officially deciphered,'' then vet
later. logos's documented inverse is \textbf{pre-register $\rightarrow$
held-out $\rightarrow$ locked artifact} (\S3.1, \S3.4; the held-out
libation-formula corpus and its low-power caveat are in Appendix~A.3).

\subsection{Identifiability appears the more fundamental present constraint}

We take care not to write a defeatist framing, and equally not to
over-claim from corpus size: the toy-model unicity distance
($U \approx 204$--$415$ signs; Table~\ref{tab:denoms}, \S4) is a
\emph{precondition} diagnostic
only, explicitly withdrawn as any sufficiency or refutation claim ---
``$U \ll N$'' does \textbf{not} license the statement that there is enough
text to pin a decipherment (\S4; Appendix~C.4). The
results are consistent with \emph{identifiability} being the more
fundamental present constraint --- the absence of a known cognate language
to map to, an uncorrected multiple-testing search burden, and a mixed sign
inventory --- rather than a demonstrated corpus-size threshold; a null can
arise from identifiability, low power, or size alike, and we do not
establish that more data would be irrelevant.
\citet[\S8]{salgarella2025}
reaches the field-level version of this by a different route (\S2.2).

\subsection{Study-wide error budget}

Multiplicity discipline operates at two levels, and we are explicit about
both. Within a single probe, matches are deflated for everything
searched (\S3.2--3.4). Across the study, the unit of multiplicity is the
\textbf{per-probe pre-registration}: the \textbf{five executed probes}
(morphology, metrology, phonology, cross-script image alignment, and
contamination) each carry their own pre-committed $\alpha$, and every probe
is reported regardless of outcome. \textbf{Error control is defined within
each pre-registered probe; we claim no global family-wise guarantee across
probes.} Nothing in the battery is read as a decipherment.

\subsection{Limitations}

The nulls are honestly bounded, not decisive. A dependence-adjusted
trial-count correction of the morphology z-test is deferred; the deferral
would only strengthen the null (Appendix~A.1). No held-out
\emph{decipherment} has
been achieved, and the study reports \textbf{no positive decipherment
claim} --- the deliverable is calibrated absence, not a reading. Nothing
here claims a rejected reading is \emph{false} (\S3.4): every REJECT /
NULL in this study bounds \emph{evidential support}
on held-out data, not truth. The one probe that recovers an above-chance
number, cross-script sign-image alignment, is a font-robust but circular
\emph{engineering demonstration} (a font-swap control refutes the
typeface-artifact worry; \S7.3, Appendix~B.3), truth-capped at
$\le 0.75$, never a reading. No
phonetic claim is made. The philological points in \S5 rest on the
published literature cited there rather than on independent
Aegean-epigraphic
review, which we disclose as a limitation; because those points are cited
rather than original, the argument does not depend on any new sign reading
of our own. Chronological cross-validation is largely foreclosed by the
single-horizon corpus (\S4; Appendix~A.1). And the corpus itself
--- small, partly lossy, with no bilingual yet found
\citep[\S8]{salgarella2025} --- sets the ceiling.

\subsection{Conclusion and reproducibility}

The framework withholds validation when a proposed reading does not clear
the pre-registered gates (\S1, \S3); the contribution is the
framework and the calibrated null itself --- the instrument and the honest
limits of what the corpus can support --- not a reading, and the calibrated
null is a safeguard against unsupported claims.
\ifanon
The code, pre-registrations, and all findings are open under the MIT
licence and available as an anonymized artifact (code and pre-registrations
released on acceptance).
\else
The code, pre-registrations, and all findings are open under the MIT
licence and archived on Zenodo (concept DOI 10.5281/zenodo.21087572;
repository \url{https://github.com/papadopouloskyriakos/logos}).
\fi
The licensed epigraphic corpora are not redistributed; each source's
licence and retrieval path is documented (data-provenance; Software \&
data note), so the harness is reproducible against a
locally obtained corpus. This work is independent and not affiliated with,
or endorsed by, any of the cited projects or their authors.

\bibliographystyle{acl_natbib}
\nocite{*}
\bibliography{refs}

\appendix

\section{Morphology (replication detail)}

\paragraph{A.1 Morphology: rates, floors, word-length, stratification,
sensitivity.} The probe is pre-registered, unsupervised morphology and
segmentation induction over the full Linear A corpus
(Table~\ref{tab:denoms}). Target claims are the recurring affixes and
positional grammar of the Davis/Thomas/Duhoux--Val\'{e}rio tradition:
prefixes and suffixes (\texttt{i-}, \texttt{a-}, \texttt{-te/-ti}, and
Davis's \texttt{i-*301} verb stem). Pre-specification anchoring (\S3.1):
pre-specifications were frozen in-repository before each run (commit
history); the external immutable deposit
\ifanon
([artifact DOI withheld for review])
\else
(Zenodo concept DOI 10.5281/zenodo.21087572)
\fi
certifies the current state of the registrations, not temporal priority.
Each pre-registered affix
(drawn from \citet{thomas2020} and Duhoux/Val\'{e}rio) must recur on
$\ge$2 distinct stems across independent inscriptions, above a
within-form permutation null, and survive multiplicity correction at a
target of \emph{p} $< 0.01$ corrected. Every candidate was graded against two
null floors: a within-word sign-order permutation (shuffle) floor, and an
\texttt{L\_fake} floor built from a first-order sign-bigram (Markov) corpus
with zero lexical overlap with the real inscriptions.

\emph{Pooled run.} Of the 16 pre-registered affixes, 4 beat the within-form
permutation null and the deflated multiplicity bar (\texttt{i-*301},
\texttt{a-}, \texttt{i-}, \texttt{-te/-ti}), clearing the shuffle floor
decisively (rates and Wilson intervals in Table~\ref{tab:morphrates}).
The apparent standout \texttt{i-*301} sits far above the permutation null
--- but that figure is computed against the weaker within-form permutation
null only, and it does not survive the operative \texttt{L\_fake} bigram
floor, under which the whole set fails: the fabricated Markov corpus
confirms the same affixes at a \emph{higher} rate than the real corpus
(Table~\ref{tab:morphrates}). The outcome is classified NO POWER, not
evidence against morphology; \texttt{has\_morphology\_power} is False and
the module reports the null (the Linear-B positive control that validates
the detector is reported below).

\emph{Word-length distribution.} Linear A words are short
(Table~\ref{tab:morphrates}), and on such words ``morphology'' and
``bigram order'' are statistically indistinguishable. This reconciles
with Fuls's 3.3-sign figure once the denominator is matched:
\citet[57, \S6.5]{fuls2015} computes the mean over a word \emph{list} of
``554 complete sign sequences, where duplicates are reduced to one''
(GORILA; \citealp{godart1985}) --- distinct word \emph{types}, not
tokens; restricting our corpus to distinct multi-sign types likewise
recovers 3.07. Fuls reads the 3.3-sign length as evidence that Minoan is
\emph{polysynthetic}; our direct test finds the apparent affixation
bigram-reproducible with no power, and we report the null. The direct
positive control --- the identical bigram-floor test on Mycenaean Greek
(Linear B), whose longer words carry uncontested inflectional morphology
--- was RUN on the D\=AMOS corpus (parsed by the same
\texttt{\_damos\_wordforms} extractor as \S7.3) against a pre-registered
16-affix Mycenaean panel \citep{ventris1953}. It fires
(Table~\ref{tab:morphrates}), carried by the securely-grammatical
suffixes -jo, -o-jo (gen sg -oio), -de (allative), -qe (enclitic
``and''), -o-i (dat pl), so \texttt{has\_morphology\_power = True}. The
identical test that returns NO POWER on Linear A's short words recovers
Mycenaean morphology on longer ones (a sensitivity check: an ad-hoc
parser yielding mean 2.19 signs/word did \emph{not} fire --- power tracks
word length exactly as the short-word argument predicts), which rules out
a universally insensitive detector and supports the interpretation that
Linear A's shorter forms contribute materially to its lack of power
(\texttt{scripts/\allowbreak comparison/\allowbreak linb\_morphology\_control.py}).

\emph{Segmentation positive.} Under leave-one-site-out cross-validation (52
sites; not k-fold, given formulaic dependence), a DP-unigram
(Goldwater-style) segmenter recovers the scribe's own word divisions above
a random-boundary baseline; a site-clustered bootstrap (resample the 52
sites with replacement, 4,000 draws) puts the real-minus-random gap above
zero, and the paired gap, not the wide marginal per-segmenter CIs, is the
operative test (Table~\ref{tab:morphrates}). Morfessor over-segments the
short syllabic corpus and is reported as not usable rather than hidden
(Table~\ref{tab:morphrates}).

\emph{Stratified re-run.} The strata reconcile exactly to the
1,341-inscription corpus total (Table~\ref{tab:strata}). The
99-inscription ``libation'' \emph{genre} stratum is not the same object
as the libation-\emph{formula} corpus (\S A.3). Following
\citet{salgarellajudson2024}, chronological cross-validation was explicitly
declined (988/1,341 $\approx$ 74\% single-horizon LM IB), and site --- not
scribal hand, which is not individuated for Linear A
\citep{salgarella2019} --- is the independence unit, gated by presence at
$\ge$2 distinct sites. The result is \texttt{has\_validated\_morphology} =
False (\S9.1). Stratification also showed
that even \texttt{i-*301} is bigram-reproducible in the repetitive libation
register, where the Markov corpus manufactures the apparent verb
morphology; Davis's slot grammar is not refuted, merely not separable from
sign bigrams here. Where formula-bearing sites drive a signal, power is
intrinsically low: a leave-one-site-out CV over so few formula findspots
(an effective $\sim$5-fold site partition) cannot resolve a modest effect.

\emph{Sensitivity, disclosed.} The honest 2$\times$2 over genre bucketing
$\times$ $\alpha$ threshold is Table~\ref{tab:sens}; the doubly-relaxed
corner's \texttt{i-} (locative) and \texttt{-te/-ti} (ablative) are
precisely the Duhoux/Val\'{e}rio H3 nominal affixes that
\citet{salgarella2025} \S8 independently endorses --- the named borderline
candidates, fragile to \emph{both} the $\alpha$ threshold and a defensible
genre choice (the $\alpha = 0.01$ null is itself robust to the bucketing).
The present analytical approximation uses attempted-candidate
count rather than a dependence-adjusted effective trial count --- a
recorded, deferred limitation whose correction would only strengthen the
null.

\paragraph{A.2 \texttt{S\_morph} detector calibration (\S3.2).} After the productivity-gate fix (\S3.2 --- an affix must attach to
$\ge$2 distinct residual stems), \texttt{S\_morph}'s measured false-positive
rate on within-form-shuffled corpora is consistent with the nominal
one-sided $z \ge 2$ rate within sampling error, and its planted-affix
power is a saturated point, not a full sensitivity curve
(Table~\ref{tab:calib}; contrast the phonology power curve, \S7). This is
an \textbf{indicative calibration check, not the operative gate}: 6/200
is too coarse to certify an $\alpha = 0.01$ threshold on its own, and the
gate applies the order-statistic bar over the far larger permutation and
candidate counts the deflation machinery tracks (\S3.4).

\paragraph{A.3 Held-out design: the libation-formula corpus (\S9.3).} A
leave-one-site-out split over so few formula sites (\S1) has
low statistical power --- and site-level CV is in any case necessary, not
sufficient, since within-site scribal-hand correlation must also be controlled
\citep{salgarellajudson2024}. This formula corpus should not be conflated with
the 99 stone-vessel inscriptions of the ``libation'' genre stratum (\S6): they
are different objects.

\section{Metrology, phonology, and cross-script (methods and controls)}

\paragraph{B.1 Metrology.} Direction D parses HT (LM I) documents into
balance units \texttt{[recipient][commodity logogram][integer][fraction
signs]} and solves the fraction-sign values by RANSAC over exact-\texttt{Fraction}
Gaussian elimination (robust to the unbalanced-tablet outliers that defeat
least-squares), requiring line items to balance the preserved KU-RO totals;
validation is grouped k-fold by document (no tablet straddles train and
test) against a permutation null. Results are in Table~\ref{tab:metro}: of
$\sim$7 fraction-bearing constraints only one is satisfied by the solved
system, the sole determinable value is J = 1/2, held-out fraction balance
is not separated from the permutation null, and the synthetic positive
control (planted values, all recovered) separates. The null traces to sparse
automated parse coverage (multi-commodity blocks, illegible number glyphs,
KI-RO deficit sub-lists, editorial restorations); \citet[p.~2]{corazza2021}
abandoned the balance-to-totals method for the same reason. The surviving J
= 1/2 (\S7.1) gains Schrijver's
\citeyearpar{schrijver2014} three further tablets (HT Zd 156, PE 1, HT
123a): at least four corroborating documents.

\paragraph{B.2 Distributional phonology.} Track 2 asks whether a sign's
distributional context --- which signs sit beside it --- carries
information about its phonetic value, measured on the known \texttt{AB}
signs whose Linear B sound values can be borrowed. Features are
PPMI-weighted, L2-normalised left/right co-occurrence counts; labels are
the Linear-B-transferred (C, V) parsed from the GORILA name, so features
never see the labels. The test is leave-one-out 1-NN against a
label-permutation null, \v{S}id\'{a}k-deflated over the two tests; 50 of
$\sim$55 AB signs met the $\ge$3-occurrence threshold.

Per-test results are in Table~\ref{tab:phon}. The honest verdict is
\emph{data-limited},
not \emph{no bridge exists}, because the positive control disciplines the
reading: the positive-control power curve over planted strengths shows
the test firing at strength $\ge$ 2 but not at 1. The test
works but is not infinitely sensitive at n = 50 / 13 classes, so we cannot
distinguish ``no distribution$\rightarrow$phonetics bridge'' from ``too few
signs to detect a weak one'' (\S7.2, \S B.3).

\paragraph{B.3 Cross-script image alignment.} The palaeographic probe
renders 88 Linear-A and 74 Linear-B glyphs and tests held-out recovery of
shared A$\leftrightarrow$B anchors. Classical descriptors (HOG + Hu moments
+ shape-context) sit far above image chance, and a from-scratch I-JEPA is
excluded from the claims as over-parameterized for $\sim$90 glyphs
(results and controls in Table~\ref{tab:xscript}). Two controls decide how
to read the image-vs-sequence asymmetry (\S7.3), and together they
\emph{downgrade} it from a positive to a demonstration.

\emph{First, a font-swap control.} Rendering both scripts in one typeface
(Aegean, George Douros) risks making the alignment a single-designer
house-style artifact rather than a fact about the signs. We re-render Linear
A in Noto Sans Linear A and Linear B in Noto Sans Linear B --- two
independent type designers, each font covering only its own Unicode block
(Noto-A carries zero Linear B glyphs and vice-versa), so no single hand
draws both scripts. The recovery survives and is if anything stronger
(Table~\ref{tab:xscript}; both legs $\gg$ chance): the font-artifact
hypothesis is \emph{refuted}, and ``distances reflect shape, not font'' is
now demonstrated rather than asserted
(\texttt{scripts/\allowbreak palaeo/\allowbreak font\_control.py}).

\emph{Second, and decisively, the circularity the font control cannot
remove.} Linear B was historically adapted from Linear A, and the
A$\leftrightarrow$B transcription \emph{values} that define the anchor set
were assigned by twentieth-century epigraphers largely on the basis of
sign-shape similarity, so the recovery re-derives a shape-defined
correspondence (\S7.3): it confirms a known palaeographic fact (borrowed
signs look alike) and validates a font-robust descriptor pipeline, but it
reads no Minoan. It remains truth-layer-capped in any case --- the
\citet{ferrara2021} ``Archanes Formula'' hard negative (\texttt{CH 095}
resembles \texttt{AB 60}/\emph{ra} but is not its phonetic counterpart)
shows a plausible shape match can be phonetically wrong, so a shape match is
a $\le 0.75$ signal (invariant 5), never a reading. The honest asymmetry is
thus: the image path re-derives a shape-defined mapping (expected, circular,
capped), while the sequence and cognate paths \textbf{remain at chance even
after ingesting the full D\=AMOS Mycenaean corpus} --- 13,562 wordforms,
$\approx$15$\times$ the earlier cognate file, 56 anchors
(Table~\ref{tab:xscript}) --- while the \emph{positive control
recovers}, so the harness works and the null is a real null: the
A$\leftrightarrow$B phonetic correspondence is simply not carried by
cross-script sign co-occurrence, even at full Linear-B scale
(\S B.2). And the A-only signs on which
positive-decipherment claims would pivot have no cross-script
anchor at all (\S5), so images cannot impute them even in principle.

\section{Contamination and sufficiency (detail)}

\paragraph{C.1 LLM-ablation reproduction gradient.} The ablation
(\S3.3) is the remaining gate on the contamination number, matching a
model's cognate glosses against the indexed lexical claims. We ran it
against \texttt{qwen2.5:72b} --- a large open-weight model, a far
stronger memoriser than a local 12B --- on a rented H100, with an
NW-Semitic candidate and a Hebrew skeleton lexicon for the model-free arm
(scale and quantities in Table~\ref{tab:canary}) (\artref). The
first-pass contamination rate is an \textbf{artifact, reported as such}:
the mechanical arm never shares a literature-matching correspondence with
the LLM (their value vocabularies barely overlap), so the rate is forced
to 1.0 by construction and is not reported as a contamination result
(Table~\ref{tab:canary}; \artref).

\emph{Generation settings (replication).} Checkpoint \texttt{qwen2.5:72b}
as served by Ollama (Ollama's default \texttt{q4\_K\_M} quantization;
Ollama version not recorded), on an ephemeral vast.ai H100 SXM 80\,GB
instance (torn down after the run). One \texttt{/api/generate} call per
seed, stream off, timeout 180\,s; temperature 0.8; sampling seed = the
replicate seed; top-p, top-k, maximum output length, and stop conditions
left at server defaults (not recorded). Each
seed's 40 held-out forms are drawn deterministically
(\texttt{numpy.default\_rng(seed)}) and rendered
as canonical sign names in a fixed JSON-request prompt
(\texttt{\_PROMPT\_TEMPLATE} in
\texttt{scripts/\allowbreak comparison/\allowbreak ablation.py}). Parsing:
the JSON \texttt{partial\_map} is extracted; keys not naming a sign shown
to the model are dropped (the proposer cannot invent signs); a dead or
garbled call contributes an empty proposal without crashing the batch
(fail-closed). The gate-2 rerun
prepends the 7 indexed probe forms to every seed. Per-call logs:
\texttt{runtime/ablation/llm\_calls.jsonl}; harness commits
\texttt{0582966} (ablation) and \texttt{9cd6d04} (gate-2).

The gate-2 rerun, reconciled to a common consonant-skeleton vocabulary,
yields the honest deliverable: the per-item reproduction gradient over 40
seeds (Table~\ref{tab:repro}; \artref). An
external red-team correction split \texttt{KU-RO} between the
Semitic-root route (\emph{kull}) and the general-knowledge meaning
``total'' (confirmed by tablet arithmetic), so \texttt{KU-RO} = ``total''
is \textbf{excluded} as a witness and \texttt{KU-RO} = \emph{kull} is
downweighted (Table~\ref{tab:repro}; \artref). The surviving gradient ---
famous readings reproduced, obscure ones not --- is what the ablation
detects (\S8.1). The gate-2 rate itself remains confounded by the baseline's
representational ceiling and is not reported as a contamination measure
(Table~\ref{tab:canary}).

\paragraph{C.2 \texttt{L\_not\_indexed} abstention asymmetry.} The
\texttt{L\_not\_indexed} probe shows a sign in real context (\S8.2); the
not-indexed signs are absent from the indexed
sources, hence not memorisable from them (never ``impossible to have seen''
in any absolute sense) (\artref). Across 40 seeds, \texttt{qwen2.5:72b}
abstains $\sim$97\% of
the time on the untouched signs while confidently valuing the memorisable
ones (Table~\ref{tab:abstain}). The apparent known$-$not-indexed
consistency gap is coverage-confounded and was flagged, not asserted: the
matched-\emph{n} robustness check returned \texttt{no\_power}, so the
rigorous comparison could not run (Table~\ref{tab:abstain}).
The interpretable signal is the abstention asymmetry itself, read under
the standing rule of \S8.2 (\artref).

\paragraph{C.3 Information-sufficiency CSA sweep (design and deposit).}
The sweep downsamples known-answer benchmarks --- Linear B /
Mycenaean Greek as the primary syllabic analog, sourced through D\=AMOS,
and further syllabic corpora --- and measures held-out cognate-ID
recovery against per-size permutation floors, using the published
cognate-ID matcher of \S2.3 \citep{tamburini2025}, validated in the
artifact against its published Table-3 values (\artref). The reading
rules and the reported curve are in the body (Registered extension); the
full per-size report and per-cell artifacts are deposited in the
repository (\texttt{results/csa/}).
The phonological / sound path is no longer data-deferred: even the full
ingested D\=AMOS corpus leaves the cross-script distributional alignment
at chance (\S7.3, \S B.3).

\paragraph{C.4 Information floor: corpus denominators and the unicity
computation (\S4).} The corpus envelope, the four
sign-count denominators, the two occurrence channels, and the measured
unicity range are tabulated in Table~\ref{tab:denoms}
\citep[Schoep 2002;][\S1--3, \S6]{salgarella2025}; the envelope spans ca.
1800--1450 BCE (Middle Minoan II--Late Minoan IB; \S4). Each denominator
must be labelled where it is used (Table~\ref{tab:denoms}). The channel through
which structure is observed is a \emph{sign-occurrence} count, not a
document count or a sign inventory (invariant 7): the broader $\sim$7,400
occurrences count every sign type, while the narrower
$N \approx 5{,}792$ parsed syllable-signs --- after logograms and
numerals are stripped --- is what the unicity computation consumes. Two
further properties tighten the budget: Linear A words are short and
formulaic (the six-position libation-formula template;
\citealp[\S7]{salgarella2025}), and a single language/horizon dominates.

On the logos-normalized corpus, Shannon's unicity-distance criterion
yields $U \approx 204$--$415$ signs across the whole plausible $(V, P)$
parameter space, everywhere far below $N$ (Table~\ref{tab:denoms}) ---
retained only as a toy-model precondition diagnostic (\S4). Unicity
distance presupposes a \emph{known plaintext-language oracle}; because we
estimate the redundancy $D$ from the ciphertext itself, that oracle
collapses into Linear A up to relabelling, so the computation
demonstrates only measurable recurrent structure, not that any reading is
identifiable; and a lossy, non-injective sign$\rightarrow$sound channel
may possess no unicity distance at all (\S9.4).

\paragraph{C.5 Literature index and \texttt{L\_fake} canary (\S3.3).} \emph{Index seed and adversarial population.}
The seed is deliberately conservative --- a sign wrongly omitted is wrongly
granted \texttt{L\_not\_indexed} status, the dangerous direction --- and
every verifiably-published disputed entry is indexed to \emph{catch}
regurgitation, never as an
accepted reading. Population is adversarial: of the seeded West-Semitic
lexical proposals, five are Gordon's equations and a-sa-sa-ra-me = ``oh
Asherah!'' is \citet{best1981}, not Gordon; provenance-unverifiable
candidates --- \texttt{qa-pa} = kappu and a Semitic \texttt{ki-ro}, whose
attribution to Gordon is not supported by the primary account
\citep{rendsburg1996} --- were excluded rather than indexed. Exposure
indexing is independent of evidential validity --- a false public claim
can still contaminate. Generated from the
seed (invariant 12), the current census (Table~\ref{tab:canary}) is a
deliberately conservative, non-exhaustive seed, not yet the full \S C.1
literature census.
\emph{\texttt{L\_fake} calibration.} The canary is calibrated to a real
candidate language's phoneme frequencies, root-template structure, and
lexicon size (\S3.3). Its own divergence from the calibration
targets is reported, never minimised; with a Semitic root-template sampler
and an external Hebrew (ETCBC/bhsa) reject set, the root-template
total-variation distance fell markedly (Table~\ref{tab:canary}), and any
residual real-lexicon collision is measured and published rather than
asserted to be zero.

\section{The order-statistic graduation bar: formula, assumptions,
calibration}

The operative deflation clause (\S3.4) is a normal-order-statistic
\textbf{approximation} to the expected maximum of the multiplicity-null over
the attempted candidates. For a null of mean $\mu_0$ and standard deviation
$\sigma_0$ and $n$ attempted candidates,
\[
  \resizebox{0.95\columnwidth}{!}{$\displaystyle
  E[\max] \approx \mu_0 + \sigma_0\bigl[(1-\gamma)\,\Phi^{-1}(1-\tfrac{1}{n})
  + \gamma\,\Phi^{-1}(1-\tfrac{1}{ne})\bigr]$}
\]
with $\gamma = 0.5772\ldots$ the Euler--Mascheroni constant --- the
blended expected-maximum expression used by \citet{bailey2014}; the
general order-statistics treatment is \citet[\emph{Order Statistics}, 3rd
ed.]{david2003}. The
pipeline exposes the bar as
\texttt{expected\_max\_order\_stat}($\mu_0, \sigma_0, N$); the
deduplicating \texttt{SearchLog}
(\S3.2, invariant 12) reports the exact distinct-candidate count
\texttt{n\_trials} a scanner produces. The bar fails closed on a degenerate null
($\sigma_0 \le 0 \Rightarrow E[\max] = \mu_0 \Rightarrow$ no graduation).
The \texttt{generalizes\_to\_not\_indexed} clause (\S3.4) requires recovery
above a pre-committed threshold and on a minimum number of distinct
not-indexed signs, and fails closed when no threshold or no supported
not-indexed sign is present. The full clause conjunction: the candidate
must separate from the \texttt{L\_fake} negative-control distribution,
its corrected margin must be cleared, and the observed statistic must
beat the order-statistic bar over the Packard/random-lexeme multiplicity
null.

\textbf{Assumptions and the direction of their error.} The approximation is
exact only under (i) an approximately Gaussian null and (ii) independence
among the $n$ attempted candidates. Neither holds exactly. (i) The
permutation / \texttt{L\_fake} nulls are not perfectly Gaussian; the blended
Euler-$\gamma$ form is a first-order correction, not a guarantee. (ii) The
candidates a real search tries are typically \textbf{positively correlated}
(overlapping sign-value assignments). Under a common-variance Gaussian
comparison model, positive dependence lowers the expected maximum relative
to independence (a Slepian-type comparison; \citealp{slepian1962}), so
using the
raw attempted-candidate count $n$ in place of a smaller effective
independent-trial count places the bar \emph{above} the correlated truth;
this does \textbf{not} establish conservatism for arbitrary heterogeneous
or non-Gaussian dependence --- the empirical best-of-100 calibration
therefore carries the operational error-control claim. We use the
raw deduplicated candidate count; a dependence-adjusted effective count is a
recorded, deferred refinement (\S9.6) whose correction would only strengthen
the nulls.

\textbf{Calibration and its interval.} Because the analytic bar rests on
assumptions that do not hold exactly, the error-control claim is carried not
by the formula but by the direct empirical calibration (\S3.4): a
Monte-Carlo of the full gate on a best-of-100 random-map null-generating
process (\texttt{scripts/\allowbreak gate\_null\_calibration.py}). The
interval on that calibration is the \textbf{one-sided exact
Clopper--Pearson} bound (the Beta-quantile inversion of the binomial), which
needs no normal approximation to the calibration count --- for 3 graduations
in 500 replicates, a 95\% upper bound of $\approx 1.54\%$; the verdict
partition of the 500 replicates is in Table~\ref{tab:calib} (no
INCOMPLETE, the search being instrumented).
\citet{bailey2014} supply the selective-inference lineage (\S2.4) and the
blended expression above; the operational error-control claim nonetheless
rests on the empirical calibration, not on the formula.

\textbf{Cognate-ID baseline (\S3.4).} The match-rate the deflation is
applied to is supplied by the published cognate-ID matcher identified in
\S2.3 (reproduced in the artifact), so the gate does not depend on any
single citation.

\begin{center}\rule{0.5\columnwidth}{0.4pt}\end{center}

\noindent\emph{Software \& data.}
\ifanon
The artifact. Code, pre-registrations, and findings are released as an
anonymized artifact on acceptance. Licensed corpora (GORILA-, SigLA-,
lineara.xyz-, D\=AMOS-derived) are not redistributed; a data-provenance note
documents each source's licence and retrieval path.
\else
logos (v0.1.0). Zenodo. Concept DOI 10.5281/zenodo.21087572; repository
\url{https://github.com/papadopouloskyriakos/logos}. Licensed corpora
(GORILA-, SigLA-, lineara.xyz-, D\=AMOS-derived) are not redistributed; see
\texttt{docs/data-provenance.md}.
\fi

\section{Complementary tables}

\begin{table}[h]
\centering\small
\resizebox{\columnwidth}{!}{%
\begin{tabular}{@{}lrrrl@{}}
\toprule
\textbf{Stratum} & \textbf{Inscr.} & \textbf{Words} & \textbf{Sites} &
\textbf{Role} \\
\midrule
admin & 290 & 1,486 & 14 & induction \\
libation (stone vessels) & 99 & 370 & 15 & induction \\
other & 129 & 187 & 40 & induction \\
seal & 740 & --- & --- & excluded \\
emptied by abbrev.-strip & 83 & --- & --- & excluded \\
\midrule
total & 1,341 & & & \\
\bottomrule
\end{tabular}}
\caption{Morphology corpus strata (Appendix~A.1). Seals are structurally
excluded as an abbreviation channel; the strata reconcile exactly to the
1,341-inscription corpus total.}
\label{tab:strata}
\end{table}

\begin{table}[h]
\centering\small
\resizebox{\columnwidth}{!}{%
\begin{tabular}{@{}lll@{}}
\toprule
\textbf{Bucketing} & $\alpha=0.01$ & $\alpha=0.05$ \\
\midrule
strict (stone vessels only) & NULL & NULL \\
loose (stone objects + inked) & NULL & \texttt{i-}, \texttt{-te/-ti}
validate \\
\bottomrule
\end{tabular}}
\caption{Morphology sensitivity $2\times2$ (Appendix~A.1): genre bucketing
$\times$ significance threshold. Only the doubly-relaxed corner validates
the H3 nominal affixes \texttt{i-} (locative) and \texttt{-te/-ti}
(ablative).}
\label{tab:sens}
\end{table}

\begin{table}[h]
\centering\small
\resizebox{\columnwidth}{!}{%
\begin{tabular}{@{}ll@{}}
\toprule
\textbf{Quantity} & \textbf{Value} \\
\midrule
Linear A confirm rate (16 affixes) & 0.250 (4/16; [0.10, 0.50]) \\
\texttt{L\_fake} bigram-floor rate & 0.375 (6/16; [0.18, 0.61]) \\
\texttt{i-*301} vs.\ within-form null & $z \approx 8.0$ (6 environments) \\
\midrule
Linear B (D\=AMOS) confirm rate & 0.562 (9/16) \\
Linear B within-word-shuffle floor & $\sim$0.03 \\
Linear B \texttt{L\_fake} bigram floor & $\sim$0.30 (5 draws) \\
\midrule
word-token mean length & 1.84 signs (median/mode 1) \\
word tokens $\le$2 signs / 1 sign & 76.1\% / 56.5\% \\
multi-sign \emph{type} mean length & 3.07 (Fuls 3.3, 554-type list) \\
Linear B mean length & 3.23 (13,562 wordforms) \\
\midrule
DP-unigram segmenter micro-F1 & 0.436 \\
random-boundary baseline & 0.389 (base rate 0.406) \\
real$-$random gap, 95\% CI & [0.021, 0.099]; P(gap${>}$0) = 0.998 \\
Morfessor & 0.156 (over-segments; not usable) \\
\bottomrule
\end{tabular}}
\caption{Morphology-induction replication quantities (Appendix~A.1):
confirm rates with Wilson 95\% intervals, word-length denominators, and
the segmentation positive with its site-clustered bootstrap gap (4,000
draws over 52 sites).}
\label{tab:morphrates}
\end{table}

\begin{table}[h]
\centering\small
\resizebox{\columnwidth}{!}{%
\begin{tabular}{@{}ll@{}}
\toprule
\textbf{Quantity} & \textbf{Value} \\
\midrule
documents / balance units & 32 / 35 \\
integer-level balance & 7/28 (25\%) \\
determinable fraction value & J = 1/2 (HT104: $2 \cdot J = 1$) \\
held-out fraction balance (real) & 0.0 \\
permutation null & mean 0.113 (max 0.143); $p = 1.0$ \\
planted positive control & 0.9 vs.\ null 0.22; $p = 0.003$ \\
\bottomrule
\end{tabular}}
\caption{Metrology (Direction D) replication quantities (Appendix~B.1).
The real corpus is not separated from the permutation null; the planted
control separates, so the null is not a machinery bug.}
\label{tab:metro}
\end{table}

\begin{table}[h]
\centering\small
\resizebox{\columnwidth}{!}{%
\begin{tabular}{@{}lrrrr@{}}
\toprule
\textbf{Test} & \textbf{Classes} & \textbf{Acc.} & \textbf{Null} & $p$ \\
\midrule
consonant class & 13 & 0.080 & 0.065 & 0.40 \\
vowel class & 5 & 0.140 & 0.201 & 0.86 \\
\bottomrule
\end{tabular}}
\caption{Distributional-phonology per-test results (Appendix~B.2):
leave-one-out 1-NN vs.\ a label-permutation null, \v{S}id\'{a}k-deflated
over the two tests. The positive-control power curve fires ($p<.05$) at
planted strength $\ge 2$ but not at 1.}
\label{tab:phon}
\end{table}

\begin{table}[h]
\centering\small
\resizebox{\columnwidth}{!}{%
\begin{tabular}{@{}lll@{}}
\toprule
\textbf{Leg} & \textbf{Accuracy} & \textbf{Floor} \\
\midrule
classical, direct-NN & 0.410 [0.18, 0.64] & 0.0135 \\
classical, Procrustes & 0.185 [0.00, 0.36] & 0.0135 \\
I-JEPA (3 seeds; excluded) & $0.324 \pm 0.025$ & 0.0135 \\
font-swap, direct-NN & $0.367 \rightarrow 0.416$ & $\gg$ chance \\
font-swap, Procrustes & $0.156 \rightarrow 0.195$ & $\gg$ chance \\
sequence/cognate (D\=AMOS) & 0.025 & 0.011 \\
\quad positive control & 0.947 & --- \\
\bottomrule
\end{tabular}}
\caption{Cross-script sign-image alignment and controls (Appendix~B.3).
Font-swap rows give same-font $\rightarrow$ independent-fonts accuracy on
the control's matched anchor set; the sequence/cognate row is the best of
five alignment methods over 56 anchors
(\texttt{clearly\_above\_chance} = False).}
\label{tab:xscript}
\end{table}

\begin{table}[h]
\centering\small
\resizebox{\columnwidth}{!}{%
\begin{tabular}{@{}llrr@{}}
\toprule
\textbf{Item} & \textbf{Source} & \textbf{LLM} & \textbf{Baseline} \\
\midrule
\texttt{JA-NE} = \emph{yn} ``wine'' & Gordon & 14/40 & 0/40 \\
\texttt{A-SA-SA-RA-ME} & \citet{best1981} & 2/40 & 0/40 \\
\texttt{KA-RO-PA} & Gordon & 0/40 & 0/40 \\
\texttt{SU-PA-RA} & Gordon & 0/40 & 0/40 \\
\texttt{KU-RO} = \emph{kull} (downwt.) & Gordon & 8/40 & --- \\
\texttt{KU-RO} = ``total'' (excl.) & gen.\ knowledge & 7/40 & --- \\
\bottomrule
\end{tabular}}
\caption{LLM-ablation per-item reproduction over 40 seeds,
\texttt{qwen2.5:72b} vs.\ the model-free edit-distance baseline
(Appendix~C.1). The two \texttt{KU-RO} routes overlap on one seed; three
single-word items are among Gordon's five West-Semitic equations, while
\texttt{A-SA-SA-RA-ME} is Best's.}
\label{tab:repro}
\end{table}

\begin{table}[h]
\centering\small
\resizebox{\columnwidth}{!}{%
\begin{tabular}{@{}lrr@{}}
\toprule
\textbf{Sign set} & \textbf{Signs} & \textbf{Valued (of 40)} \\
\midrule
known AB (indexed) & 20 & $\sim$20 \\
not-indexed *-series & 11 & $\sim$1 ($\sim$97\% abstention) \\
\bottomrule
\end{tabular}}
\caption{\texttt{L\_not\_indexed} abstention asymmetry (Appendix~C.2).
The apparent consistency gap 0.601 (permutation $p = 0.0005$) is
coverage-confounded; the matched-$n$ check returned \texttt{no\_power}
(0 of 11 not-indexed signs reached $n \ge 10$).}
\label{tab:abstain}
\end{table}

\begin{table}[h]
\centering\small
\resizebox{\columnwidth}{!}{%
\begin{tabular}{@{}ll@{}}
\toprule
\textbf{Denominator / channel} & \textbf{Value} \\
\midrule
token-frequent syllabograms & 97 \\
full syllabary estimate & $\sim$150 \\
working repertoire ($V$; conflates syll./subscr./logograms) & 259 \\
incl.\ ideograms (GORILA: 390) & $\sim$350 \\
\midrule
sign occurrences, all types & $\sim$7,400 \\
parsed syllable-signs (unicity input) & $N \approx 5{,}792$ \\
\midrule
corpus envelope & $\sim$2,500 finds; $\sim$1,400 readable \\
findplaces (logos-normalized) & 56 ($\rightarrow$ 52 sites) \\
readable inscriptions & 1,341 \\
distinct word-forms & 539 \\
\midrule
unicity distance over $(V,P)$ space & $U \approx 204$--$415$ signs \\
\bottomrule
\end{tabular}}
\caption{Corpus denominators, occurrence channels, and the unicity range
(Appendix~C.4). The $\sim$7,400 all-type occurrence channel never enters
the unicity arithmetic, which consumes only the parsed syllable-signs.}
\label{tab:denoms}
\end{table}

\begin{table}[h]
\centering\small
\resizebox{\columnwidth}{!}{%
\begin{tabular}{@{}ll@{}}
\toprule
\textbf{Calibration quantity} & \textbf{Value} \\
\midrule
\texttt{S\_morph} FPR (within-form shuffle) & 3.0\% (6/200; [1.4, 6.4]\%) \\
nominal one-sided $z \ge 2$ rate & $\approx$2.3\% \\
\texttt{S\_morph} power (planted affix) & 100\% (50/50; saturated) \\
\midrule
gate false-graduation (best-of-100 null) & 0.6\% (3/500; CP $\le$ 1.54\%) \\
gate verdict partition (500 replicates) & 3 GRAD / 10 NULL / 487 REJ \\
\bottomrule
\end{tabular}}
\caption{Detector and gate calibration quantities (Appendices~A.2
and~D): Wilson interval for the \texttt{S\_morph} false-positive rate;
one-sided exact Clopper--Pearson 95\% upper bound for the gate's
false-graduation rate under the tested best-of-100 random-map null.}
\label{tab:calib}
\end{table}

\begin{table}[h]
\centering\small
\resizebox{\columnwidth}{!}{%
\begin{tabular}{@{}ll@{}}
\toprule
\textbf{Quantity} & \textbf{Value} \\
\midrule
literature-index census & 63 claims / 62 entries (55 single) \\
root-template TV distance & $\sim$0.84 $\rightarrow$ $\sim$0.33
  ($\sim$0.23 standalone) \\
\midrule
ablation scale & 40 seeds $\times$ 40 held-out forms \\
model-free lexicon & 7,682 Hebrew skeletons \\
LLM errors & 0 (40/40 seeds) \\
value-vocabulary overlap & 10 of 134 \\
first-pass contamination rate & 1.000 (33/33 defined; artifact) \\
gate-2 rate & mean 0.646 (confounded; not reported) \\
\bottomrule
\end{tabular}}
\caption{Literature-index, \texttt{L\_fake}-canary, and LLM-ablation
quantities (Appendices~C.1 and~C.5). The TV distance is measured
end-to-end on the Ugaritic$\leftrightarrow$Hebrew gold; the first-pass
contamination rate is forced to 1.0 by construction and reported as an
artifact.}
\label{tab:canary}
\end{table}

\end{document}
